% Worked solutions for Listas/lista-r10-macro-implications.md (15 exercises, Tirole Ch. 13-16 and Part VII).
% Compile: pdflatex twice; check: pdffonts (Type 1 only, uni = yes).
% Every number is asserted in lista-r10-macro-implications-solutions_check.py.
\documentclass[11pt,a4paper]{article}
\usepackage[utf8]{inputenc}
\usepackage[T1]{fontenc}
\usepackage{lmodern}   % Type 1 vector fonts (replaces PK bitmaps)
\usepackage{cmap}      % ToUnicode CMap -> copyable text and accents
\usepackage[english]{babel}
\usepackage{amsmath,amssymb,amsfonts}
\usepackage[a4paper,margin=2.2cm]{geometry}
\usepackage{booktabs}
\usepackage{enumitem}
\usepackage{fancyhdr}
\usepackage{xcolor}
\usepackage{microtype}

\pagestyle{fancy}\fancyhf{}
\lhead{\footnotesize Finan\c{c}as II --- Corporativas (MPE Insper)}
\rhead{\footnotesize List R10 --- Worked solutions}
\cfoot{\footnotesize\thepage}
\renewcommand{\headrulewidth}{0.3pt}
\setlength{\headheight}{14pt}
\setlist{topsep=2pt,itemsep=1pt}
\allowdisplaybreaks

\newcommand{\exsec}[1]{\section*{\color{blue!55!black}#1}\addcontentsline{toc}{section}{#1}}
\newcommand{\notation}{\par\noindent\textbf{Notation.}}
\newcommand{\step}[2]{\par\medskip\noindent\textbf{Step #1 --- #2.}~}
\newcommand{\intuition}{\par\medskip\noindent\textbf{Intuition.}~}
\newcommand{\bookans}{\par\medskip\noindent\textbf{Book answer.}~}
\newcommand{\refbib}{\par\smallskip\noindent\textbf{Ref:}~}
\newcommand{\asked}{\par\medskip\noindent\textbf{What is asked.}~}
\newcommand{\approach}{\par\smallskip\noindent\textbf{Approach.}~}
\newcommand{\answer}[1]{\par\smallskip\noindent\fbox{\textbf{Answer\if\relax\detokenize{#1}\relax\else\ #1\fi.}}~}
\newcommand{\pitfalls}{\par\medskip\noindent\textbf{Common mistakes.}}
\newcommand{\Rc}{\mathcal{R}}
\newcommand{\Pc}{\mathcal{P}}

\usepackage[hyphens]{url}
%% STATEMENT-PREAMBLE BEGIN (tools/statements_into_solutions.py)
\usepackage{longtable,array,calc}
\providecommand{\tightlist}{\setlength{\itemsep}{0pt}\setlength{\parskip}{0pt}}
\providecommand{\href}[2]{#2}
\providecommand{\real}[1]{#1}
\makeatletter\@ifundefined{c@none}{\newcounter{none}}{}\makeatother  % pandoc's \LTcaptype{none}
\newcommand{\stmtbegin}{\par\medskip\noindent\textbf{Statement.}\par\nopagebreak\begingroup\small\setlength{\tabcolsep}{3pt}\sloppy}
\newcommand{\stmtend}{\par\endgroup\nopagebreak\noindent\rule{\linewidth}{0.4pt}\par\medskip}
\DeclareUnicodeCharacter{2192}{\ensuremath{\rightarrow}}
\DeclareUnicodeCharacter{2212}{\ensuremath{-}}
\DeclareUnicodeCharacter{2022}{\textbullet{}}
\DeclareUnicodeCharacter{2248}{\ensuremath{\approx}}
\DeclareUnicodeCharacter{2032}{\ensuremath{^{\prime}}}
\DeclareUnicodeCharacter{2081}{\ensuremath{_{1}}}
\DeclareUnicodeCharacter{03B8}{\ensuremath{\theta}}
\DeclareUnicodeCharacter{03BB}{\ensuremath{\lambda}}
\DeclareUnicodeCharacter{03C8}{\ensuremath{\psi}}
\DeclareUnicodeCharacter{03B1}{\ensuremath{\alpha}}
\DeclareUnicodeCharacter{25A0}{\rule{0.7em}{0.7em}}
%% STATEMENT-PREAMBLE END
\begin{document}
\begin{center}{\Large\bfseries Exercise list R10 --- Macroeconomic implications of corporate-finance frictions}\\[4pt]{\large Worked solutions}\end{center}

\noindent Statements are in \url{Listas/lista-r10-macro-implications.md}; each exercise is cited by number and printed page and its full text is reproduced right before its solution. Short answers: \url{lista-r10-macro-implications-solutions-short.pdf}. All references are to Tirole, \emph{The Theory of Corporate Finance} (Princeton, 2005); ``book answer'' = Part VII, ``Answers to Selected Exercises''. Where a parameter example is given, it satisfies every assumption of the exercise and each of its numbers is asserted in \url{lista-r10-macro-implications-solutions_check.py}.

\tableofcontents

% =============================================================================
\exsec{1. Investor protection and the interest rate --- Tirole 13.1, p.~493}
%% BEGIN-STATEMENT 1
\stmtbegin
\textbf{Source:} Tirole, Ch. 13, Exercise 13.1, p.~493 (PDF 506) · Answer in book: yes, pp.~616--617

\textbf{Exercise 13.1 (improved governance).} There are two dates, \emph{t} = 0, 1, and a continuum
of mass 1 of firms. Firms are identical except for the initial wealth \(A\) initially owned by their
entrepreneur. \(A\) is distributed according to continuous cumulative distribution \(G(A)\) with
density \(g(A)\) on \([0, I]\).

Each entrepreneur has a fixed-size project, and must invest \(I\), and therefore borrow \(I - A\),
at date 0 in order to undertake it. Those entrepreneurs who do not invest themselves, invest their
wealth in other firms. The savings function of nonentrepreneurs (consumers) is an increasing
function \(S(r)\), where \(r\) is the interest rate, with \(S(r) = 0\) for \(r < 0\) (so total
savings equal \(S(r)\) plus the wealth of unfinanced entrepreneurs). Entrepreneurs have utility
\(c_0 + c_1\) from consumptions \(c_0\) and \(c_1\).

A project, if financed, yields \(R > 0\) at date 1 with probability \(p\) and 0 with probability
\(1 - p\). The probability of success is \(p_H\) if the entrepreneur works and
\(p_L = p_H - \Delta p\) if she shirks. The entrepreneur obtains private benefit \(B\) by shirking
and 0 otherwise. Assume \(p_H R > I > p_H(R - B/\Delta p)\), that financing cannot occur if the
entrepreneur is provided with incentives to misbehave, and that the equilibrium interest rate is
strictly positive.

\begin{enumerate}
\def\labelenumi{(\roman{enumi})}
\item
  What is the pledgeable income? Write the financing condition.
\item
  Give the expression determining the market rate of interest. How does this interest rate change
  when improved investor protection lowers \(B\)?
\end{enumerate}

\paragraph{Context needed to solve it}

None: the statement is self-contained (it is the fixed-investment model of Tirole §3.2 with an
endogenous interest rate, as in §13.2).
\stmtend
%% END-STATEMENT

\asked (i) How much of the project's return can be promised to investors, and which entrepreneurs (which wealth levels $A$) obtain financing when investors require the gross return $1+r$. (ii) Which equation pins down the market interest rate $r$, and whether $r$ rises or falls when investor protection improves (lower $B$).

\approach Each firm is the fixed-investment model of Tirole \S3.2: incentives cap the \emph{pledgeable income}, and only entrepreneurs with enough net worth pass the solvency test. With an endogenous interest rate the cutoff wealth $A^*(r)$ moves with $r$, and $r$ clears the market for loanable funds (Tirole \S13.2). The comparative static follows from the implicit-function theorem.

\notation
\begin{itemize}
  \item $A\in[0,I]$: entrepreneur's wealth, c.d.f.\ $G$, density $g$; $I$: fixed investment; $R$: return if success.
  \item $p_H$, $p_L=p_H-\Delta p$: success probabilities when working / shirking; $B$: private benefit from shirking.
  \item $R_b$: entrepreneur's reward if success; $\Pc\equiv p_H(R-B/\Delta p)$: pledgeable income.
  \item $r>0$: interest rate; $S(r)$, $S'>0$: consumers' savings; $A^*(r)$: cutoff wealth for financing.
\end{itemize}

\step{1}{Incentive constraint and pledgeable income} The entrepreneur works iff her expected reward when working is at least the one when shirking plus $B$:
\begin{align}
  p_HR_b &\geq p_LR_b+B \nonumber\\
  (\Delta p)R_b &\geq B \nonumber\\
  R_b &\geq \frac{B}{\Delta p}. \label{r10e1-ic}
\end{align}
Investors receive $R-R_b$ in case of success, so the most they can expect is, using the smallest $R_b$ in \eqref{r10e1-ic},
\begin{align}
  \Pc = p_H\left(R-\frac{B}{\Delta p}\right). \label{r10e1-P}
\end{align}

\step{2}{Financing condition} Investors lend $I-A$ and require $(1+r)(I-A)$ in expectation; financing is possible iff
\begin{align}
  (1+r)(I-A) &\leq p_H\left(R-\frac{B}{\Delta p}\right) \label{r10e1-fc}\\
  \iff A &\geq A^*(r)\equiv I-\frac{p_H(R-B/\Delta p)}{1+r}. \label{r10e1-Astar}
\end{align}
By the assumption $I>\Pc$, $A^*(r)>0$ for every $r\geq0$. A financed entrepreneur gets $p_HR-(1+r)(I-A)$ instead of $(1+r)A$ from lending her wealth; she prefers to invest iff $p_HR\geq(1+r)I$, which holds for the moderate $r$ of the equilibrium (it holds in the example below).

\answer{(i)} Pledgeable income $p_H(R-B/\Delta p)$; financing iff $(1+r)(I-A)\leq p_H(R-B/\Delta p)$, i.e.\ $A\geq A^*(r)$.

\step{3}{Market for loanable funds} Financed firms (mass $1-G(A^*)$) demand $I-A$ each; consumers supply $S(r)$ and unfinanced entrepreneurs lend their wealth:
\begin{align}
  S(r)+\int_0^{A^*(r)}A\,g(A)\,dA = \int_{A^*(r)}^{I}(I-A)\,g(A)\,dA. \label{r10e1-mc}
\end{align}
Adding $\int_{A^*}^{I}A\,g(A)\,dA$ to both sides and using $\int_{A^*}^{I}I\,g\,dA=[1-G(A^*)]I$:
\begin{align}
  S(r)+\int_0^{I}A\,g(A)\,dA = [1-G(A^*(r))]\,I. \label{r10e1-mc2}
\end{align}
Total savings (left) must equal total investment (right).

\step{4}{Uniqueness} The left side of \eqref{r10e1-mc2} increases with $r$ ($S'>0$). Differentiating \eqref{r10e1-Astar},
\begin{align}
  \frac{\partial A^*}{\partial r} = \frac{p_H(R-B/\Delta p)}{(1+r)^2}>0, \label{r10e1-dAr}
\end{align}
so the right side, $[1-G(A^*)]I$, decreases with $r$. The two curves cross at most once.

\step{5}{Effect of $B$} Write \eqref{r10e1-mc2} as $\Phi(r,B)\equiv S(r)+\int_0^IA\,g\,dA-[1-G(A^*(r,B))]I=0$. Its partial derivatives are
\begin{align}
  \Phi_r &= S'(r)+g(A^*)\,I\,\frac{\partial A^*}{\partial r}>0, \label{r10e1-Phir}\\
  \Phi_B &= g(A^*)\,I\,\frac{\partial A^*}{\partial B}, \qquad \frac{\partial A^*}{\partial B}=\frac{p_H}{(1+r)\Delta p}>0. \label{r10e1-PhiB}
\end{align}
By the implicit-function theorem,
\begin{align}
  \frac{dr}{dB} = -\frac{\Phi_B}{\Phi_r} = -\frac{g(A^*)\,I\,p_H/[(1+r)\Delta p]}{S'(r)+g(A^*)\,I\,\Pc/(1+r)^2}<0. \label{r10e1-drdB}
\end{align}
A lower $B$ (better protection) raises $r$.

\step{6}{Example} $I=1$, $G$ uniform on $[0,1]$, $S(r)=2r$, $p_H=0.8$, $R=1.5$, $\Delta p=0.4$. Then $\int_0^1A\,dA=0.5$, $1-G(A^*)=\Pc/(1+r)$, and \eqref{r10e1-mc2} reads $(2r+0.5)(1+r)=\Pc$.
With $B=0.3$: $\Pc=0.8(1.5-0.75)=0.6$, so $2r^2+2.5r-0.1=0$ and $r=(-2.5+\sqrt{6.25+0.8})/4=0.038795$; $A^*=1-0.6/1.038795=0.4224$; financed share $0.6/1.038795=0.577592$.
With $B=0.2$: $\Pc=0.8(1.5-0.5)=0.8$, so $2r^2+2.5r-0.3=0$, $r=0.110272$, $A^*=0.279456$, financed share $0.720544$. Also $(1+r)I=1.11\leq p_HR=1.2$.

\answer{(ii)} $r$ solves $S(r)+\int_0^IA\,g(A)\,dA=[1-G(A^*(r))]I$ with $A^*(r)=I-p_H(R-B/\Delta p)/(1+r)$; lowering $B$ raises $r$ (eq.~\eqref{r10e1-drdB}).

\intuition Better governance raises the income that can be pledged, so more entrepreneurs pass the solvency test and the demand for loanable funds rises. With savings that respond to $r$, the interest rate goes up; this transfers part of the gain to savers and raises the cutoff $A^*$ again, so general equilibrium dampens the expansion of credit (it does not reverse it: the financed share still rises).

\pitfalls
\begin{itemize}
  \item Writing the financing condition as $\Pc\geq I-A$, forgetting that investors now require $1+r$.
  \item Omitting the wealth of unfinanced entrepreneurs from the supply side of \eqref{r10e1-mc}.
  \item Concluding that investment rises one-for-one with the partial-equilibrium fall in $A^*$: the rise in $r$ pushes $A^*$ back up.
\end{itemize}

\bookans Matches the book's answer (pp.~616--617): same cutoff, same clearing equation in both forms, and ``an increase in investor protection raises the equilibrium interest rate.'' The derivative \eqref{r10e1-drdB} and the example are additions.
\refbib{Tirole \S3.2 (fixed investment); \S13.2 ``Capital Squeezes and Economic Activity: The Balance-Sheet Channel'', pp.~471--478.}

% =============================================================================
\exsec{2. Time-consistent government policy --- Tirole 16.2, pp.~564--565}
%% BEGIN-STATEMENT 2
\stmtbegin
\textbf{Source:} Tirole, Ch. 16, Exercise 16.2, pp.~564--565 (PDF 577) · Answer in book: yes, p.~622

\textbf{Exercise 16.2 (time-consistent government policy).} Consider a unit mass of identical
entrepreneurs with variable-investment projects. The timing is summarized in Figure 16.11.

Figure 16.11 (timeline, left to right; p.~564, PDF 577)

{\def\LTcaptype{none} % do not increment counter
\begin{longtable}[]{@{}
  >{\raggedright\arraybackslash}p{(\linewidth - 6\tabcolsep) * \real{0.2500}}
  >{\raggedright\arraybackslash}p{(\linewidth - 6\tabcolsep) * \real{0.2500}}
  >{\raggedright\arraybackslash}p{(\linewidth - 6\tabcolsep) * \real{0.2500}}
  >{\raggedright\arraybackslash}p{(\linewidth - 6\tabcolsep) * \real{0.2500}}@{}}
\toprule\noalign{}
\begin{minipage}[b]{\linewidth}\raggedright
Step 1
\end{minipage} & \begin{minipage}[b]{\linewidth}\raggedright
Step 2
\end{minipage} & \begin{minipage}[b]{\linewidth}\raggedright
Step 3
\end{minipage} & \begin{minipage}[b]{\linewidth}\raggedright
Step 4
\end{minipage} \\
\midrule\noalign{}
\endhead
\bottomrule\noalign{}
\endlastfoot
Representative entrepreneur invests \(I\), borrows \(I - A\). & Government selects \(\tau\). &
Entrepreneurial moral hazard: behaves --- Pr(success) \(p_H + \tau\), private benefit 0; misbehaves
--- Pr(success) \(p_L + \tau\), private benefit \(BI\). & Outcome: \(RI\) with probability
\(p + \tau\), 0 with probability \(1 - (p + \tau)\). \\
\end{longtable}
}

The cost of the policy for the country is \(\gamma(\tau)I\) (where \(\gamma'(0) = 0\),
\(\gamma'(\tau) > 0\), for \(\tau > 0\), \(\gamma'(1 - p_H) = \infty\), \(\gamma'' > 0\)).

All investors are domestic investors (there are no foreign lenders and, when choosing \(\tau\), the
government maximizes social welfare, equal to entrepreneurs' welfare plus investors' welfare).

Assume that

\[(p_H + \tau)R > 1 > (p_H + \tau)\left(R - \frac{B}{\Delta p}\right)\]

in the relevant range of values of \(\tau\), and that it is never optimal to induce entrepreneurs to
misbehave. Everyone is risk neutral, and the entrepreneurs are protected by limited liability.

\begin{enumerate}
\def\labelenumi{(\roman{enumi})}
\tightlist
\item
  • Show that, when expecting policy \(\tau\), entrepreneurs invest
\end{enumerate}

\[I(\tau) = \frac{A}{1 - (p_H + \tau)(R - B/\Delta p)}.\]

• What is the equilibrium value \(\tau^*\)?

\begin{enumerate}
\def\labelenumi{(\roman{enumi})}
\setcounter{enumi}{1}
\item
  What value would the government choose if it selected \(\tau\) \emph{before} entrepreneurs borrow?
\item
  Informally explain how your answer to (i) would change if investors were foreign investors and the
  government discounted their welfare relative to that of domestic residents.
\end{enumerate}

\paragraph{Context needed to solve it}

None beyond Figure 16.11, transcribed above (the variable-investment model is restated in the
statement).
\stmtend
%% END-STATEMENT

\asked (i) The investment that entrepreneurs undertake when they expect policy $\tau$, and the policy the government actually picks once investment is sunk. (ii) The policy a government that moves first would choose. (iii) How the answer to (i) changes if investors are foreign and the government gives them less weight.

\approach Each firm is the variable-investment model (Tirole \S3.4) with success probability shifted by $\tau$. The government's policy raises everyone's success probability but costs $\gamma(\tau)I$. Comparing the ex post choice (investment fixed) with the ex ante choice (investment responds to $\tau$) is the classic time-consistency problem of Tirole \S16.3.2: ex post the government ignores the effect of policy on borrowing capacity.

\notation
\begin{itemize}
  \item $\tau\geq0$: policy (raises the probability of success to $p_H+\tau$ or $p_L+\tau$); $\gamma(\tau)I$: its social cost, $\gamma'>0$, $\gamma''>0$.
  \item $R$, $B$, $\Delta p$, $A$, $I$: as in the variable-investment model; $R_b$, $R_l=RI-R_b$: entrepreneur's and investors' payoffs if success.
  \item $W$: social welfare = entrepreneurs' + investors' welfare.
\end{itemize}

\step{1}{Incentive constraint} With policy $\tau$, behaving gives success probability $p_H+\tau$, misbehaving $p_L+\tau$ and $BI$:
\begin{align}
  (p_H+\tau)R_b &\geq (p_L+\tau)R_b+BI \nonumber\\
  [(p_H+\tau)-(p_L+\tau)]R_b &\geq BI \nonumber\\
  R_b &\geq \frac{BI}{\Delta p}. \label{r10e2-ic}
\end{align}
The policy does not change the incentive constraint: $\tau$ cancels.

\step{2}{Borrowing capacity} Investors break even (rate 0) if $(p_H+\tau)(RI-R_b)=I-A$. With \eqref{r10e2-ic} binding:
\begin{align}
  (p_H+\tau)\left(R-\frac{B}{\Delta p}\right)I &= I-A \nonumber\\
  I\left[1-(p_H+\tau)\left(R-\frac{B}{\Delta p}\right)\right] &= A \nonumber\\
  I(\tau) &= \frac{A}{1-(p_H+\tau)(R-B/\Delta p)}. \label{r10e2-I}
\end{align}
The denominator is positive by the assumption $1>(p_H+\tau)(R-B/\Delta p)$, and $I'(\tau)>0$.

\step{3}{Ex post policy} Once $I$ is sunk, welfare is expected output minus the cost of the policy minus the (sunk) investment:
\begin{align}
  W^{\text{ex post}}(\tau) &= (p_H+\tau)RI-\gamma(\tau)I-I, \label{r10e2-Wexpost}\\
  \frac{\partial W^{\text{ex post}}}{\partial\tau} &= RI-\gamma'(\tau)I = 0 \;\Longrightarrow\; \gamma'(\tau^*)=R. \label{r10e2-tau}
\end{align}
The second-order condition holds ($-\gamma''I<0$). Investors anticipate $\tau^*$, so $I=I(\tau^*)$ from \eqref{r10e2-I}.

\answer{(i)} $I(\tau)=A/[1-(p_H+\tau)(R-B/\Delta p)]$; equilibrium policy $\gamma'(\tau^*)=R$, investment $I(\tau^*)$.

\step{4}{Ex ante policy} Choosing $\tau$ before borrowing, the government knows investment will be $I(\tau)$. Investors break even, so welfare is the NPV:
\begin{align}
  W(\tau) &= [(p_H+\tau)R-1-\gamma(\tau)]\,I(\tau), \label{r10e2-W}\\
  W'(\tau^c) &= [R-\gamma'(\tau^c)]\,I(\tau^c)+[(p_H+\tau^c)R-1-\gamma(\tau^c)]\,I'(\tau^c)=0, \nonumber\\
  [\gamma'(\tau^c)-R]\,I(\tau^c) &= [(p_H+\tau^c)R-1-\gamma(\tau^c)]\,\frac{dI}{d\tau}(\tau^c). \label{r10e2-tauc}
\end{align}
The right side is positive (NPV per unit positive, $I'>0$), so $\gamma'(\tau^c)>R=\gamma'(\tau^*)$ and, $\gamma'$ being increasing, $\tau^c>\tau^*$.

\step{5}{Example} $p_H=0.5$, $R=2$, $B/\Delta p=1$, $A=1$, $\gamma(\tau)=10\tau^2$ (so $\gamma'=20\tau$; this quadratic only illustrates the relevant range). Then $\tau^*=R/20=0.1$, $I(\tau^*)=1/(1-0.6)=2.5$, $W(\tau^*)=0.25$. Ex ante, $W(\tau)=(2\tau-10\tau^2)/(0.5-\tau)$; $W'=0$ gives $(2-20\tau)(0.5-\tau)+(2\tau-10\tau^2)=0$, i.e.\ $10\tau^2-10\tau+1=0$, so $\tau^c=(10-\sqrt{60})/20=0.112702>\tau^*$ and $W(\tau^c)=0.254033$.

\answer{(ii)} $\tau^c$ solves \eqref{r10e2-tauc}; $\gamma'(\tau^c)>R$, so $\tau^c>\tau^*$: commitment yields a more investor-friendly policy.

\step{6}{Foreign investors} Suppose the government puts weight $\omega<1$ on investors' payoff $(p_H+\tau)R_l$. Ex post it maximizes $(p_H+\tau)(RI-R_l)+\omega(p_H+\tau)R_l-\gamma(\tau)I$:
\begin{align}
  RI-(1-\omega)R_l-\gamma'(\tau)I &= 0 \nonumber\\
  \gamma'(\tau_\omega) &= R-(1-\omega)\frac{R_l}{I} = R-(1-\omega)\left(R-\frac{B}{\Delta p}\right)<R, \label{r10e2-omega}
\end{align}
using $R_l=(R-B/\Delta p)I$ at the binding \eqref{r10e2-ic}. In the example, $\omega=0.5$ gives $\gamma'=1.5$ and $\tau_\omega=0.075<0.1$.

\answer{(iii)} The ex post policy is less ambitious ($\tau<\tau^*$), because part of its benefit accrues to foreigners; anticipating it, investors lend less ($I(\tau)$ is increasing).

\intuition A policy that raises the probability of success also raises pledgeable income and therefore borrowing capacity. Ex post, investment is sunk and that benefit is gone, so the government equates marginal cost only to the marginal output gain $R$. Commitment lets it internalise the investment response. When the gains of investors do not count fully (foreign lenders), the ex post temptation to under-provide is even stronger.

\pitfalls
\begin{itemize}
  \item Putting $\tau$ into the incentive constraint: it raises both success probabilities and cancels.
  \item Forgetting $I'(\tau)$ in the ex ante first-order condition, which makes $\tau^c=\tau^*$.
  \item Maximising the entrepreneur's payoff only: here welfare includes investors (all domestic).
\end{itemize}

\bookans Matches the book's answer (p.~622): $\gamma'(\tau^*)=R$, condition \eqref{r10e2-tauc}, and $\tau<\tau^*$ in (iii). The sign argument and the example are additions.
\refbib{Tirole \S16.3.2 ``Basics of Time Inconsistency in Corporate Finance'', pp.~546--549; \S3.4.}

% =============================================================================
\exsec{3. Imperfectly correlated liquidity shocks --- Tirole 15.3, p.~532}
%% BEGIN-STATEMENT 3
\stmtbegin
\textbf{Source:} Tirole, Ch. 15, Exercise 15.3, p.~532 (PDF 545) · Answer in book: yes, p.~621

\textbf{Exercise 15.3 (imperfectly correlated shocks).} This exercise extends the analysis of
Section 15.3 to allow for imperfect correlation among the shocks faced by the firms. As in Section
15.3.1, there is a mass 1 of \emph{ex ante} identical entrepreneurs. Each entrepreneur has a
constant-returns-to-scale project. An investment of size \(I\) at date 0 yields \(\rho_1 I\) at date
2, of which \(\rho_0 I\) is pledgeable, provided that the liquidity shock \(\rho I\) is met at date
1. \(\rho\) is equal to \(\rho_L\) with probability \((1 - \lambda)\) and \(\rho_H\) with
probability \(\lambda\), with \(\rho_L < \rho_0 < \rho_H < \rho_1\) and
\((1 - \lambda)(\rho_H - \rho_0) < 1\). As usual, entrepreneurs and investors are risk neutral, and
the latter demand a rate of return equal to 0 (see Figure 15.11).

Figure 15.11 (timeline, left to right; p.~532, PDF 545)

{\def\LTcaptype{none} % do not increment counter
\begin{longtable}[]{@{}
  >{\raggedright\arraybackslash}p{(\linewidth - 6\tabcolsep) * \real{0.2500}}
  >{\raggedright\arraybackslash}p{(\linewidth - 6\tabcolsep) * \real{0.2500}}
  >{\raggedright\arraybackslash}p{(\linewidth - 6\tabcolsep) * \real{0.2500}}
  >{\raggedright\arraybackslash}p{(\linewidth - 6\tabcolsep) * \real{0.2500}}@{}}
\toprule\noalign{}
\begin{minipage}[b]{\linewidth}\raggedright
Date 0
\end{minipage} & \begin{minipage}[b]{\linewidth}\raggedright
Date 1
\end{minipage} & \begin{minipage}[b]{\linewidth}\raggedright
Date 1 → 2
\end{minipage} & \begin{minipage}[b]{\linewidth}\raggedright
Date 2
\end{minipage} \\
\midrule\noalign{}
\endhead
\bottomrule\noalign{}
\endlastfoot
Representative entrepreneur borrows \(I - A\), invests \(I\). & Liquidity need \(\rho I\) realized.
Then either ``Continue'' or ``Liquidate''. & (If continue) Moral hazard. & Expected income
\(\rho_1 I\), of which \(\rho_0 I\) is pledgeable. \\
\end{longtable}
}

The new feature is that shocks are imperfectly correlated: for a fraction \(1 - \theta\) of
entrepreneurs, shocks are drawn independently (\(\theta = 0\) in Section 15.2.1). A fraction
\(\theta\) of entrepreneurs face the same shock, \(\rho_L\) with probability \((1 - \lambda)\) and
\(\rho_H\) with probability \(\lambda\) (\(\theta = 1\) in Section 15.3.1).

There is no outside store of value, and the long-term projects are the only investment projects
available to the corporate sector.

Show that the private sector is self-sufficient (i.e., the efficient allocation can be implemented
using the inside liquidity created by the long-term projects) if and only if
\(\theta \leqslant \theta^*\), where

\[(1 - \theta^*)(I - A) = \theta^*(\rho_H - \rho_0)I,\]

where \(I\) is independent of \(\theta\).

\paragraph{Context needed to solve it}

\emph{From Tirole §15.2.1 ``The Sufficiency Result'', pp.~518--520 (PDF 531--533):}

There are three periods, \(t = 0, 1, 2\). Investors are risk neutral and the market rate of interest
in the economy is 0. The economy is also populated by a large number---technically a continuum of
mass 1---of ex ante identical, risk-neutral entrepreneurs. The representative entrepreneur at date 0
has wealth \(A\), borrows \(I - A\), and invests \(I\). At date 1, a given firm faces liquidity
shock \(\rho I\), with

\[\rho = \begin{cases} \rho_L & \text{with probability } 1 - \lambda \text{ (healthy firm)},\\ \rho_H > \rho_L & \text{with probability } \lambda \text{ (firm in distress)}.\end{cases}\]

The firm can continue only if it finds funds to defray its liquidity shock; otherwise, it is
liquidated. We normalize the liquidation value at 0.

In the case of continuation, the firm's date-2 expected income is denoted by \(\rho_1 I\), of which
only \(\rho_0 I < \rho_1 I\) is pledgeable to investors (see Figure 15.1, identical to Figure 15.11
above).\textsuperscript{3}

Let us assume that

\[\rho_L < \rho_0 < \rho_H < \rho_1 \qquad (15.1)\]

and

\[(1 - \lambda)(\rho_H - \rho_L) < 1. \qquad (15.2)\]

{[}\ldots{]} Let us in a first step ignore this difficulty {[}\ldots{]} and assume that somehow
investors can commit to any probability \(x\) in \([0, 1]\) of continuation in the adverse state of
nature and that they do not demand an extra return for this (that is, they just want to recoup the
extra loss \(x(\rho_H - \rho_0)I\) induced by continuation in that state). {[}\ldots{]}

The investors' outlay at date 0 is equal to \(I - A\), their expected outlay at date 1 is
\(((1 - \lambda)\rho_L + \lambda\rho_H x)I\), and their expected income is
\((1 - \lambda + \lambda x)\rho_0 I\). Hence, the investors' breakeven constraint is

\[[1 + (1 - \lambda)\rho_L + \lambda\rho_H x]I - A \leqslant [1 - \lambda + \lambda x]\rho_0 I.\]

The efficient allocation is defined as the one that maximizes the representative entrepreneur's
utility subject to the constraint that investors break even at their IMRS. {[}\ldots{]} This
efficient allocation solves

\[\max_{\{I, x\}} \{(1 - \lambda + \lambda x)(\rho_1 - \rho_0)I\} \quad \text{s.t.} \quad [1 + (1 - \lambda)\rho_L + \lambda\rho_H x]I - A \leqslant [1 - \lambda + \lambda x]\rho_0 I.\]

{[}\ldots{]} the unit cost of effective investment is

\[c(x) \equiv \frac{1 + (1 - \lambda)\rho_L + \lambda\rho_H x}{1 - \lambda + \lambda x}.\]

The program yielding the efficient allocation becomes
\(\max_{\{x\}} \left\{\dfrac{(\rho_1 - \rho_0)A}{c(x) - \rho_0}\right\}\). The optimal \(x\) must
therefore minimize \(c(x)\), which, together with condition (15.2), implies that \(x = 1\).

Let us now write the investors' breakeven constraint for the policy of never liquidating:

\[[(1 - \lambda)(\rho_0 - \rho_L) + \lambda(\rho_0 - \rho_H)]I = I - A \quad\text{or}\quad (\rho_0 - \bar\rho)I = I - A,\]

where \(\bar\rho \equiv (1 - \lambda)\rho_L + \lambda\rho_H\) is the expected liquidity shock per
unit of investment. Lastly, we assume that \(\rho_1 > 1 + \bar\rho\) {[}\ldots{]}. To sum up, the
efficient allocation is given by

\[x = 1 \quad\text{and}\quad I = \frac{A}{(1 + \bar\rho) - \rho_0}.\]

\emph{15.2.1.3 The Sufficiency Result.} Next we assume that investors cannot pledge their future
earnings and therefore cannot directly commit to reinject cash at date 1 in the firm when they lose
money on this reinjection (\(\rho_H > \rho_0\)); we ask whether the efficient allocation can
nevertheless be implemented.

Assume that the shocks are drawn independently across firms, and so there is no macroeconomic
uncertainty. Because the corporate sector is a net borrower, \(I - A > 0\), the investors' breakeven
condition implies that investors' profit in the healthy firms (those facing a low shock at date 1)
more than offsets the loss that they incur in the others:

\[(1 - \lambda)(\rho_0 - \rho_L) > \lambda(\rho_H - \rho_0). \qquad (15.3)\]

We can define (gross) inside liquidity as the value, \((1 - \lambda)(\rho_0 - \rho_L)I\), of healthy
firms. Condition (15.3) states that the (gross) inside liquidity exceeds the net refinancing need,
\(\lambda(\rho_H - \rho_0)I\), of firms facing a high liquidity shock {[}\ldots{]}.

\emph{From Tirole footnote 3, p.~519 (PDF 532):} The wedge between date-2 value and pledgeable
income can, as in Chapter 5, be motivated by moral-hazard considerations: {[}\ldots{]}
\(\rho_1 \equiv p_H R\) and \(\rho_0 = p_H[R - B/\Delta p]\).

\emph{From Tirole §15.3.1 ``Aggregate Shocks and the Value of Outside Liquidity'', p.~523 (PDF
536):}

Suppose now that liquidity shocks are perfectly correlated across firms in the model of Section 15.2
{[}\ldots{]}. In this polar case, and in the absence of outside liquidity, firms cannot continue
when \(\rho = \rho_H > \rho_0\): \(x = 0\). {[}\ldots{]} Because all firms are valueless in the bad
state of nature, distressed firms cannot meet liquidity shocks by selling (even indirectly through a
financial intermediary) shares in healthy ones.
\stmtend
%% END-STATEMENT

\asked Show that inside liquidity (claims on the firms themselves) is enough to implement the efficient allocation exactly when the correlated fraction $\theta$ of firms is at most $\theta^*$, with $\theta^*$ given by the stated equation.

\approach The efficient allocation of Tirole \S15.2.1 continues all firms ($x=1$) and invests $I=A/(1+\bar\rho-\rho_0)$, independent of $\theta$. Inside liquidity is the value of healthy firms at date 1; it must cover the net refinancing need of distressed firms in \emph{every} aggregate state. Only the state in which the correlated block is hit is binding.

\notation
\begin{itemize}
  \item $\rho_L<\rho_0<\rho_H<\rho_1$: low / high liquidity shock, pledgeable and total income per unit; $\lambda$: probability of the high shock; $\bar\rho=(1-\lambda)\rho_L+\lambda\rho_H$.
  \item $\theta$: fraction of firms with a common shock; $1-\theta$: fraction with independent shocks.
\end{itemize}

\step{1}{Efficient allocation} From the context (\S15.2.1.2), the breakeven with $x=1$ is
\begin{align}
  [(1-\lambda)(\rho_0-\rho_L)-\lambda(\rho_H-\rho_0)]I &= I-A, \label{r10e3-be}\\
  I &= \frac{A}{1+\bar\rho-\rho_0}, \label{r10e3-I}
\end{align}
which does not involve $\theta$. ($x=1$ is optimal by (15.2), $(1-\lambda)(\rho_H-\rho_L)<1$; we take the section's allocation as the target.)

\step{2}{State where the common shock is high} The correlated block (mass $\theta$) and a fraction $\lambda$ of the independent block are distressed; each needs $(\rho_H-\rho_0)I$ beyond what it can raise by issuing claims:
\begin{align}
  \text{Need} &= [\theta+(1-\theta)\lambda](\rho_H-\rho_0)I. \label{r10e3-need}
\end{align}
Healthy firms are only in the independent block, mass $(1-\theta)(1-\lambda)$, each worth $(\rho_0-\rho_L)I$ net of its own shock:
\begin{align}
  \text{Supply} &= (1-\theta)(1-\lambda)(\rho_0-\rho_L)I. \label{r10e3-supply}
\end{align}

\step{3}{Use the breakeven} From \eqref{r10e3-be}, $(1-\lambda)(\rho_0-\rho_L)I=(I-A)+\lambda(\rho_H-\rho_0)I$. Substituting into \eqref{r10e3-supply}, Supply $\geq$ Need becomes
\begin{align}
  (1-\theta)\bigl[(I-A)+\lambda(\rho_H-\rho_0)I\bigr] &\geq [\theta+(1-\theta)\lambda](\rho_H-\rho_0)I \nonumber\\
  (1-\theta)(I-A)+(1-\theta)\lambda(\rho_H-\rho_0)I &\geq \theta(\rho_H-\rho_0)I+(1-\theta)\lambda(\rho_H-\rho_0)I \nonumber\\
  (1-\theta)(I-A) &\geq \theta(\rho_H-\rho_0)I. \label{r10e3-cond}
\end{align}

\step{4}{Threshold} The left side of \eqref{r10e3-cond} decreases in $\theta$ and the right side increases (with $I$ fixed by \eqref{r10e3-I}), so \eqref{r10e3-cond} holds iff $\theta\leq\theta^*$ with
\begin{align}
  (1-\theta^*)(I-A) = \theta^*(\rho_H-\rho_0)I \iff \theta^*=\frac{I-A}{(I-A)+(\rho_H-\rho_0)I}. \label{r10e3-thstar}
\end{align}

\step{5}{The other aggregate state} If the common shock is low, the need is $(1-\theta)\lambda(\rho_H-\rho_0)I$ and the supply $[\theta+(1-\theta)(1-\lambda)](\rho_0-\rho_L)I\geq(1-\theta)(1-\lambda)(\rho_0-\rho_L)I$. By (15.3), $(1-\lambda)(\rho_0-\rho_L)>\lambda(\rho_H-\rho_0)$ (net borrowing), so supply exceeds need for every $\theta$: no shortage there.

\step{6}{Example} $\lambda=0.3$, $\rho_L=0.2$, $\rho_0=0.6$, $\rho_H=1$, $\rho_1=1.8$, $A=1$: $\bar\rho=0.44$, $I=1/0.84=1.190476$, $\theta^*=0.190476/(0.190476+0.476190)=2/7=0.285714$; at $\theta^*$ supply = need $=0.238095$. Here $(1-\lambda)(\rho_H-\rho_L)=0.56<1$.

\answer{} Self-sufficiency iff $(1-\theta)(I-A)\geq\theta(\rho_H-\rho_0)I$, i.e.\ $\theta\leq\theta^*=(I-A)/[(I-A)+(\rho_H-\rho_0)I]$, with $I=A/(1+\bar\rho-\rho_0)$. (Liquidity must be pooled, e.g.\ via credit lines, as in \S15.2.2.)

\intuition Inside liquidity works by moving value from healthy to distressed firms. Independent shocks diversify away; a common shock cannot be diversified, because the firms that would provide liquidity are hit at the same time. The more of the economy is exposed to the common shock (higher $\theta$), the thinner the pool of healthy firms in the bad state, until it cannot cover the refinancing need.

\pitfalls
\begin{itemize}
  \item Checking only the average state: the binding constraint is the state in which the correlated block is hit.
  \item Letting $I$ depend on $\theta$: the efficient allocation is the same for every $\theta$.
  \item Counting the gross value $\rho_0 I$ of healthy firms instead of the net $(\rho_0-\rho_L)I$ after their own shock.
\end{itemize}

\bookans Matches the book's answer (p.~621), which writes the same need, supply and breakeven and arrives at $(1-\theta)(I-A)\geq\theta(\rho_H-\rho_0)I$. Note: the statement assumes $(1-\lambda)(\rho_H-\rho_0)<1$, while $x=1$ in \S15.2.1 rests on (15.2), $(1-\lambda)(\rho_H-\rho_L)<1$; the solution uses the section's allocation.
\refbib{Tirole \S15.2.1, pp.~518--521; \S15.3.1, p.~523.}

% =============================================================================
\exsec{4. Property-rights institutions and international finance --- Tirole Review Problem 5, pp.~628--629}
%% BEGIN-STATEMENT 4
\stmtbegin
\textbf{Source:} Tirole, Part VII, Review Problem 5, pp.~628--629 (PDF 641) · Answer in book: yes,
p.~635

\textbf{Review Problem 5 (property rights institutions and international finance).} Consider a
country with a continuum of identical firms (of mass 1). The representative firm is described as in
Section 3.4. That is, it has initial wealth \(A\) and has a variable-investment project. As usual,
let

\[\rho_0 \equiv p_H\left(R - \frac{B}{\Delta p}\right) < 1 < \rho_1 \equiv p_H R.\]

To finance their investment, the domestic firms must borrow from domestic residents and from foreign
investors. Domestic residents have limited savings \(S_D < \rho_0 A/(1 - \rho_0)\). Foreign
investors have unlimited amounts of money to lend at the market rate of return. The rate of return
demanded by foreigners and domestic residents is 0.

After the financing has been secured, the country's government chooses a tax rate \(t \geqslant 0\)
on income received by \emph{investors}. This tax rate does not apply to the entrepreneurs and does
not discriminate between domestic and foreign investors.

To close the model, assume that the government transforms tax proceeds \(t\mathcal{I}\) (for an
amount of investor income \(\mathcal{I}\)) into \(B_0(t)\mathcal{I}\), where

\[B_0' > 0, \quad B_0'' < 0, \quad\text{and}\quad B_0'(0) = 1.\]

(That is, tax collection is wasteful here.) Assume that these benefits \(B_0(t)\mathcal{I}\) are
returned to entrepreneurs in proportion to the tax income collected. Thus the representative
entrepreneur's equilibrium NPV is

\[(\rho_1 - 1)I - t^*\rho_0 I + B_0(t^*)\rho_0 I.\]

The government maximizes the sum of the welfares of the entrepreneurs and of the domestic investors,
and puts no weight on foreign investors (see Figure 1).

Figure 1 (timeline, left to right; p.~629, PDF 642)

{\def\LTcaptype{none} % do not increment counter
\begin{longtable}[]{@{}
  >{\raggedright\arraybackslash}p{(\linewidth - 6\tabcolsep) * \real{0.2500}}
  >{\raggedright\arraybackslash}p{(\linewidth - 6\tabcolsep) * \real{0.2500}}
  >{\raggedright\arraybackslash}p{(\linewidth - 6\tabcolsep) * \real{0.2500}}
  >{\raggedright\arraybackslash}p{(\linewidth - 6\tabcolsep) * \real{0.2500}}@{}}
\toprule\noalign{}
\begin{minipage}[b]{\linewidth}\raggedright
Step 1
\end{minipage} & \begin{minipage}[b]{\linewidth}\raggedright
Step 2
\end{minipage} & \begin{minipage}[b]{\linewidth}\raggedright
Step 3
\end{minipage} & \begin{minipage}[b]{\linewidth}\raggedright
Step 4
\end{minipage} \\
\midrule\noalign{}
\endhead
\bottomrule\noalign{}
\endlastfoot
Financing: \(S_D\) by domestic investors, \(I - A - S_D\) by foreign investors. & Moral hazard. &
Government sets tax rate \(t\). & Incomes accrues. Returns paid to investors. Taxes paid to
government. \\
\end{longtable}
}

\begin{enumerate}
\def\labelenumi{(\roman{enumi})}
\item
  Solve for a rational expectation equilibrium \((I^*, t^*, \theta^*)\), where \(I^*\) is the
  representative entrepreneur's investment, \(t^*\) is the equilibrium tax rate, and \(\theta^*\),
  equal to \(S_D/(I^* - A)\), is the fraction of external financing brought by domestic residents.
\item
  How does the entrepreneur's welfare change with domestic savings \(S_D\)?
\item
  What tax rate would prevail if the government were able to commit on the tax rate before the
  financing stage?
\item
  How would your answer to question (i) change if the government were still unable to commit to a
  tax rate and furthermore could discriminate between domestic and foreign investors?
\end{enumerate}

\paragraph{Context needed to solve it}

\emph{From Tirole §3.4.1 ``The Continuous-Investment Model'', p.~127 (PDF 140):}

{[}\ldots{]} we now assume that there are constant returns to scale in the investment technology. An
investment \(I \in [0, \infty)\) yields income \(RI\), proportional to \(I\), in the case of
success, and 0 in the case of failure. The borrower's private benefit from misbehaving is also taken
to be proportional to investment. As before, the borrower has a choice between behaving, in which
case she derives no private benefit and the probability of success is \(p_H\), and misbehaving, that
is, enjoying private benefit \(BI\), and reducing the probability of success to
\(p_L = p_H - \Delta p < p_H\). {[}\ldots{]}

The borrower initially has cash \(A\), and must therefore borrow \(I - A\) to finance a project of
size \(I\). A loan agreement specifies that the lenders (who as before are assumed to make no
profit) and the borrower receive 0 each in the case of failure, and \(R_l\) and \(R_b\),
respectively, in the case of success, where \(R_l + R_b = RI\).

As in Section 3.2, we assume that investment has positive NPV (net present value), here per unit of
investment, if the borrower behaves,

\[p_H R > 1, \qquad (3.7)\]

but negative NPV otherwise,

\[1 > p_L R + B, \qquad (3.8)\]

so that unless one can control the agency problem the investment cannot be funded. We also make an
assumption that guarantees that the equilibrium investment is finite:

\[p_H R < 1 + \frac{p_H B}{\Delta p}. \qquad (3.9)\]

\emph{From Tirole §3.4.2 ``The Lenders' Credit Analysis'', pp.~127--128 (PDF 140--141):}

Following the steps of Section 3.2, the incentive compatibility and the breakeven conditions are

\[(\Delta p)R_b \geqslant BI \qquad (IC_b)\]

and

\[p_H(RI - R_b) \geqslant I - A. \qquad (IR_l)\]

In equilibrium, competitive lenders make no profit on the contract that is most advantageous for the
borrower; the borrower's net utility is therefore equal to the social surplus brought about by the
investment:

\[U_b = (p_H R - 1)I. \qquad (3.10)\]

From (3.10) it is optimal for the borrower to invest as much as possible. {[}\ldots{]} Substituting
\((IC_b)\) into \((IR_l)\), we obtain

\[I \leqslant kA, \qquad (3.11)\]

where

\[k = \frac{1}{1 - p_H(R - B/\Delta p)} > 1. \qquad (3.12)\]
\stmtend
%% END-STATEMENT

\asked (i) The rational-expectations equilibrium: investment $I^*$, the tax rate $t^*$ the government sets after financing, and the domestic share $\theta^*$ of external finance. (ii) How entrepreneurs' welfare moves with domestic savings $S_D$. (iii) The tax rate under commitment. (iv) The equilibrium when the government can tax foreign and domestic investors differently.

\approach The firm is the variable-investment model (Tirole \S3.4): investors receive the pledgeable income $\rho_0I$. Taxing investor income ex post is a hold-up: it is a transfer from foreigners (not counted) and a costly transfer from domestic investors. Anticipated, it lowers what investors expect and thus borrowing capacity (the property-rights problem of Tirole \S16.3).

\notation
\begin{itemize}
  \item $\rho_0=p_H(R-B/\Delta p)<1<\rho_1=p_HR$; $A$: wealth; $I$: investment; $t$: tax rate on investor income.
  \item $\mathcal I=\rho_0I$: investor income; $B_0(t)\mathcal I$: benefit returned to entrepreneurs, $B_0'>0$, $B_0''<0$, $B_0'(0)=1$, and $B_0(0)=0$ (no tax, no benefit), so $B_0(t)<t$ for $t>0$.
  \item $S_D$: domestic savings; $\theta=S_D/(I-A)$: domestic share of investor claims.
\end{itemize}

\step{1}{Borrowing capacity for an expected tax $t$} The entrepreneur's incentive constraint $(\Delta p)R_b\geq BI$ is untouched (the tax hits investors only), so investors' pre-tax income is $\rho_0I$ and their after-tax income $(1-t)\rho_0I$. Breakeven at rate 0:
\begin{align}
  (1-t)\rho_0I &= I-A \nonumber\\
  I(t) &= \frac{A}{1-(1-t)\rho_0}, \qquad I(t)-A=\frac{(1-t)\rho_0A}{1-(1-t)\rho_0}. \label{r10e4-I}
\end{align}
$I(t)$ decreases in $t$.

\step{2}{Ex post tax} With $I$ and $\theta$ fixed, domestic surplus is entrepreneurs' payoff $(\rho_1-\rho_0)I+B_0(t)\rho_0I$ plus domestic investors' $\theta(1-t)\rho_0I$ (minus constants):
\begin{align}
  \frac{\partial}{\partial t}\Bigl[B_0(t)\rho_0I+\theta(1-t)\rho_0I\Bigr] &= [B_0'(t)-\theta]\rho_0I=0 \;\Longrightarrow\; B_0'(t^*)=\theta^*. \label{r10e4-foc}
\end{align}
The second-order condition is $B_0''<0$.

\step{3}{Fixed point} Rational expectations require \eqref{r10e4-foc} with $\theta^*=S_D/(I(t^*)-A)$:
\begin{align}
  B_0'(t^*) = \frac{S_D}{I(t^*)-A} = \frac{S_D[1-(1-t^*)\rho_0]}{(1-t^*)\rho_0A}, \qquad I^*=\frac{A}{1-(1-t^*)\rho_0}. \label{r10e4-fp}
\end{align}
At $t=0$ the left side is $1$ and the right side is $S_D(1-\rho_0)/(\rho_0A)<1$ by the assumption $S_D<\rho_0A/(1-\rho_0)$; the right side increases in $t$ (to $\infty$ as $t\to1$) and the left side decreases. Hence a unique $t^*\in(0,1)$: the government taxes, because foreigners hold part of the claims.

\step{4}{Example} $B_0(t)=\ln(1+t)$ (so $B_0'(t)=1/(1+t)$), $\rho_0=0.6$, $\rho_1=1.5$, $A=1$, $S_D=0.5$. At $t=1/3$: $I-A=(\tfrac23)(0.6)/(1-0.4)=\tfrac23$, $\theta=0.5/\tfrac23=0.75=1/(1+\tfrac13)$, so $t^*=0.333333$, $\theta^*=0.75$, $I^*=1.666667$. Entrepreneurs' NPV (statement): $[(\rho_1-1)-(t^*-B_0(t^*))\rho_0]I^*=[0.5-(0.333333-0.287682)(0.6)](1.666667)=0.787682$.

\answer{(i)} $B_0'(t^*)=\theta^*=S_D/(I^*-A)$ and $I^*=A/[1-(1-t^*)\rho_0]$, with $t^*\in(0,1)$.

\step{5}{More domestic savings} Raising $S_D$ shifts the right side of \eqref{r10e4-fp} up, so $t^*$ falls; then $I^*$ rises by \eqref{r10e4-I} and $\theta^*=B_0'(t^*)$ rises. Welfare
\begin{align}
  U_b(t) = \bigl[(\rho_1-1)-(t-B_0(t))\rho_0\bigr]\,I(t) \label{r10e4-Ub}
\end{align}
decreases in $t$: the bracket falls ($\frac{d}{dt}(t-B_0)=1-B_0'\geq0$) and $I(t)$ falls. So $U_b$ rises with $S_D$. Example with $S_D=0.6$: $t^*=0.276984$, $\theta^*=0.783095$, $I^*=1.766190$, $U_b=0.848673>0.787682$.

\answer{(ii)} Higher $S_D$ lowers $t^*$, raises $\theta^*$ and $I^*$, and raises entrepreneurs' welfare.

\step{6}{Commitment} Choosing $t$ before financing, the government maximises \eqref{r10e4-Ub}. Its derivative is $-(1-B_0'(t))\rho_0I(t)+[\cdots]I'(t)<0$ for all $t\geq0$ (the first term is $\leq0$, the second $<0$), so $t=0$ and $U_b=(\rho_1-1)A/(1-\rho_0)=0.5/0.4=1.25$ in the example.

\answer{(iii)} $t=0$.

\step{7}{Discrimination} Ex post the government taxes foreign income at $t_F$ and domestic at $t_D$. Foreign income does not count: maximise $B_0(t_F)$, so $t_F$ is as high as possible ($t_F=1$, full expropriation). Domestic: maximise $B_0(t_D)-t_D$, derivative $B_0'(t_D)-1\leq0$, so $t_D=0$. Foreigners anticipate and do not lend; only domestic savings finance firms:
\begin{align}
  I^*-A = S_D. \label{r10e4-disc}
\end{align}
Domestic investors receive the untaxed $\rho_0I^*=\rho_0(A+S_D)>S_D$ (equivalent to $S_D<\rho_0A/(1-\rho_0)$): a positive return. Example: $I^*=1.5$, domestic investors get $0.9$ for $0.5$ lent.

\answer{(iv)} Foreigners would be fully taxed, so none lend: $I^*=A+S_D$, domestic investors untaxed and earning a positive return.

\intuition Weak property-rights institutions let the government expropriate investors once the money is in. Since foreigners carry part of the tax and do not vote, the temptation grows with foreign financing; investors see it coming and lend less. Domestic savings act as a commitment device: the more domestic residents hold the claims, the more the government hurts its own constituents by taxing.

\pitfalls
\begin{itemize}
  \item Letting the tax enter the entrepreneur's incentive constraint: it applies to investors only.
  \item Solving the government's problem with $I$ variable in (i): ex post, $I$ is sunk.
  \item In (iv), keeping foreign lending: rational foreigners anticipate full expropriation.
\end{itemize}

\bookans Matches the book's answer (p.~635): $B_0'(t^*)=\theta^*=S_D/(I^*-A)$, $I^*=A/[1-(1-t^*)\rho_0]$; $S_D$ raises $\theta^*$ and $I^*$ and lowers $t^*$; commit to $t=0$; with discrimination $I^*-A=S_D$. Existence/uniqueness and the example are additions.
\refbib{Tirole \S3.4, pp.~127--128; \S16.3 ``Property Rights Institutions'', pp.~544--551.}

% =============================================================================
\exsec{5. Borrowing abroad --- Tirole 16.1, pp.~563--564}
%% BEGIN-STATEMENT 5
\stmtbegin
\textbf{Source:} Tirole, Ch. 16, Exercise 16.1, pp.~563--564 (PDF 576) · Answer in book: yes, p.~622

\textbf{Exercise 16.1 (borrowing abroad).} Consider a small country with a mass 1 of identical
entrepreneurs. There is a single (tradable) good. The representative entrepreneur has initial wealth
\(A\) and a variable-investment constant-returns-to-scale project. A project of size
\(I \in [0, \infty)\) at date 1 yields at date 2 verifiable revenue \(RI\) with probability \(p\)
and 0 with probability \(1 - p\). The probability \(p\) is not subject to moral hazard. There is
moral hazard, though: instead of investing \(I\) in the firm, the entrepreneur can invest it abroad
and get private return \(\mu I\), where \(\mu < 1\). The investors are unable to seize the return
from this alternative investment. Everyone is risk neutral, has discount factor 1 (i.e., has utility
equal to the undiscounted sum of consumptions at dates 1 and 2), and the entrepreneur is protected
by limited liability.

One will assume that

\[pR > 1 > pR - \mu.\]

\begin{enumerate}
\def\labelenumi{(\roman{enumi})}
\item
  Compute the representative entrepreneur's borrowing capacity and utility. Show that the outcome is
  the same as in a situation in which the entrepreneur cannot divert funds and invest them abroad,
  but can enjoy a private benefit per unit of investment \(B = \mu\). Explain why.
\item
  Adopt the convention that the payment to investors, \(R_l\), is a debt payment. Suppose that the
  entrepreneurs' projects are independent and that the government imposes a per-unit-of-income tax
  on successful projects and offers a guarantee/compensation \(\sigma\) on private debt (so
  \(\tau RI\) is the tax on successful projects and \(\sigma R_l\) is the investors' payoff in the
  case of bankruptcy). Show that the borrowing capacity and entrepreneur utility are the same as in
  (i).
\end{enumerate}

In contrast, compute the impact on entrepreneurs when the government starts at date 1 with an
inherited public debt outstanding to foreign lenders equal to \(D\) (\(\leqslant A\)) per
entrepreneur and must finance it through an income tax on successful projects.

\begin{enumerate}
\def\labelenumi{(\roman{enumi})}
\setcounter{enumi}{2}
\tightlist
\item
  Coming back to question (i), suppose that the government can through its governance institutions
  or other policies affect the return \(\mu\) on investments abroad. There are two levels
  \(\mu_L < \mu_H\) (where both levels satisfy the conditions in (i)). The choice between the two
  levels involves no cost (but affects behavior!). The government's objective function is to
  maximize the representative entrepreneur's welfare.
\end{enumerate}

Assuming that all borrowing is foreign borrowing, what is the representative entrepreneur's utility
when

\begin{enumerate}
\def\labelenumi{(\alph{enumi})}
\item
  the government can commit to \(\mu\) before foreign investors invest;
\item
  the government chooses \(\mu\) after they have invested (but before the entrepreneurs select their
  action)?
\end{enumerate}

\begin{enumerate}
\def\labelenumi{(\roman{enumi})}
\setcounter{enumi}{3}
\tightlist
\item
  Suppose now that the output \(RI\) (in the case of success) is in terms of a nontradable good (but
  the endowment \(A\) and the investment \(I\) are in tradable goods). Another sector of the economy
  (the ``export sector'') will receive \(\mathcal{R}\) in tradable goods at date 2. All domestic
  agents have utility from date-2 consumptions \(c\) and \(c^*\) of nontradable and tradable goods
  equal to \(c + c^*\) (so the two goods are perfect substitutes for domestic residents, while
  foreigners consume only the tradable good). Define the date-2 exchange rate \(e \geqslant 1\) as
  the price of tradables in terms of nontradables. Compute the borrowing capacity and the exchange
  rate. (One will, for example, assume that funds fraudulently invested abroad cannot be reimported
  and must be consumed abroad. So they yield \(\mu\) rather than \(e\mu\).)
\end{enumerate}

\paragraph{Context needed to solve it}

None: the statement is self-contained. For the comparison in (i), the basic model is the one copied
in exercise 4 above (Tirole §3.4: \((\Delta p)R_b \geqslant BI\), \(p_H(RI - R_b) \geqslant I - A\),
\(I \leqslant A/[1 - p_H(R - B/\Delta p)]\), \(U_b = (p_H R - 1)I\)).
\stmtend
%% END-STATEMENT

\asked (i) Borrowing capacity and utility when the entrepreneur can divert funds abroad, and why this is the basic model with $B=\mu$. (ii) Whether a balanced tax-plus-guarantee scheme changes anything, and the effect of inherited public debt $D$ financed by taxing successful projects. (iii) The utility when the government chooses $\mu$ before vs.\ after foreign investors lend. (iv) Borrowing capacity and the exchange rate when output is nontradable.

\approach The diversion option is an incentive constraint: the entrepreneur must earn at home at least $\mu I$. Everything follows from that constraint and the investors' breakeven, i.e.\ from pledgeable income $(pR-\mu)I$ (Gertler and Rogoff 1990, cited in Tirole \S16.7). In (iv) the country can pay foreigners only with tradables, so export revenue becomes international collateral.

\notation
\begin{itemize}
  \item $p$: success probability; $R$: return per unit; $\mu<1$: private return of funds invested abroad; $A$: wealth; $I$: investment.
  \item $R_l$: payment to investors if success; $\tau$, $\sigma$: tax rate on successful income and guarantee rate on debt; $D$: public debt per entrepreneur.
  \item $\Rc$: export sector's tradable income; $e\geq1$: price of tradables in nontradables.
\end{itemize}

\step{1}{Constraints} Diverting yields $\mu I$ and nothing at home (investors seize nothing abroad); investing yields $p(RI-R_l)$. Incentive compatibility and breakeven:
\begin{align}
  p(RI-R_l) &\geq \mu I, \label{r10e5-ic}\\
  pR_l &= I-A. \label{r10e5-be}
\end{align}
Substituting \eqref{r10e5-be} into \eqref{r10e5-ic}:
\begin{align}
  pRI-(I-A) &\geq \mu I \nonumber\\
  I[1-(pR-\mu)] &\leq A \nonumber\\
  I &\leq \frac{A}{1-(pR-\mu)}. \label{r10e5-I}
\end{align}
The denominator is positive since $pR-\mu<1$. Investors break even, so $U_b$ is the NPV $(pR-1)I$, positive since $pR>1$, so $I$ is at its maximum:
\begin{align}
  U_b = \frac{pR-1}{1-(pR-\mu)}A. \label{r10e5-Ub}
\end{align}

\step{2}{Equivalence} In the basic model the pledgeable income per unit is $p_H(R-B/\Delta p)$. With $p_H=p$, $p_L=0$ (diversion kills the domestic project) and $B=\mu$, $\Delta p=p$ and $p(R-\mu/p)=pR-\mu$, the same as in \eqref{r10e5-I}.

\answer{(i)} $I=A/[1-(pR-\mu)]$, $U_b=(pR-1)A/[1-(pR-\mu)]$. Diversion is a private benefit $\mu$ per unit with $p_L=0$: investors get nothing from the diverted funds.

\step{3}{Balanced tax and guarantee} The constraints become
\begin{align}
  p[(1-\tau)RI-R_l] &\geq \mu I, \label{r10e5-ic2}\\
  pR_l+(1-p)\sigma R_l &= I-A, \label{r10e5-be2}\\
  p\tau RI &= (1-p)\sigma R_l \quad\text{(government budget)}. \label{r10e5-gb}
\end{align}
Using \eqref{r10e5-gb} in \eqref{r10e5-be2}: $pR_l+p\tau RI=I-A$. From \eqref{r10e5-ic2}, $pR_l+p\tau RI\leq pRI-\mu I$. So $I-A\leq(pR-\mu)I$: the same bound \eqref{r10e5-I}, and the same NPV. (In the example of Step~7, with $\tau=0.1$: $R_l=15$, $\sigma=0.2$, and investors get $pR_l+(1-p)\sigma R_l=7.5+1.5=9=I-A$.)

\step{4}{Inherited debt} Now $p\tau RI=D$ leaves the country. With \eqref{r10e5-ic2} binding and \eqref{r10e5-be}:
\begin{align}
  p(1-\tau)RI-(I-A) &\geq \mu I \nonumber\\
  pRI-D-I+A &\geq \mu I \nonumber\\
  I &= \frac{A-D}{1-(pR-\mu)}. \label{r10e5-ID}
\end{align}
The entrepreneur consumes $\mu I=\mu(A-D)/[1-(pR-\mu)]$; her gain over the net-of-debt wealth $A-D$ is $(pR-1)(A-D)/[1-(pR-\mu)]$, and relative to (i) her consumption falls by $\mu D/[1-(pR-\mu)]>D$ (since $pR>1$).

\answer{(ii)} The balanced scheme leaves $I$ and $U_b$ unchanged. Public debt $D$ acts like a fall of net worth to $A-D$: $I=(A-D)/[1-(pR-\mu)]$ and each unit of debt costs entrepreneurs $\mu/[1-(pR-\mu)]>1$ units.

\step{5}{Choosing $\mu$} By \eqref{r10e5-Ub}, $U_b$ falls with $\mu$. (a) With commitment the government picks $\mu_L$. (b) Without commitment, suppose investors lend expecting $\mu_L$, with $p(RI-R_l)=\mu_LI$. Ex post, $\mu_H$ gives the entrepreneur $\max\{p(RI-R_l),\mu_HI\}=\mu_HI>\mu_LI$; foreign investors' loss does not count, so the government picks $\mu_H$. Investors anticipate it and lend as if $\mu=\mu_H$.

\answer{(iii)} (a) $U_b=(pR-1)A/[1-(pR-\mu_L)]$; (b) $U_b=(pR-1)A/[1-(pR-\mu_H)]$, lower.

\step{6}{Nontradable output} Investors are paid in nontradables $pR_l$ (aggregate, independent projects), which must be turned into tradables. The only date-2 tradables are the export sector's $\Rc$; domestic residents value both goods equally, so at $e>1$ they sell all their tradables. Tradable market and breakeven in tradables:
\begin{align}
  e\Rc &= pR_l, \label{r10e5-fx}\\
  \frac{pR_l}{e} &= I-A, \label{r10e5-be3}\\
  p(RI-R_l) &\geq \mu I. \label{r10e5-ic3}
\end{align}
From \eqref{r10e5-fx}--\eqref{r10e5-be3}, $I-A=\Rc$. From \eqref{r10e5-be3} and \eqref{r10e5-ic3}: $pRI-e(I-A)\geq\mu I$, i.e.\ $I\leq A/[1-(pR-\mu)/e]$. Equating the two,
\begin{align}
  I &= \Rc+A = \frac{A}{1-(pR-\mu)/e}, \label{r10e5-I4}\\
  1-\frac{pR-\mu}{e} &= \frac{A}{A+\Rc} \;\Longrightarrow\; e=\left(1+\frac{A}{\Rc}\right)(pR-\mu). \label{r10e5-e}
\end{align}
This is valid iff $e\geq1$, i.e.\ $(1+A/\Rc)(pR-\mu)\geq1$; otherwise $e=1$ and $I$ is as in (i) (tradables are in excess supply).

\step{7}{Example} $p=0.5$, $R=3$, $\mu=0.6$, $A=1$: $pR-\mu=0.9$, $I=1/0.1=10$, $U_b=0.5(10)=5$. With $D=0.2$: $I=8$, book's gain $0.5(8)=4$, consumption loss $0.6(0.2)/0.1=1.2$. With $\mu_H=0.8$: $U_b=0.5/0.3=1.666667$. With $\Rc=2$: $e=1.5(0.9)=1.35$, $I=3=1/(1-0.9/1.35)$, $pR_l=e(I-A)=2.7=e\Rc$.

\answer{(iv)} $I=A+\Rc$ and $e=(1+A/\Rc)(pR-\mu)$ when this is $\geq1$; otherwise $e=1$ and $I=A/[1-(pR-\mu)]$.

\intuition The country can borrow only against what it can credibly hand over. Funds that can be spirited abroad behave like a private benefit; public debt owed abroad eats into net worth with a multiplier; a government that cannot commit to good institutions erodes pledgeable income; and when output is nontradable, export earnings are the binding collateral --- borrowing capacity is fixed by $\Rc$ and the exchange rate adjusts.

\pitfalls
\begin{itemize}
  \item Treating the guarantee scheme as a subsidy: the government budget \eqref{r10e5-gb} makes it a pure transfer between investors and taxpayers.
  \item Reading the book's $(pR-1)(A-D)/[\cdot]$ as consumption: it is the gain over $A-D$; consumption is $\mu I$.
  \item In (iv), valuing diverted funds at $e\mu$: the statement says they are consumed abroad and yield $\mu$.
\end{itemize}

\bookans Matches the book's answer (p.~622) in every part: $I=A/[1-(pR-\mu)]$, the same pledgeable income in (ii), $I=(A-D)/[1-(pR-\mu)]$ and $U_b=(pR-1)(A-D)/[1-(pR-\mu)]$ (gain measured over $A-D$, see Step~4), $\mu=\mu_H$ without commitment, $e\Rc=pR_l$, $I=\Rc+A$, $e\geq1\iff(1+A/\Rc)(pR-\mu)\geq1$.
\refbib{Tirole \S3.4; \S16.3 (property rights, time inconsistency), pp.~544--551; \S15.2.2 (international collateral, Caballero--Krishnamurthy), p.~522; Gertler and Rogoff (1990).}

% =============================================================================
\exsec{6. Shleifer--Vishny with decreasing returns --- Tirole 14.1, p.~515}
%% BEGIN-STATEMENT 6
\stmtbegin
\textbf{Source:} Tirole, Ch. 14, Exercise 14.1, p.~515 (PDF 528) · Answer in book: no

\textbf{Exercise 14.1 (investment externalities in an industry with decreasing returns to scale).}
Suppose that the entrepreneur's limited attention, say, induces decreasing returns to scale. Income
in the case of success is \(R(I)\), where \(R' > 0\), \(R'' < 0\), \(R'(0) = \infty\),
\(R'(\infty) = 0\). Redo the analysis of the Schleifer--Vishny model with this modification, and
determine the sign of the investment externality.

\paragraph{Context needed to solve it}

\emph{From Tirole §14.2.2 ``Industry-Wide Shocks and Distress Sales: The Shleifer--Vishny Model'',
pp.~500--503 (PDF 513--516):}

This subsection, building on Shleifer and Vishny (1992), endogenizes the resale price in the
two-firm, continuous-investment version of the model of Section 4.3.1. Let us restate the key
ingredients of this model.

\emph{Investment and redeployability.} There are two firms in the industry. The ``industry'' is here
defined as a group of symmetric firms using the same equipment/assets. For simplicity, we assume
that the two firms do not compete in the same product market (see below for a discussion of this
hypothesis). Each firm is run by an entrepreneur, who has initial cash \(A_i\).

Initially, firm \(i\) invests \(I_i\), and therefore borrows \(I_i - A_i\) from some ex ante
competitive lender (``lender \(i\)''). Then, there is a costless ``learning period.'' At the end of
the learning period, each firm learns whether it is ``productive'' (which has probability \(x\)), or
``unproductive,'' i.e., ``in distress'' (which has probability \(1 - x\)). Being unproductive means
that the firm will always be unsuccessful regardless of whether the entrepreneur behaves. For
example, there may be no demand for the firm's output. Its assets are then useless if left in place.
A productive firm is described as in the variable-investment model of Section 3.4 (which,
incidentally, corresponds to the case \(x = 1\)).

If both firms are productive, each manages its initial investment. Firm \(i\)'s profit is either 0
or \(RI_i\). Borrower \(i\)'s private benefit is 0 (if she behaves) or \(BI_i\) (if she does not).
The associated probabilities of success are \(p_H\) and \(p_L\), respectively.

If firm \(j\) is in distress, it sells its assets, which now have no internal
use.\textsuperscript{10} We assume that potential buyers outside the industry do not have the
knowledge to operate these assets. Only firm \(i\) (if it itself is not in distress) can buy it.
There has been no initial contract that would specify the transfer price in the case of distress.
Rather, this transfer price is determined through bargaining after distress occurs. We will later
determine the per-unit transfer price \(P\). The entrepreneur in firm \(i\) then manages
\((I_1 + I_2)\) units of assets, and obtains private benefit 0 (if she behaves) or \(B(I_1 + I_2)\)
(if she does not). Similarly, the income is either 0 or \(R(I_1 + I_2)\). Probabilities of success
are \(p_H\) if the entrepreneur behaves and \(p_L\) if she does not. Firm \(i\) has just grown
bigger through the acquisition.

As usual, we assume that in the absence of adverse shock (\(x = 1\)), projects are viable only if
the borrower behaves,

\[\rho_1 \equiv p_H R > 1 > p_L R + B, \qquad (14.1)\]

and we make a further assumption guaranteeing that loans are finite (even for \(x\) close to 1):

\[p_H R < 1 + \frac{p_H B}{\Delta p} \quad\text{or}\quad \rho_0 \equiv p_H\left(R - \frac{B}{\Delta p}\right) < 1. \qquad (14.2)\]

\emph{Loan agreements.} Lender \(i\) and entrepreneur \(i\) secretly sign at the start a loan
agreement specifying the amount of the loan \(I_i - A_i\) and the stake \(R_{bi}\) of entrepreneur
\(i\) in the case of success (in the absence of purchase of firm \(j\)'s assets). {[}\ldots{]}

\emph{Summary of timing.} The timing is summarized in Figure 14.1 (where ``MH\(i\)'' stands for
``moral hazard in firm \(i\)'').

Figure 14.1 (two parallel timelines, one per firm, left to right; p.~501, PDF 514): Firm 1: \(I_1\)
→ Loan agreement 1 → Learning period: No distress (prob. \(x\)) or Distress (prob. \(1 - x\)) →
Transfer of asset (from a firm in distress to the other) → MH1 → Outcome 1. Firm 2: \(I_2\) → Loan
agreement 2 → Learning period: No distress (prob. \(x\)) or Distress (prob. \(1 - x\)) → Transfer of
asset → MH2 → Outcome 2.

\emph{Correlation of shocks.} The shocks affecting the demands for the two products may be
correlated. We allow for an arbitrary level of correlation. The conditional probabilities (given
firm \(i\)'s state) that firm \(j\) is productive or in distress are stated in Table 14.1. For
consistency, the parameters must be such that the probability that firm \(j\) is productive is
\(x\):

\[x\mu + (1 - x)(1 - \nu) = x \iff x(1 - \mu) = (1 - x)(1 - \nu). \qquad (14.3)\]

Table 14.1 (conditional probability that firm \(j\) is \ldots)

{\def\LTcaptype{none} % do not increment counter
\begin{longtable}[]{@{}lll@{}}
\toprule\noalign{}
when firm \(i\) is & productive & in distress \\
\midrule\noalign{}
\endhead
\bottomrule\noalign{}
\endlastfoot
productive (prob. \(x\)) & \(\mu\) & \(1 - \mu\) \\
in distress (prob. \(1 - x\)) & \(1 - \nu\) & \(\nu\) \\
\end{longtable}
}

\emph{Nonconcurrent risks.} In the first polar case, at most one firm is in distress. Put
differently, if firm \(i\) is in distress then firm \(j\) is not: \(\nu = 0\). The consistency
condition (14.3) then implies \(\mu = (2x - 1)/x\) (which naturally requires that
\(x \geqslant \frac12\)).

\emph{Identical shock.} The other polar case is that of perfectly correlated environments. There are
only two states of nature: either both firms are productive or both are in distress. This
corresponds to \(\mu = \nu = 1\).

\emph{Transfer price.} If both firms are in distress (which has probability \((1 - x)\nu\)), the
four participants (entrepreneurs, lenders) receive no ex post revenue. If neither is in distress
(which has probability \(x\mu\)), no sale occurs and the model is the standard variable-investment
one.

So let us consider the more interesting case in which firm 1, say, is in distress and firm 2 is not.
We then assume that lender 1 makes a take-it-or-leave-it offer to lender 2 (see the third remark
below for more general bargaining powers). Let \(P\) denote the per-unit price demanded by lender 1.

Note that lender 2 must adjust entrepreneur 2's incentive scheme to account for the increased
investment and therefore for the increased private benefit from not behaving (now equal to
\(B(I_1 + I_2)\) instead of \(BI_2\)). Assume that entrepreneur 2's incentive compatibility
constraint is binding in the absence of a purchase (\((\Delta p)R_{b2} = BI_2\)), which actually
turns out to be optimal. Then, lender 2 must raise entrepreneur 2's income in the case of success by
\(\delta R_{b2}\) such that\textsuperscript{12}

\[(\Delta p)(\delta R_{b2}) = BI_1.\]

So, entrepreneur 2's rent increases by \([p_H B/\Delta p]I_1\) and the transfer price is

\[PI_1 = p_H\left(R - \frac{B}{\Delta p}\right)I_1 = \rho_0 I_1.\]

The transfer price is simply the pledgeable income: \(P = \rho_0\). {[}\ldots{]} From (14.2), we see
that \(P < 1\).

So, the asset is sold at a discount even though the seller has the bargaining power. {[}\ldots{]}
Because firm 1 has the bargaining power, lender 2 actually gains nothing from firm 1's distress,
while entrepreneur 2 pockets the extra agency rent \(p_H BI_1/\Delta p = (\rho_1 - \rho_0)I_1\).

\emph{Entrepreneurs' expected utility.} Suppose entrepreneur \(i\) maximizes her net utility. As
usual, the lenders' zero-profit condition implies that the borrower receives the full surplus
associated with investment:

\[U_{bi} = [xp_H RI_i - I_i] + [(1 - x)(1 - \nu)PI_i] + \left[x(1 - \mu)\frac{p_H B}{\Delta p}I_j\right]. \qquad (14.4)\]

The first term in brackets in the expression of \(U_{bi}\) corresponds to the case in which distress
sales are impossible. The second term comes from lender \(i\)'s revenue from the sales of assets if
only firm \(i\) is in distress. The third term represents the expected windfall gain from firm
\(j\)'s distress.

We can rewrite (14.4) as

\[U_{bi} = [x\rho_i + (1 - x)(1 - \nu)\rho_0 - 1]I_i + [x(1 - \mu)(\rho_1 - \rho_0)]I_j. \qquad (14.5)\]

So let

\[\alpha \equiv x\rho_1 + (1 - x)(1 - \nu)\rho_0 - 1 \quad\text{and}\quad \kappa \equiv x(1 - \mu)(\rho_1 - \rho_0) > 0.\]

Using (14.4), we can rewrite borrower \(i\)'s net utility \(U_{bi}\) as

\[U_{bi} = \alpha I_i + \kappa I_j. \qquad (14.6)\]

\emph{Borrowing capacity.} Because lender \(i\) expects no surplus from the purchase of firm \(j\)'s
assets when the latter is in distress, lender \(i\)'s expected profit is

\[xp_H(RI_i - R_{bi}) + (1 - x)(1 - \nu)PI_i - (I_i - A_i) = 0,\]

where \(R_{bi}\) is borrower \(i\)'s income when firm \(i\) does not purchase firm \(j\)'s assets
and is successful. Incentive compatibility requires that \((\Delta p)R_{bi} \geqslant BI_i\), and,
as usual, this inequality is satisfied with equality in order to maximize pledgeable income and
therefore debt capacity. Using these two equations, firm \(i\)'s maximal investment is
\(I_i = kA_i\), where

\[k = \frac{1}{1 - \rho_0[x + (1 - x)(1 - \nu)]}. \qquad (14.7)\]

{[}\ldots{]} (a) \emph{Low correlation} (\(\nu < \nu^*\)). Then \(\alpha > 0\) and the firms invest
up to their borrowing capacity: \(I_i = kA_i\). While each firm's investment is independent of the
existence of the other firm, each firm derives a positive externality from this existence in the
form of asset redeployability. (b) \emph{High correlation} (\(\nu > \nu^*\)). Then \(\alpha < 0\),
and no firm invests, even though coordinated investment could be profitable (if
\(\alpha + \kappa > 0\)). {[}\ldots{]}

\emph{Remark (decreasing returns to scale).} Investment externalities are here positive. A firm's
investment allows it to stand by to purchase the other firm's assets if the latter is in distress.
On the other hand, if (due to entrepreneurial limited attention, for example) returns to investment
were decreasing rather than constant, investment externalities could become negative, as a more
active firm is less eager to take on new tasks. That is, the transfer price a firm in distress can
get for its assets is lower, the higher the level of existing investment by the other firm.

\emph{From Tirole footnotes 10 and 12, p.~500 (PDF 513):} 10. In particular, the entrepreneur in
firm \(j\) cannot enjoy private benefit \(BI_j\) by keeping the assets. 12. The new incentive
constraint is \((\Delta p)(R_{b2} + \delta R_{b2}) \geqslant B(I_2 + I_1)\).
\stmtend
%% END-STATEMENT

\asked Redo the two-firm resale model of \S14.2.2 with success income $R(I)$ concave, and find the sign of the effect of one firm's investment on the other firm's utility (the investment externality).

\approach Keep every step of \S14.2.2: the seller's lender makes a take-it-or-leave-it offer, so the price equals the buyer's increase in \emph{pledgeable income}; the buyer's entrepreneur pockets the extra agency rent. With concave $R$, the buyer's increase in pledgeable income from extra assets depends on how much she already runs: the resale price now falls with the buyer's own investment.

\notation
\begin{itemize}
  \item $I_i$, $A_i$: investment and wealth of firm $i\in\{1,2\}$; $R(\cdot)$: success income, $R'>0>R''$.
  \item $x$: probability of being productive; $\mu$, $\nu$: conditional probabilities of Table 14.1; $\pi_s\equiv x(1-\mu)=(1-x)(1-\nu)$ by (14.3): probability that a given firm is productive and the other in distress.
  \item $b\equiv p_HB/\Delta p$: agency rent per unit of assets; $\Pc(I)\equiv p_HR(I)-bI$: pledgeable income of a firm running $I$.
\end{itemize}

\step{1}{Transfer price} Let $j$ be in distress and $i$ productive. With $(\Delta p)R_{bi}=BI_i$ binding, running $I_i+I_j$ requires $(\Delta p)(R_{bi}+\delta R_{bi})\geq B(I_i+I_j)$, i.e.\ $\delta R_{bi}=BI_j/\Delta p$; the rent rises by $bI_j$. Lender $i$'s willingness to pay is the increase in income net of that rent:
\begin{align}
  P_jI_j = p_H[R(I_i+I_j)-R(I_i)]-bI_j. \label{r10e6-price}
\end{align}
With $R(I)=RI$ this is $\rho_0I_j$ (the book's $P=\rho_0$). Now $\partial(P_jI_j)/\partial I_i=p_H[R'(I_i+I_j)-R'(I_i)]<0$.

\step{2}{Entrepreneur $j$'s utility} Lenders break even, so $U_{bj}$ is the NPV of firm $j$: output when both productive ($x\mu$), output plus the acquisition when $i$ is in distress ($\pi_s$), sale revenue when $j$ is in distress and $i$ productive ($\pi_s$), minus $I_j$. When $j$ buys, it pays $P_iI_i=p_H[R(I_j+I_i)-R(I_j)]-bI_i$, so its net gain is $p_HR(I_j+I_i)-P_iI_i=p_HR(I_j)+bI_i$:
\begin{align}
  U_{bj} &= x\mu\,p_HR(I_j)+\pi_s\bigl[p_HR(I_j)+bI_i\bigr]+\pi_s\bigl\{p_H[R(I_i+I_j)-R(I_i)]-bI_j\bigr\}-I_j \nonumber\\
         &= xp_HR(I_j)-I_j+\pi_s\,bI_i+\pi_s\bigl\{p_H[R(I_i+I_j)-R(I_i)]-bI_j\bigr\}, \label{r10e6-Ub}
\end{align}
using $x\mu+\pi_s=x$. Check: with $R(I)=RI$, \eqref{r10e6-Ub} gives $[x\rho_1+(1-x)(1-\nu)\rho_0-1]I_j+x(1-\mu)(\rho_1-\rho_0)I_i=\alpha I_j+\kappa I_i$, which is (14.5).

\step{3}{The externality} Differentiate \eqref{r10e6-Ub} with respect to $I_i$:
\begin{align}
  \frac{\partial U_{bj}}{\partial I_i} &= \pi_s b+\pi_s p_H[R'(I_i+I_j)-R'(I_i)] \nonumber\\
  &= \pi_s\,p_H\left[\frac{B}{\Delta p}-\bigl(R'(I_i)-R'(I_i+I_j)\bigr)\right]. \label{r10e6-ext}
\end{align}
Two effects: the \emph{windfall} $\pi_sb>0$ (as in the CRS model, $\kappa$) and the \emph{price effect} $\pi_sp_H[R'(I_i+I_j)-R'(I_i)]<0$ (new): a bigger potential buyer pays less for $j$'s assets.

\step{4}{Borrowing capacity} Lender $j$'s breakeven, with $j$ selling at \eqref{r10e6-price} when in distress:
\begin{align}
  x\,\Pc(I_j)+\pi_s\bigl\{p_H[R(I_i+I_j)-R(I_i)]-bI_j\bigr\} = I_j-A_j. \label{r10e6-bc}
\end{align}
The left side falls with $I_i$ (derivative $\pi_sp_H[R'(I_i+I_j)-R'(I_i)]<0$): firm $i$'s investment lowers $j$'s collateral value and borrowing capacity. In the CRS model it did not (multiplier (14.7) independent of $I_i$).

\step{5}{Example} $R(I)=2\sqrt I$, nonconcurrent risks $\nu=0$, $x=0.8$ (so $\mu=0.75$, $\pi_s=0.2$), $p_H=0.8$, $B/\Delta p=0.5$. At $I_i=I_j=1$: $R'(1)-R'(2)=1-0.707107=0.292893$ and \eqref{r10e6-ext} $=0.16(0.5-0.292893)=0.033137>0$; the price revenue is $0.8(2\sqrt2-2)-0.4=0.262742>0$. At $I_i=I_j=0.25$: $R'(0.25)-R'(0.5)=2-1.414214=0.585786$ and \eqref{r10e6-ext} $=0.16(0.5-0.585786)=-0.013726<0$.

\answer{} $\partial U_{bj}/\partial I_i=\pi_sp_H\{B/\Delta p-[R'(I_i)-R'(I_i+I_j)]\}$: ambiguous. Positive when returns decrease slowly (the CRS limit gives $\kappa>0$), negative when concavity is strong ($R'(I_i)-R'(I_i+I_j)>B/\Delta p$). Through borrowing capacity, $I_i$ always hurts firm $j$.

\intuition A firm's investment makes it a ready buyer for its rival's assets, which helps the rival only through the agency rent the buyer earns on acquired assets. With decreasing returns, a firm that already runs a lot gains little from extra assets, so the rival in distress fetches a lower price. The more concave the technology, the more this ``less eager buyer'' effect dominates and investment becomes a negative externality.

\pitfalls
\begin{itemize}
  \item Keeping $P=\rho_0$: with concave $R$ the price depends on both investment levels.
  \item Forgetting that the buyer's private benefit scales with all assets run, $B(I_i+I_j)$, which is what generates the rent $bI_j$.
  \item Mixing the two probabilities: by (14.3) the windfall and the sale states have the same probability $\pi_s$.
\end{itemize}

\bookans The book gives no answer to this exercise. The result agrees with the text's remark on decreasing returns (p.~503): ``investment externalities could become negative \ldots{} the transfer price a firm in distress can get for its assets is lower, the higher the level of existing investment by the other firm.''
\refbib{Tirole \S14.2.2, pp.~500--503; Exercise 3.5 (decreasing returns).}

% =============================================================================
\exsec{7. Inefficiently low volume of asset reallocations --- Tirole 14.4, p.~516}
%% BEGIN-STATEMENT 7
\stmtbegin
\textbf{Source:} Tirole, Ch. 14, Exercise 14.4, p.~516 (PDF 529) · Answer in book: yes, p.~619

\textbf{Exercise 14.4 (inefficiently low volume of asset reallocations).} This exercise applies the
logic of corporate risk management developed in Chapter 5 to show that, even with frictionless
resale markets, there will be an inefficiently low volume of transactions in the secondary market.

There are three dates, \(t = 0, 1, 2\), and at least two firms \(i = 1, 2\).

Firm 1, the firm of interest, is managed by a risk-neutral entrepreneur, who owns initial wealth
\(A\) at date 0 and is protected by limited liability. This firm invests at a variable investment
level \(I \in [0, \infty)\). The per-unit profitability of investment is random and learned at date
1. The investment yields \(RI\) with probability \(p + \tau\) and 0 with probability
\(1 - (p + \tau)\). The random variable \(\tau\) is drawn from a continuous distribution. The
variable \(p\) is equal to \(p_H\) if the entrepreneur behaves (no private benefit) and \(p_L\) if
the entrepreneur misbehaves (private benefit \(BI\)). Let

\[\rho_1 = (p_H + \tau)R\]

and

\[\rho_0 = (p_H + \tau)\left(R - \frac{B}{\Delta p}\right) \equiv \rho_1 - \Delta\rho\]

denote the random continuation per-unit NPV and pledgeable income when the entrepreneur behaves and
the realization of profitability is \(\tau\). The distribution on \(\tau\) induces a cumulative
distribution function \(F(\rho_0)\) on \([\underline\rho_0, \bar\rho_0]\).

At date 1, the firm may either continue or resell assets \(I\) to firm 2 (or to a competitive
market). Firm 2 has a known level \(\hat\rho_0\) of per-unit pledgeable income per unit of
investment (its NPV per unit of investment is in general larger than this).

Firms 1 and 2 do not contract with each other at date 0. Rather, investors in firm 1 make a
take-it-or-leave-it offer to firm 2 at date 1 if firm 1's initial contract specifies that assets
ought to be reallocated.

Assume for simplicity that the contract between firm 1's investors and the entrepreneur can be
contingent on the realization of \(\rho_0\).

Show that at the optimal contract assets are resold whenever \(\rho_0 < \rho_0^*\), where

\[\rho_0^* < \hat\rho_0,\]

and so the volume of asset reallocations is inefficiently low.

\paragraph{Context needed to solve it}

\emph{From Tirole §14.2.2, pp.~501--502 (PDF 514--515)} --- the take-it-or-leave-it price paid by an
acquirer is its pledgeable income per unit:

{[}\ldots{]} lender 1 makes a take-it-or-leave-it offer to lender 2 {[}\ldots{]}. Note that lender 2
must adjust entrepreneur 2's incentive scheme to account for the increased investment and therefore
for the increased private benefit from not behaving {[}\ldots{]}. So, entrepreneur 2's rent
increases by \([p_H B/\Delta p]I_1\) and the transfer price is
\(PI_1 = p_H(R - B/\Delta p)I_1 = \rho_0 I_1\). The transfer price is simply the pledgeable income:
\(P = \rho_0\). {[}\ldots{]} So, the asset is sold at a discount even though the seller has the
bargaining power. In this sense, the asset market exhibits some degree of \emph{illiquidity}.
Indeed, while reallocation of assets is here efficient as assets in the firm in distress have a zero
productivity, it could be inefficient in an environment in which productivity does not fall to 0,
but is still lower than the productivity following a reallocation to the other firm (see Exercise
14.4).
\stmtend
%% END-STATEMENT

\asked Show that firm 1's optimal contract resells its assets to firm 2 exactly when the realized pledgeable income $\rho_0$ falls below a cutoff $\rho_0^*$ that is strictly below firm 2's pledgeable income $\hat\rho_0$, and explain why that is too few resales.

\approach This is the continuation-vs-liquidation trade-off of Chapter 5 (risk management) with resale as the alternative to continuation. Firm 1's investors get $\hat\rho_0$ per unit if they sell (take-it-or-leave-it price = firm 2's pledgeable income, \S14.2.2) and $\rho_0$ if they continue; the entrepreneur gets the extra rent $\Delta\rho$ only if the firm continues. The cutoff maximizes the NPV subject to breakeven.

\notation
\begin{itemize}
  \item $\rho_0\sim F$ on $[\underline\rho_0,\bar\rho_0]$, density $f$: firm 1's pledgeable income per unit; $\rho_1=\rho_0+\Delta\rho$: its NPV per unit.
  \item $\hat\rho_0$: firm 2's pledgeable income per unit (= resale price); $\hat\rho_1\geq\hat\rho_0$: firm 2's NPV per unit.
  \item $c$: resale cutoff (resell iff $\rho_0<c$); $\pi(c)$, $v(c)$: pledgeable income and value per unit of investment.
\end{itemize}

\step{1}{Contract} Assets are resold iff $\rho_0<c$ (a cutoff is optimal since continuation becomes more valuable as $\rho_0$ rises). Resale proceeds $\hat\rho_0I$ go to the investors: after a resale there is no moral hazard in firm 1, so paying the entrepreneur would only lower pledgeable income. Per unit of $I$:
\begin{align}
  \pi(c) &= F(c)\hat\rho_0+\int_c^{\bar\rho_0}\rho_0\,dF(\rho_0), \label{r10e7-pi}\\
  v(c) &= F(c)\hat\rho_0+\int_c^{\bar\rho_0}(\rho_0+\Delta\rho)\,dF(\rho_0) = \pi(c)+\Delta\rho[1-F(c)]. \label{r10e7-v}
\end{align}

\step{2}{Investment and utility} Breakeven: $\pi(c)I=I-A$, so
\begin{align}
  I(c) = \frac{A}{1-\pi(c)}. \label{r10e7-I}
\end{align}
The entrepreneur receives the NPV:
\begin{align}
  U_b(c) = [v(c)-1]\,I(c) = \frac{v(c)-1}{1-\pi(c)}A. \label{r10e7-Ub}
\end{align}

\step{3}{First-order condition} Differentiating \eqref{r10e7-pi}--\eqref{r10e7-v} (Leibniz rule): $\pi'(c)=f(c)(\hat\rho_0-c)$ and $v'(c)=f(c)(\hat\rho_0-c-\Delta\rho)$. Setting $dU_b/dc=0$ in \eqref{r10e7-Ub}:
\begin{align}
  v'(c)[1-\pi(c)]+[v(c)-1]\pi'(c) &= 0 \nonumber\\
  (\hat\rho_0-c-\Delta\rho)[1-\pi]+[v-1](\hat\rho_0-c) &= 0 \nonumber\\
  (\hat\rho_0-c)[1-\pi+v-1] &= \Delta\rho[1-\pi] \nonumber\\
  (\hat\rho_0-c)\,\Delta\rho[1-F(c)] &= \Delta\rho[1-\pi] \nonumber\\
  \rho_0^* &= \hat\rho_0-\frac{1-\pi(\rho_0^*)}{1-F(\rho_0^*)}, \label{r10e7-foc}
\end{align}
using $v-\pi=\Delta\rho(1-F)$ from \eqref{r10e7-v}.

\step{4}{Bounds} Since $\pi<1$ (finite investment), the fraction in \eqref{r10e7-foc} is positive: $\rho_0^*<\hat\rho_0$. And $(1-\pi)/(1-F)<\Delta\rho\iff1-\pi<\Delta\rho(1-F)=v-\pi\iff v>1$, which holds because the project has positive NPV. Hence
\begin{align}
  \hat\rho_0-\Delta\rho<\rho_0^*<\hat\rho_0. \label{r10e7-bounds}
\end{align}

\step{5}{Why ``inefficiently low''} Moving a unit to firm 2 creates $\hat\rho_1-\rho_1=\hat\rho_1-\rho_0-\Delta\rho$. If firm 2's agency rent is at least firm 1's, $\hat\rho_1-\hat\rho_0\geq\Delta\rho$ (e.g.\ equal per-unit agency costs), reallocation is efficient whenever $\rho_0<\hat\rho_0$; firm 1 resells only below $\rho_0^*<\hat\rho_0$. It ignores the rent $\hat\rho_1-\hat\rho_0$ it leaves to firm 2's insiders, and it keeps assets in some states to protect its own entrepreneur's rent.

\step{6}{Example} $\rho_0\sim U[0,1]$, $\hat\rho_0=0.7$, $\Delta\rho=0.5$: $\pi(c)=0.7c+(1-c^2)/2$, so \eqref{r10e7-foc} becomes $(0.7-c)(1-c)=0.5-0.7c+0.5c^2$, i.e.\ $c^2-2c+0.4=0$, and $\rho_0^*=1-\sqrt{0.6}=0.225403\in(0.2,0.7)$; $v=1.019677$, $I/A=2.720193$, $U_b/A=0.053526$.

\answer{} Resell iff $\rho_0<\rho_0^*$ with $\rho_0^*=\hat\rho_0-[1-\pi(\rho_0^*)]/[1-F(\rho_0^*)]$, so $\hat\rho_0-\Delta\rho<\rho_0^*<\hat\rho_0$: assets with $\rho_0^*<\rho_0<\hat\rho_0$ stay with firm 1 though they are worth more in firm 2.

\intuition Reselling raises pledgeable income (investors get $\hat\rho_0>\rho_0$) but destroys the entrepreneur's continuation rent; as in Chapter 5's liquidity choice, the cutoff sits between the pure pledgeable-income rule ($\hat\rho_0$) and the pure NPV rule ($\hat\rho_0-\Delta\rho$). Because the buyer pays only its pledgeable income, part of the gain from trade is invisible to the seller, so even a frictionless resale market sees too little reallocation.

\pitfalls
\begin{itemize}
  \item Taking the resale price to be firm 2's NPV: under take-it-or-leave-it with a financially constrained buyer it is $\hat\rho_0$.
  \item Giving resale proceeds to the entrepreneur: it lowers pledgeable income for no incentive gain.
  \item Maximising gross value $v(c)I(c)$ instead of the NPV $[v(c)-1]I(c)$ (see the book answer).
\end{itemize}

\bookans Matches the book's answer (p.~619) on the structure (resale iff $\rho_0<\rho_0^*$, proceeds to investors, breakeven $[F(\rho_0^*)\hat\rho_0+\int_{\rho_0^*}\rho_0\,dF]I=I-A$) and on the bounds $\hat\rho_0-\Delta\rho<\rho_0^*<\hat\rho_0$. The book writes $U_b=[F\hat\rho_0+\int(\rho_0+\Delta\rho)dF]\,I(\rho_0^*)$, which omits $-I$; maximising that gross expression gives a different cutoff (0.591681 instead of 0.225403 in the example), still inside the bounds. The correct objective is the NPV \eqref{r10e7-Ub}.
\refbib{Tirole \S14.2.2--14.2.3, pp.~500--505; Ch.~5 (liquidity and risk management).}

% =============================================================================
\exsec{8. Inside liquidity --- Tirole Review Problem 6, pp.~629--630}
%% BEGIN-STATEMENT 8
\stmtbegin
\textbf{Source:} Tirole, Part VII, Review Problem 6, pp.~629--630 (PDF 642) · Answer in book: yes,
p.~635

\textbf{Review Problem 6 (inside liquidity).} Consider the variable-investment model with two
possible values of liquidity shocks (0 and \(\rho\) per unit of investment). The timing is described
in Figure 2.

Figure 2 (timeline, left to right; p.~629, PDF 642)

{\def\LTcaptype{none} % do not increment counter
\begin{longtable}[]{@{}
  >{\raggedright\arraybackslash}p{(\linewidth - 6\tabcolsep) * \real{0.2500}}
  >{\raggedright\arraybackslash}p{(\linewidth - 6\tabcolsep) * \real{0.2500}}
  >{\raggedright\arraybackslash}p{(\linewidth - 6\tabcolsep) * \real{0.2500}}
  >{\raggedright\arraybackslash}p{(\linewidth - 6\tabcolsep) * \real{0.2500}}@{}}
\toprule\noalign{}
\begin{minipage}[b]{\linewidth}\raggedright
Date 0
\end{minipage} & \begin{minipage}[b]{\linewidth}\raggedright
Date 1
\end{minipage} & \begin{minipage}[b]{\linewidth}\raggedright
Date 1 → 2
\end{minipage} & \begin{minipage}[b]{\linewidth}\raggedright
Date 2
\end{minipage} \\
\midrule\noalign{}
\endhead
\bottomrule\noalign{}
\endlastfoot
Entrepreneur has wealth \(A\), invests \(I\), borrows \(I - A\). & Firm is ``intact'' (no
reinvestment needed) with probability \(1 - \lambda\), and ``distressed'' (reinvestment \(\rho\) per
salvaged unit) with probability \(\lambda\). & Moral hazard. & Success (\(R\) per unit) with
probability \(p\) and failure (0) with probability \(1 - p\), for the salvaged investment. \\
\end{longtable}
}

Investors and entrepreneurs are risk neutral, the entrepreneur is protected by limited liability and
the rate of interest in the economy is 0. If the firm is in distress (suffers a liquidity shock), a
reinvestment \(\rho xI\) allows it to salvage a fraction \(x \in [0, 1]\) of the investment (so
there is no constraint to salvage all or nothing, even though, as we will see, the solution will be
a ``corner solution'').

Continuation is subject to moral hazard. The probability of success is \(p_H\) if the entrepreneur
behaves and \(p_L\) is she misbehaves. The private benefit of misbehaving is \(BxI\). The project
yields \(RxI\) in the case of success and 0 in the case of failure. Let

\[\rho_0 \equiv p_H\left(R - \frac{B}{\Delta p}\right) < c \equiv \min\left\{1 + \lambda\rho, \frac{1}{1 - \lambda}\right\} < \rho_1 \equiv p_H R.\]

In the first two questions, one will assume that there is a costless outside store of value (there
are assets that at per-unit cost \(q = 1\) at date 0 yield a return equal to 1 at date 1).

\begin{enumerate}
\def\labelenumi{(\roman{enumi})}
\tightlist
\item
  Show that, when choosing \(x\), the entrepreneur can borrow up to
\end{enumerate}

\[I = \frac{A}{(1 + \lambda\rho x) - [1 - \lambda + \lambda x]\rho_0}.\]

\begin{enumerate}
\def\labelenumi{(\roman{enumi})}
\setcounter{enumi}{1}
\tightlist
\item
  Compute the borrower's utility and show that
\end{enumerate}

\[x = 1 \quad\text{if and only if}\quad (1 - \lambda)\rho \leqslant 1\]

(and \(x = 0\) otherwise).

(Hint: write the borrower's utility as a function of the ``average unit cost of preserved
investment.'')

\begin{enumerate}
\def\labelenumi{(\roman{enumi})}
\setcounter{enumi}{2}
\item
  Suppose now that there is no outside store of value. There is mass 1 of (\emph{ex ante} identical)
  entrepreneurs. The only liquidity in the economy is the inside liquidity created by the securities
  issued by the firms. One will assume that \(\rho_0 < \rho < 1/[1 - \lambda]\). Is there enough
  liquidity if the firms' liquidity shocks are perfectly correlated? If not, what is the level of
  the liquidity shortage?
\item
  Consider question (iii) except that the liquidity shocks faced by the entrepreneurs are
  independently distributed. Show that the firms' holding the stock index may not be optimal. What
  should be done?
\end{enumerate}

\paragraph{Context needed to solve it}

\emph{From Tirole §15.2.2 ``Wasting Inside Liquidity'', pp.~521--522 (PDF 534--535)} --- what
``holding the stock index'' means and the remedy:

{[}\ldots{]} the ``natural implementation,'' namely, the policy that consists in each firm holding
the stock index, that is, a representative portfolio of claims on all firms in the economy, in
general does not work. {[}\ldots{]} Why does this ``self-provision'' of liquidity result in a waste
of liquidity? Firms that face a low liquidity shock at date 1 have excess cash {[}\ldots{]}. At date
1, they either redistribute the excess cash to or invest it on behalf of their owners; they have no
incentive to invest in distressed firms, in which the reinvestment cost exceeds the pledgeable
income. {[}\ldots{]} in order to force healthy firms to redispatch cash to distressed ones at date
1, firms must at date 0 pool their liquidity and organize a system of (or akin to) credit lines. For
example, shares are deposited with one (or an arbitrary number of) financial institutions; each firm
is then entitled to draw on a credit line {[}\ldots{]}. The financial institution can raise the cash
needed to honor the credit lines by selling at date 1 shares it holds in firms to date-1 investors.
\stmtend
%% END-STATEMENT

\asked (i) Borrowing capacity for a given salvaged fraction $x$. (ii) The borrower's utility and the optimal $x$. (iii) Whether inside liquidity suffices under perfectly correlated shocks, and the size of the shortage. (iv) Why holding the stock index may fail under independent shocks, and the remedy.

\approach Chapter 5's liquidity model: the cost of continuing distressed firms is weighed against scale through the \emph{unit cost of preserved investment} $c(x)$. Chapter 15 then asks whether claims issued by the firms themselves can fund the reinvestment: never under a common shock; in principle yes under independent shocks, but only if liquidity is pooled.

\notation
\begin{itemize}
  \item $\lambda$: probability of distress; $\rho$: reinvestment per salvaged unit; $x\in[0,1]$: fraction salvaged in distress.
  \item $\rho_0=p_H(R-B/\Delta p)$, $\rho_1=p_HR$: pledgeable income and value per salvaged unit; $c(x)$: average cost per unit of preserved investment.
\end{itemize}

\step{1}{Borrowing capacity} Investors pay $I-A$ at date 0 and $\lambda\rho xI$ expected at date 1 (via the free store of value); they receive $\rho_0$ per salvaged unit, of which there are $(1-\lambda)I+\lambda xI$:
\begin{align}
  (1+\lambda\rho x)I-A &= [1-\lambda+\lambda x]\rho_0I \nonumber\\
  I &= \frac{A}{(1+\lambda\rho x)-[1-\lambda+\lambda x]\rho_0}. \label{r10e8-I}
\end{align}

\answer{(i)} $I=A/\{(1+\lambda\rho x)-[1-\lambda+\lambda x]\rho_0\}$.

\step{2}{Utility} $U_b$ = value of preserved units minus total cost:
\begin{align}
  U_b &= [(1-\lambda+\lambda x)\rho_1-(1+\lambda\rho x)]I. \label{r10e8-Ub0}
\end{align}
Define $c(x)=(1+\lambda\rho x)/(1-\lambda+\lambda x)$, so $1+\lambda\rho x=(1-\lambda+\lambda x)c(x)$. Then \eqref{r10e8-I} reads $I=A/[(1-\lambda+\lambda x)(c(x)-\rho_0)]$ and
\begin{align}
  U_b = (1-\lambda+\lambda x)[\rho_1-c(x)]\cdot\frac{A}{(1-\lambda+\lambda x)[c(x)-\rho_0]} = \frac{\rho_1-c(x)}{c(x)-\rho_0}A. \label{r10e8-Ub}
\end{align}

\step{3}{Optimal $x$} $\frac{d}{dc}\frac{\rho_1-c}{c-\rho_0}=-\frac{\rho_1-\rho_0}{(c-\rho_0)^2}<0$, so minimise $c(x)$:
\begin{align}
  c'(x) &= \frac{\lambda\rho(1-\lambda+\lambda x)-\lambda(1+\lambda\rho x)}{(1-\lambda+\lambda x)^2} = \frac{\lambda[(1-\lambda)\rho-1]}{(1-\lambda+\lambda x)^2}. \label{r10e8-cprime}
\end{align}
The sign does not depend on $x$: $c$ is decreasing iff $(1-\lambda)\rho<1$. Hence $x=1$ iff $(1-\lambda)\rho\leq1$ (indifference at equality), $x=0$ otherwise. The assumption $\rho_0<c=\min\{1+\lambda\rho,1/(1-\lambda)\}<\rho_1$ says that at the chosen corner, $c(1)=1+\lambda\rho$ or $c(0)=1/(1-\lambda)$, investment is finite ($c>\rho_0$) and profitable ($c<\rho_1$).

\answer{(ii)} $U_b=\dfrac{\rho_1-c(x)}{c(x)-\rho_0}A$ with $c(x)=\dfrac{1+\lambda\rho x}{1-\lambda+\lambda x}$; $x=1$ iff $(1-\lambda)\rho\leq1$.

\step{4}{Correlated shocks} With $\rho<1/(1-\lambda)$, $x=1$ is desired. Under a common shock all firms are distressed together; each can raise $\rho_0I$ by issuing new claims on its own date-2 income, and claims on other firms are worth nothing more than what those firms are themselves issuing. The need is $\rho I$, so
\begin{align}
  \text{shortage} = (\rho-\rho_0)I > 0, \qquad I=\frac{A}{1+\lambda\rho-\rho_0}. \label{r10e8-short}
\end{align}

\answer{(iii)} No: the shortage is $(\rho-\rho_0)I$.

\step{5}{Independent shocks and the index} A fraction $\lambda$ is distressed. After distressed firms dilute their old shareholders completely (they raise $\rho_0I<\rho I$), only intact firms have value for old holders, so a share of the index is worth $(1-\lambda)\rho_0I$ at date 1. A distressed firm that holds the index can raise $\rho_0I+(1-\lambda)\rho_0I$, and needs $\rho I$:
\begin{align}
  \text{index suffices} \iff \rho-\rho_0\leq(1-\lambda)\rho_0. \label{r10e8-index}
\end{align}
In aggregate, though, liquidity is ample: the net-borrower breakeven $(1+\lambda\rho)I-A=\rho_0I$ gives $I-A=[(1-\lambda)\rho_0-\lambda(\rho-\rho_0)]I>0$, i.e.\ intact firms' value exceeds distressed firms' net need. The index fails because healthy firms keep their excess liquidity.

\step{6}{Example} $\lambda=0.4$, $\rho=1.2$, $\rho_0=0.5$, $\rho_1=2$, $A=1$: $(1-\lambda)\rho=0.72\leq1$ so $x=1$; $c(1)=1.48$; $I=1/0.98=1.020408$; $U_b=0.52/0.98=0.530612$. Shortage under correlation: $0.7(1.020408)=0.714286$. Under independence: $\rho-\rho_0=0.7>(1-\lambda)\rho_0=0.3$, so the index is not enough, although $(1-\lambda)\rho_0=0.3>\lambda(\rho-\rho_0)=0.28$ in aggregate.

\answer{(iv)} If $\rho-\rho_0>(1-\lambda)\rho_0$ the index leaves distressed firms short even though aggregate inside liquidity suffices; firms should pool liquidity (credit lines from an intermediary that holds the shares and resells them at date 1).

\intuition Liquidity is about moving resources across firms and states. A common shock leaves no one to move them from. Independent shocks do, but a passive portfolio spreads the healthy firms' value thinly across everyone, including firms that do not need it; a credit line concentrates it where the shock hits.

\pitfalls
\begin{itemize}
  \item Optimising $x$ inside $(0,1)$: $c(x)$ is monotone, so the solution is a corner.
  \item Counting distressed firms' old shares in the index value at date 1: they are diluted to zero.
  \item Concluding from (iv) that there is an aggregate shortage: the problem is allocation, not quantity.
\end{itemize}

\bookans Matches the book's answer (p.~635): the breakeven and $I$, $U_b=(\rho_1-c)/(c-\rho_0)A$, $x=1\iff(1-\lambda)\rho\leq1$, shortage $(\rho-\rho_0)I$, index worth $(1-\lambda)\rho_0I$ and insufficient if $\rho-\rho_0>(1-\lambda)\rho_0$, remedy a liquidity pool.
\refbib{Tirole Ch.~5 (liquidity needs); \S15.2--15.3, pp.~518--525.}

% =============================================================================
\exsec{9. Time consistency and the soft budget constraint --- Tirole 16.4, p.~566}
%% BEGIN-STATEMENT 9
\stmtbegin
\textbf{Source:} Tirole, Ch. 16, Exercise 16.4, p.~566 (PDF 579) · Answer in book: no

\textbf{Exercise 16.4 (time consistency and the soft budget constraint).} A firm is run by a
risk-neutral entrepreneur with wealth \(A\), and has a fixed-size project with investment cost
\(I\). The project, if undertaken at date 0, will deliver a verifiable income,
\(y \in \{y_L, y_H\}\) in the case of success and 0 in the case of failure, at date 2, provided that
one worker is employed in the firm. The project yields nothing if it is interrupted (the worker is
laid off). \(y = y_H\) with probability \(\rho\) and \(y = y_L < y_H\) with probability
\(1 - \rho\).

Moreover, in the case of ``continuation'' and regardless of the value of \(y\), the entrepreneur may
behave (the income is \(y\) for certain, the entrepreneur receives no private benefit) or misbehave
(the income is \(y\) with probability \(p_L\) and 0 with probability \(1 - p_L\), the entrepreneur
receives private benefit \(B\)). The entrepreneur is protected by limited liability. Let

\[\mathcal{R} \equiv \frac{B}{1 - p_L}.\]

(One will assume that \(B\) is small enough that it is worth inducing the entrepreneur to behave in
the case of continuation.) The (risk-neutral) worker is paid \(w\) in the case of continuation and 0
otherwise. He obtains unemployment benefit paid by the state \(w_u < w\) when laid off. We take
\(w\) and \(w_u\) as given. (Note: they could be endogenized through some efficiency wage and
incentive-to-search stories, but take these as exogenous for this exercise.)

Figure 16.12 (timeline, left to right; p.~566, PDF 579)

{\def\LTcaptype{none} % do not increment counter
\begin{longtable}[]{@{}
  >{\raggedright\arraybackslash}p{(\linewidth - 6\tabcolsep) * \real{0.2500}}
  >{\raggedright\arraybackslash}p{(\linewidth - 6\tabcolsep) * \real{0.2500}}
  >{\raggedright\arraybackslash}p{(\linewidth - 6\tabcolsep) * \real{0.2500}}
  >{\raggedright\arraybackslash}p{(\linewidth - 6\tabcolsep) * \real{0.2500}}@{}}
\toprule\noalign{}
\begin{minipage}[b]{\linewidth}\raggedright
Date 0
\end{minipage} & \begin{minipage}[b]{\linewidth}\raggedright
Date 1
\end{minipage} & \begin{minipage}[b]{\linewidth}\raggedright
Date 1 → 2
\end{minipage} & \begin{minipage}[b]{\linewidth}\raggedright
Date 2
\end{minipage} \\
\midrule\noalign{}
\endhead
\bottomrule\noalign{}
\endlastfoot
Entrepreneur invests \(I\), borrows \(I - A\). Entrepreneur and investors choose liquidity \(L\).
Entrepreneur hires one worker. & Productivity \(y \in \{y_L, y_H\}\) is revealed. Firm continues
(keeps worker, pays him \(w\)), or stops (lays worker off). & Moral hazard (in the case of
continuation): the entrepreneur behaves or misbehaves. & Realization of income (\(y\) or 0). \\
\end{longtable}
}

Assume that the interest rate in the economy is 0 and that

\[w < y_L < w + \mathcal{R}\]

and

\[I - A \leqslant \rho(y_H - w - \mathcal{R}) + (1 - \rho)(y_L - w - \mathcal{R}).\]

\begin{enumerate}
\def\labelenumi{(\roman{enumi})}
\item
  Write the firm's NPV depending on whether the firm continues (\(x = 1\)) or stops (\(x = 0\)) when
  productivity is low (\(y = y_L\)). Show that \(x^* = 1\). Assuming a perfectly functioning capital
  market at date 1, what is the amount of liquidity that is needed to complement capital market
  refinancing?
\item
  Introduce a government that can at date 1 bring a subsidy \(s \geqslant 0\) to the firm (it is a
  pure subsidy: the government takes no ownership stake in exchange). The shadow cost of public
  funds is \(\lambda\), and so the cost of subsidy \(s\) for the taxpayers is \((1 + \lambda)s\).
  The government maximizes total welfare (entrepreneur, investors, worker, taxpayers). Assuming that
\end{enumerate}

\[\lambda[(w - w_u) + \mathcal{R}] \leqslant (1 + \lambda)y_L\]

and that the government selects its subsidy at date 1 (having observed the realization of \(y\)),
what is the liquidity \(L\) chosen by entrepreneur and investors at date 0?

How would the government (contingent) choice of \(s\) be affected if the government could commit to
\(s\) at date 0, before the investors and the entrepreneur write their contract?

\paragraph{Context needed to solve it}

None beyond Figure 16.12, transcribed above.
\stmtend
%% END-STATEMENT

\asked (i) The firm's NPV with and without continuation in the low state, why continuing is optimal, and how much liquidity the firm must hoard on top of date-1 refinancing. (ii) The liquidity chosen at date 0 when a welfare-maximising government decides on a subsidy at date 1, and how a government that commits at date 0 would behave.

\approach The project is the fixed-size version of Chapter 5's liquidity model: continuation costs the wage $w$, of which the capital market will fund only the pledgeable part $y-\Rc$. The government's date-1 problem is a bailout decision with sunk liquidity choices; anticipating a bailout, private parties under-hoard --- the \emph{soft budget constraint} (Tirole \S5.5, from p.~220; time inconsistency, \S16.3.2).

\notation
\begin{itemize}
  \item $y\in\{y_L,y_H\}$, $\Pr(y_H)=\rho$; $w$: wage if continuation; $w_u<w$: unemployment benefit; $x\in\{0,1\}$: continuation in the low state.
  \item $\Rc=B/(1-p_L)$: minimum stake that makes the entrepreneur behave; $L$: liquidity hoarded at date 0; $s$: subsidy; $\lambda$: shadow cost of public funds.
\end{itemize}

\step{1}{Incentives and pledgeable income at date 1} Behaving gives $y$ for sure; misbehaving gives $y$ w.p.\ $p_L$ plus $B$. With stake $R_b$ in $y$: $R_b\geq p_LR_b+B\iff R_b\geq B/(1-p_L)=\Rc$. Continuing at productivity $y$ yields income $y$, of which $y-\Rc$ is pledgeable, and costs $w$.

\step{2}{NPV and $x^*$} In the high state the firm always continues ($y_H>y_L>w$):
\begin{align}
  \text{NPV}(x) = \rho(y_H-w)+(1-\rho)\,x\,(y_L-w)-I. \label{r10e9-npv}
\end{align}
$\text{NPV}(1)-\text{NPV}(0)=(1-\rho)(y_L-w)>0$. Financing with $x=1$ requires pledgeable income $\geq I-A$:
\begin{align}
  \rho(y_H-w-\Rc)+(1-\rho)(y_L-w-\Rc)\geq I-A, \label{r10e9-fin}
\end{align}
which is the second assumption. So $x^*=1$.

\step{3}{Liquidity} At date 1 the firm can raise $y-\Rc$ from new investors. In the low state, $y_L-\Rc<w$ (first assumption), so the shortfall
\begin{align}
  L^* = w+\Rc-y_L>0 \label{r10e9-L}
\end{align}
must be secured at date 0 (hoarded or a credit line). In the high state no liquidity is needed: with $I>A$, \eqref{r10e9-fin} and $y_L-w-\Rc<0$ imply $y_H-w-\Rc>0$. With a costless store of value, unused liquidity is returned to investors, so \eqref{r10e9-fin} remains the breakeven.

\answer{(i)} $\text{NPV}(1)=\rho(y_H-w)+(1-\rho)(y_L-w)-I>\text{NPV}(0)=\rho(y_H-w)-I$, financeable, so $x^*=1$; liquidity $L^*=w+\Rc-y_L$ for the low state.

\step{4}{Date-1 bailout} In the low state with hoarded $L<L^*$, continuation needs $s=w+\Rc-y_L-L$. Welfare (entrepreneur + investors + worker + taxpayers), counting the wage and the subsidy as transfers:
\begin{align}
  W_{\text{continue}} &= (y_L-w+s)+w-(1+\lambda)s = y_L-\lambda s, \label{r10e9-Wc}\\
  W_{\text{stop}} &= w_u-(1+\lambda)w_u = -\lambda w_u. \label{r10e9-Ws}
\end{align}
The government bails out iff $y_L-\lambda s\geq-\lambda w_u\iff\lambda(s-w_u)\leq y_L$. The largest possible subsidy, with $L=0$, is $s=w+\Rc-y_L$, and then the condition reads
\begin{align}
  \lambda(w+\Rc-y_L-w_u)\leq y_L \iff \lambda[(w-w_u)+\Rc]\leq(1+\lambda)y_L, \label{r10e9-bail}
\end{align}
exactly the assumption. So the government bails out whatever $L$ is, with $s=\max\{0,L^*-L\}$.

\step{5}{Private choice of $L$} Every unit of liquidity not hoarded is replaced one-for-one by a subsidy, a gift to the entrepreneur and investors. They choose
\begin{align}
  L = 0, \qquad s = w+\Rc-y_L \text{ in the low state}. \label{r10e9-L0}
\end{align}

\step{6}{Commitment} If the government commits at date 0, it compares a subsidy (cost $\lambda s$ in deadweight) with private liquidity (free at $q=1$, and the firm can afford it by \eqref{r10e9-fin}). It commits to $s=0$ in both states; firms then hoard $L^*$ and still continue. Without commitment, the expected deadweight loss is $(1-\rho)\lambda(w+\Rc-y_L)$.

\step{7}{Example} $y_H=3$, $y_L=1$, $w=0.8$, $\Rc=0.5$, $\rho=0.5$, $I=1$, $A=0.5$, $\lambda=0.3$, $w_u=0.4$: pledgeable income $0.7\geq0.5$; $\text{NPV}(1)=0.2>\text{NPV}(0)=0.1$; $L^*=0.3$; \eqref{r10e9-bail}: $0.27\leq1.3$; with $L=0$, $W_{\text{continue}}-W_{\text{stop}}=(1-0.09)+0.12=1.03>0$; expected deadweight $0.5(0.3)(0.3)=0.045$.

\answer{(ii)} Without commitment: $L=0$ and the government pays $s=w+\Rc-y_L$ when $y=y_L$ (soft budget constraint). With commitment: $s=0$ in every state, and firms hoard $L^*=w+\Rc-y_L$.

\intuition The government cannot stand idle when a viable firm is about to lay off a worker, because shutting down costs output and unemployment benefits. Knowing this, the firm and its investors stop insuring themselves and let taxpayers provide the liquidity at a cost $1+\lambda$ per unit. A government that can tie its hands avoids the deadweight loss without losing a single job.

\pitfalls
\begin{itemize}
  \item Counting the wage or the subsidy as a social cost: both are transfers; only $\lambda s$ and $\lambda w_u$ are deadweight.
  \item Taking the liquidity need to be $w$: the firm can raise $y_L-\Rc$ from the market at date 1.
  \item Concluding that the bailout destroys value ex post: given $L=0$ it is optimal ex post; the inefficiency is ex ante.
\end{itemize}

\bookans The book gives no answer to this exercise. The result is the soft-budget-constraint logic of Tirole \S5.5 applied to a government (\S16.3.2): ex post rescue is optimal, so ex ante self-insurance is crowded out.
\refbib{Tirole \S5.5 ``Endogenous Liquidity Needs, the Sensitivity of Investment to Cash Flow, and the Soft Budget Constraint'', from p.~220; \S16.3.2, pp.~546--549; Exercise 15.2 (same technology).}

% =============================================================================
\exsec{10. News about prospects and aggregate liquidity --- Tirole 15.2, pp.~531--532}
%% BEGIN-STATEMENT 10
\stmtbegin
\textbf{Source:} Tirole, Ch. 15, Exercise 15.2, pp.~531--532 (PDF 544) · Answer in book: yes,
pp.~620--621

\textbf{Exercise 15.2 (news about prospects and aggregate liquidity).} Consider an economy with a
continuum of identical risk-neutral entrepreneurs. The representative entrepreneur has a fixed-size
investment project costing \(I\), and limited personal wealth \(A < I\). The project, if undertaken,
will deliver a random but verifiable income \(y \in [0, 1]\), with cumulative distribution function
\(G(y)\) and density \(g(y)\), provided that a reinvestment \(J\) is made after \(y\) is learned,
but before \(y\) is produced. The project yields nothing if it is interrupted.

Moreover, in the case of ``continuation'' (that is, if \(J\) is sunk), and regardless of the value
of \(y\), the entrepreneur may behave, in which case income is \(y\) for certain, or misbehave, in
which case income is \(y\) with probability \(p_L\) and 0 with probability \(1 - p_L\). The
entrepreneur, who is protected by limited liability, obtains private benefit \(B\) when misbehaving
(and no private benefit otherwise). Let

\[\mathcal{R} \equiv \frac{B}{1 - p_L}\]

(one will assume that \(B\) is small enough that, in the relevant range, it is worth inducing the
entrepreneur to behave in the case of continuation).

The timing is summarized in Figure 15.10.

Figure 15.10 (timeline, left to right; p.~531, PDF 544)

{\def\LTcaptype{none} % do not increment counter
\begin{longtable}[]{@{}
  >{\raggedright\arraybackslash}p{(\linewidth - 6\tabcolsep) * \real{0.2500}}
  >{\raggedright\arraybackslash}p{(\linewidth - 6\tabcolsep) * \real{0.2500}}
  >{\raggedright\arraybackslash}p{(\linewidth - 6\tabcolsep) * \real{0.2500}}
  >{\raggedright\arraybackslash}p{(\linewidth - 6\tabcolsep) * \real{0.2500}}@{}}
\toprule\noalign{}
\begin{minipage}[b]{\linewidth}\raggedright
Date 0
\end{minipage} & \begin{minipage}[b]{\linewidth}\raggedright
Date 1
\end{minipage} & \begin{minipage}[b]{\linewidth}\raggedright
Date 1 (continued)
\end{minipage} & \begin{minipage}[b]{\linewidth}\raggedright
Date 2
\end{minipage} \\
\midrule\noalign{}
\endhead
\bottomrule\noalign{}
\endlastfoot
Representative entrepreneur invests \(I\)? If so, she borrows \(I - A\). & Firms' idiosyncratic
productivities \(y\) revealed (drawn from \(G(\cdot)\)). Depending on its productivity realization,
each firm decides whether to continue (and spend \(J\)). & Moral hazard (in the case of
continuation): the entrepreneur behaves or misbehaves. & Realization of income (\(y\) or 0). \\
\end{longtable}
}

The rate of interest in the economy is equal to 0.

\begin{enumerate}
\def\labelenumi{(\roman{enumi})}
\item
  Compute the NPV and the investors' net income as functions of the threshold \(y^*\) for
  continuation.
\item
  Let \(y_0^* \equiv J\) and \(y_1^* = J + \mathcal{R}\).
\end{enumerate}

Define \(A_0^*\) and \(A_1^*\) by

\[I - A_k^* \equiv \int_{y_k^*}^1 y\,dG(y) - [1 - G(y_k^*)][J + \mathcal{R}],\]

for \(k \in \{0, 1\}\).

• What are the date-0 investment policy (investment/no investment) and the date-1 reinvestment
policy (the threshold \(y^*\)) as functions of \(A\)? (Hint: distinguish three regions
\(A \geqslant A_0^*\), \(A \leqslant A_1^*\), and \(A_1^* < A < A_0^*\).)

• Argue that, for \(A > A_1^*\), the entrepreneur must arrange at date 0 for her firm's date-1
liquidity.

\begin{enumerate}
\def\labelenumi{(\roman{enumi})}
\setcounter{enumi}{2}
\item
  • Is there enough inside liquidity if productivities are drawn \emph{independently} from the
  distribution \(G(\cdot)\)? Why?
\item
  Suppose, in contrast, that there is a macroeconomic shock \(\theta\) that is revealed at the
  beginning of date 1. (One will denote by \(E_\theta[\cdot]\) the date-0 expectations over the
  random variable \(\theta\).) Let \(y^*(\theta)\) denote the state-contingent threshold.
\end{enumerate}

• Write the date-0 financing constraint.

• Show that the optimal threshold when liquidity is abundant is actually state independent: there
exists \(y^*\) such that

\[y^*(\theta) = y^* \quad\text{for all } \theta.\]

• Show that the second-best allocation can be implemented when there are at least

\[\min_{\{\theta\}} \int_{y^*}^1 (y - J - \mathcal{R})\,dG(y \mid \theta)\]

units of outside liquidity delivering 1 unit of good for certain at date 1.

• What would happen if there were few such stores of value?

\paragraph{Context needed to solve it}

None beyond Figure 15.10, transcribed above.
\stmtend
%% END-STATEMENT

\asked (i) NPV and investors' net income as functions of the continuation threshold $y^*$. (ii) The investment decision and threshold for each wealth level, and why firms with $A>A_1^*$ need liquidity arranged at date 0. (iii) Whether firms' own claims provide enough liquidity with independent productivities. (iv) With a macro shock: the financing constraint, why the optimal threshold is the same in every state, how much outside liquidity is needed, and what happens if it is scarce.

\approach Fixed investment with a date-1 reinvestment $J$ (Chapter 5): continuing has NPV $y-J$ but pledgeable income $y-J-\Rc$. Weak firms must promise continuation only in high-$y$ states (raise $y^*$) to pledge more. Then Chapter 15's question: who supplies the date-1 cash, the firms themselves (inside liquidity) or outside stores of value?

\notation
\begin{itemize}
  \item $y\in[0,1]$, c.d.f.\ $G$, density $g$; $J$: reinvestment; $\Rc=B/(1-p_L)$: minimum stake; $y^*$: continue iff $y\geq y^*$.
  \item $A_0^*$, $A_1^*$: wealth cutoffs defined in the statement; $\theta$: macro shock, $G(\cdot\mid\theta)$; $\mu$: multiplier of the financing constraint.
\end{itemize}

\step{1}{NPV and investors' income} The incentive stake is $\Rc$ (as in exercise 9). With threshold $y^*$:
\begin{align}
  \text{NPV}(y^*) &= \int_{y^*}^1y\,dG(y)-[1-G(y^*)]J-I, \label{r10e10-npv}\\
  \text{Investors' net income} &= \int_{y^*}^1y\,dG(y)-[1-G(y^*)][J+\Rc]-(I-A). \label{r10e10-inv}
\end{align}

\answer{(i)} Equations \eqref{r10e10-npv} and \eqref{r10e10-inv}.

\step{2}{Maximisers} Differentiating, $d\,\text{NPV}/dy^*=-y^*g(y^*)+Jg(y^*)=(J-y^*)g(y^*)$: the NPV peaks at $y_0^*=J$. The investors' income derivative is $(J+\Rc-y^*)g(y^*)$: it peaks at $y_1^*=J+\Rc$. By the definitions, $I-A_k^*$ is the investors' gross income at $y_k^*$.

\step{3}{Policy by region}
(a) $A\geq A_0^*$: financing at $y^*=J$ is feasible; invest, $y^*=J$ (first best).
(b) $A_1^*<A<A_0^*$: the NPV falls in $y^*$ on $(J,J+\Rc)$ and investors' income rises there, so pick the smallest $y^*\in(J,J+\Rc)$ with
\begin{align}
  \int_{y^*}^1y\,dG(y)-[1-G(y^*)][J+\Rc] = I-A. \label{r10e10-bind}
\end{align}
(c) $A<A_1^*$: even $y^*=J+\Rc$ does not pledge enough: no investment ($A=A_1^*$: invest with $y^*=J+\Rc$).

\step{4}{Date-1 liquidity} For $A>A_1^*$, $y^*<J+\Rc$. In states $y\in[y^*,J+\Rc)$ the firm must spend $J$ but new investors would get at most $y-\Rc<J$: the market alone will not refinance. The firm needs up to $J+\Rc-y^*$ arranged at date 0 (credit line or reserves).

\step{5}{Example} $G$ uniform, $J=0.3$, $\Rc=0.2$, $I=0.5$: investors' gross income at $y^*$ is $(1-y^{*2})/2-(1-y^*)(0.5)$, so $A_0^*=0.5-0.105=0.395$ and $A_1^*=0.5-0.125=0.375$. For $A=0.385$, \eqref{r10e10-bind} is $y^*(1-y^*)/2=0.115$, i.e.\ $y^{*2}-y^*+0.23=0$, so $y^*=(1-\sqrt{0.08})/2=0.358579$; the largest liquidity need is $0.5-0.358579=0.141421$.

\answer{(ii)} Invest with $y^*=J$ if $A\geq A_0^*$; invest with $y^*\in(J,J+\Rc)$ solving \eqref{r10e10-bind} if $A_1^*<A<A_0^*$; do not invest if $A<A_1^*$. For $A>A_1^*$, liquidity of up to $J+\Rc-y^*$ must be arranged at date 0.

\step{6}{Independent productivities} With a continuum of independent firms, the realised distribution is $G$. The aggregate net pledgeable income at date 1 is, by \eqref{r10e10-bind},
\begin{align}
  \int_{y^*}^1(y-J-\Rc)\,dG(y) = I-A>0, \label{r10e10-agg}
\end{align}
so high-$y$ firms' excess pledgeable income covers low-$y$ firms' shortfalls. In the example: shortfall $\int_{y^*}^{0.5}(0.5-y)dy=(0.141421)^2/2=0.01$, surplus $\int_{0.5}^1(y-0.5)dy=0.125$, net $0.115=I-A$.

\answer{(iii)} Yes, provided liquidity is pooled: the corporate sector is a net borrower, so its aggregate date-1 pledgeable income is positive.

\step{7}{Macro shock: constraint and threshold} With $q=1$ (abundant stores of value), the date-0 financing constraint is
\begin{align}
  E_\theta\left[\int_{y^*(\theta)}^1(y-J-\Rc)\,dG(y\mid\theta)\right]\geq I-A. \label{r10e10-fc}
\end{align}
Maximise $E_\theta[\int_{y^*(\theta)}^1(y-J)\,dG(y\mid\theta)]-I$ subject to \eqref{r10e10-fc} with multiplier $\mu$. The derivative of the Lagrangian with respect to $y^*(\theta)$, state by state, is
\begin{align}
  -[y^*(\theta)-J]g(y^*\mid\theta)-\mu[y^*(\theta)-J-\Rc]g(y^*\mid\theta) &= 0 \nonumber\\
  (1+\mu)y^*(\theta) &= (1+\mu)J+\mu\Rc \nonumber\\
  y^*(\theta) &= J+\frac{\mu}{1+\mu}\Rc, \label{r10e10-ystar}
\end{align}
the same in every state.

\step{8}{Outside liquidity} In state $\theta$ the firms' aggregate net pledgeable income is $\int_{y^*}^1(y-J-\Rc)\,dG(y\mid\theta)$; if negative, the gap must be filled with outside stores of value. The plan is implementable with
\begin{align}
  L\geq\max\Bigl\{0,\,-\min_\theta\int_{y^*}^1(y-J-\Rc)\,dG(y\mid\theta)\Bigr\} \label{r10e10-L}
\end{align}
units (the book writes the minimum itself, meaning its absolute value when negative). Example: two equally likely states, $y\sim U[0,1]$ or $U[0,0.6]$, $A=0.45$: \eqref{r10e10-fc} with a common threshold gives $y^{*2}-y^*+0.225=0$, $y^*=(1-\sqrt{0.1})/2=0.341886$ (so $\mu/(1+\mu)=0.209431$); state-by-state net income $0.1125$ (good) and $-0.0125$ (bad), so $L\geq0.0125$. (Perturbing the thresholds across states while keeping \eqref{r10e10-fc} lowers the objective, as the script checks.)

\step{9}{Scarce liquidity} If fewer stores exist, they trade at $q>1$; the premium $(q-1)$ times the liquidity held enters the cost and the constraint (as in \S15.3.1). Firms then economise on liquidity: thresholds rise in the bad states (more interruption) and/or investment falls.

\answer{(iv)} \eqref{r10e10-fc}; $y^*(\theta)=J+\frac{\mu}{1+\mu}\Rc$ for all $\theta$; outside liquidity of at least the worst-state shortfall \eqref{r10e10-L}; with few stores, a liquidity premium $q>1$, more interruption and lower investment.

\intuition The entrepreneur trades NPV for pledgeability by interrupting moderately productive projects. Whether the date-1 cash can come from the firms themselves depends on diversification: idiosyncratic news washes out, a macro shock does not. The optimal continuation rule does not react to the macro state because the marginal trade-off between NPV and pledgeable income is the same in every state once liquidity is priced at par.

\pitfalls
\begin{itemize}
  \item Using $y^*=J$ for every $A$: it is optimal only when financing at $J$ is feasible.
  \item Assuming the market will refinance any $y\geq J$: investors need $y-\Rc\geq J$.
  \item Reading the book's $\min_\theta\int(\cdots)$ literally as a positive quantity of stores.
\end{itemize}

\bookans Matches the book's answer (pp.~620--621) in every part, including $y^*(\theta)=J+\mu\Rc/(1+\mu)$. The sign remark in Step~8 and the examples are additions.
\refbib{Tirole Ch.~5 (liquidity needs); \S15.2--15.3, pp.~518--525.}

% =============================================================================
\exsec{11. Market conditions with and without credit rationing --- Tirole 13.3, p.~494}
%% BEGIN-STATEMENT 11
\stmtbegin
\textbf{Source:} Tirole, Ch. 13, Exercise 13.3, p.~494 (PDF 507) · Answer in book: yes, pp.~617--618

\textbf{Exercise 13.3 (impact of market conditions with and without credit rationing).} This
analysis pursues that of Section 13.5.1. There, we compared the sensitivity of investment with the
output price (or installed-base investment) in the presence or absence of credit rationing, focusing
on either the fixed-investment variant or the constant-returns-to-scale variant. We now assume
decreasing returns to scale.

The representative entrepreneur (there is a unit mass of such entrepreneurs) has initial wealth
\(A\), is risk neutral and protected by limited liability, and invests \(I + K\), where \(I\) is the
scale of investment and \(K\) a fixed cost that is unrelated to scale. We assume that
\(K \geqslant A\), and so investors are unable to finance by themselves even a small investment.

An entrepreneur is successful with probability \(p\) and fails with probability \(1 - p\). We assume
that the shocks faced by the entrepreneurs are independent. This hypothesis is consistent with the
assumption made below that the output price is deterministic. When successful, the entrepreneur
produces \(R(I)\) units of a good (with \(R(0) = 0\), \(R' > 0\), \(R'' < 0\), \(R'(0) = \infty\),
\(R'(\infty) = 0\)); an entrepreneur who fails produces nothing. For concreteness, let

\[R(I) = I^\alpha, \quad\text{with } 0 < \alpha < 1.\]

As usual, the probability of success is endogenous: \(p \in \{p_L, p_H\}\). Misbehavior, \(p = p_L\)
(respectively, good behavior, \(p = p_H\)), brings about private benefit \(BI\) (respectively, no
private benefit). To prevent moral hazard, the entrepreneur must receive reward \(R_b\) in the case
of success, such that

\[(\Delta p)R_b \geqslant BI.\]

The product sells at price \(P\) per unit. Presumably, investors are risk neutral and demand rate of
return 0.

Suppose that the fixed cost \(K\) is ``not too large'' (so that the entrepreneur wants to invest in
the absence of credit rationing), and that

\[\frac{p_H B}{\Delta p} < \frac{1 - \alpha}{\alpha}.\]

\begin{enumerate}
\def\labelenumi{(\roman{enumi})}
\item
  Derive the first- and second-best investment levels as functions of \(P\). Show that they coincide
  for \(P \geqslant P_0\) for some \(P_0\).
\item
  Using a diagram, argue that there exists a region of output prices in which the second-best
  investment is more responsive than the first-best investment to the output price.
\item
  How would you analyze the impact of the existence of an installed-base level of investment
  \(I_0\)?
\end{enumerate}

\paragraph{Context needed to solve it}

\emph{From Tirole §13.5.1 ``Heuristics'', pp.~489--491 (PDF 502--504):}

To obtain some first intuition as to how past investment crowds out current investment, let us start
with a static model, with first a fixed investment, and then variable investment.

\emph{13.5.1.1 Fixed Investment Size.} Consider, thus, the fixed-investment model. There is a mass 1
of entrepreneurs. At investment cost \(I\), an entrepreneur can produce \(R\) units of a good with
probability \(p\) (and 0 units with probability \(1 - p\)). The final price per unit of output is
\(P\). Presumably, \(P\) depends on past industry investment, but we do not need to go into detail
at this stage. The probability of success is \(p_H\) if the entrepreneur behaves (no private
benefit) and \(p_L = p_H - \Delta p\) if she misbehaves (private benefit \(B\)). We assume that the
output realizations are independent across entrepreneurs (there is no aggregate uncertainty); this
assumption is consistent with the assumption made above that the output price is deterministic.

Entrepreneurs are risk neutral and are protected by limited liability. The distribution of assets in
the population of entrepreneurs is given by the cumulative distribution function \(G(A)\) on
\([0, \infty)\). Investors are risk neutral and demand rate of return equal to 0. Assume that it is
optimal to provide the entrepreneurs with incentives to behave.

Varying \(P\), let us compare the level of aggregate investment under credit rationing (\(B > 0\))
and in its absence (\(B = 0\)), and show that the first- and second-best levels of investment are as
depicted in Figure 13.11.

Figure 13.11 (p.~490, PDF 503): investment \(I\) (vertical axis) against the price \(P\) (horizontal
axis). \(I^{\text{FB}}\) is 0 for \(P < I/p_H R\) and jumps to the flat level \(I\) at
\(P = I/p_H R\). Below \(I/p_H R\), \(I^{\text{SB}} = I^{\text{FB}} = 0\). Above it,
\(I^{\text{SB}}\) starts on the vertical segment at \(P = I/p_H R\) and rises, concave, toward
\(I^{\text{FB}}\) from below.

Conditional on receiving financing, an entrepreneur's incentive constraint is
\((\Delta p)R_b \geqslant B\), and so the financing condition can be written as
\(A \geqslant \bar A\), where

\[p_H\left[PR - \frac{B}{\Delta p}\right] = I - \bar A.\]

The fraction of entrepreneurs who are able to raise funds is

\[X \equiv 1 - G\left(I - p_H\left[PR - \frac{B}{\Delta p}\right]\right), \qquad (13.13)\]

or, equivalently, total investment is \(I^{\text{SB}} = XI\), as long as the NPV is positive, i.e.,
as long as \(p_H PR - I \geqslant 0\).

\emph{In the positive-NPV range} (\(P \geqslant I/p_H R\)), aggregate investment increases with
\(P\) as long as \(B > 0\). {[}\ldots{]} \emph{aggregate investment is always more responsive to the
output price under credit rationing than with no (or little) credit rationing}, since investment
does not move with the output price in this range in the absence of credit rationing.

This illustrates the \emph{net worth effect} {[}\ldots{]}.

To complete this analysis, we can now endogenize the product price by assuming the existence of a
prior ``fraction'' \(X_0\) of similar firms that were able to raise financing in the past. The
output price is then a decreasing function of total output, \(p_H(X_0 + X)R\):

\[P = P(p_H(X_0 + X)R) \quad\text{with } P' < 0. \qquad (13.14)\]

An equilibrium is then a ``level of investment'' \(XI\), where \(X\) is obtained from (13.13) and
(13.14):

\[X = 1 - G\left(I - p_H\left(P(p_H(X_0 + X)R)R - \frac{B}{\Delta p}\right)\right).\]

Note that as \(X_0\) grows, \(X\) must decrease (a crowding-out effect (if \(X\) increased, then
\(P\) would decrease, and so \(X\) would decrease after all)) but \(X_0 + X\) must increase (there
is less than full crowding out (if \(X_0 + X\) decreased, then \(P\) would increase, and \(X\) would
increase after all)).

The increased sensitivity of investment under credit rationing, however, is not a general
conclusion, as can also be seen from Figure 13.11. A small increase in \(P\) from the negative-NPV
to the positive-NPV region raises the first-best investment dramatically, but the second-best one
less so as not all firms get on board. {[}\ldots{]}

\emph{13.5.1.2 Variable Investment Size.} Next consider the variable-investment model of Section
3.4: the number of units of output produced in the case of success is \(RI\), the private benefit in
the case of misbehavior \(BI\). To be incentivized to behave, an entrepreneur with investment size
\(I\) must receive \(R_b\) in the case of success, such that \((\Delta p)R_b \geqslant BI\). For a
given output price \(P\), an entrepreneur with assets \(A\) can borrow up to the level at which
pledgeable income is equal to investors' outlay:

\[p_H\left(PR - \frac{B}{\Delta p}\right)I = I - A \quad\text{or}\quad I = \frac{A}{1 + (p_H B/\Delta p) - p_H PR}.\]

Assuming, without loss of generality, that all entrepreneurs have the same net worth \(A\), let us
analyze the impact of a prior level of investment \(I_0\) on current investment. Total output is
then \(p_H R(I_0 + I)\). Because \(P = P(p_H R(I_0 + I))\), with \(P' < 0\),

\[I = I^{\text{SB}} = \frac{A}{1 + (p_H B/\Delta p) - p_H P(p_H R(I_0 + I))R}. \qquad (13.15)\]

By the same reasoning as in the fixed-investment version, condition (13.15) implies that previous
investment partially crowds out current investment:
\(-1 < \partial I^{\text{SB}}/\partial I_0 < 0\).

Let us now compare this sensitivity to that obtained in the absence of credit rationing. In this
first-best benchmark, firms maximize their NPV, regardless of their solvency:
\(\max_I \{(p_H PR - 1)I\}\). Competitive equilibrium in this constant-returns-to-scale environment
implies that unit revenue is equal to unit cost, or

\[p_H P(p_H R(I_0 + I))R = 1. \qquad (13.16)\]

Thus, in the absence of credit constraint, past investment fully crowds out current investment:
\(\partial I/\partial I_0 = -1\). This is due to what might be labeled a \emph{hindering effect} of
credit rationing: because part of the benefit from investment expansion accrues to the entrepreneurs
and is therefore nonpledgeable, investors are less keen than entrepreneurs to expand as the market
becomes more profitable. And so credit rationing may make investment relatively less responsive to
market conditions.

Exercise 13.3 pursues this analysis in the intermediate context of variable investment and
decreasing returns to scale.
\stmtend
%% END-STATEMENT

\asked (i) First-best and second-best investment as functions of the output price $P$, and the price $P_0$ above which they coincide. (ii) That for some prices the second-best investment responds more to $P$ than the first-best. (iii) How an installed base $I_0$ (which lowers $P$) would be analysed.

\approach Variable investment with decreasing returns and a fixed cost $K\geq A$. First best: maximise NPV. Second best: maximise NPV subject to pledgeable income $\geq$ investors' outlay (Tirole \S13.5.1). Comparing slopes $dI/dP$ at the point where the constraint starts to bind settles (ii) analytically.

\notation
\begin{itemize}
  \item $R(I)=I^\alpha$, $0<\alpha<1$; $P$: output price; $K$: fixed cost; $A\leq K$: wealth.
  \item $b\equiv p_HB/\Delta p<(1-\alpha)/\alpha$; $I^{\text{FB}}(P)$, $I^{\text{SB}}(P)$: first- and second-best investment.
  \item $\Pc(I)\equiv p_H[PR(I)-BI/\Delta p]$: pledgeable income.
\end{itemize}

\step{1}{First best} $U_b=p_HPR(I)-I-K$; first-order condition
\begin{align}
  p_HP\alpha I^{\alpha-1} &= 1 \;\Longrightarrow\; I^{\text{FB}}(P)=(\alpha p_HP)^{1/(1-\alpha)}, \label{r10e11-FB}
\end{align}
provided $K\leq p_HPR(I^{\text{FB}})-I^{\text{FB}}$. By \eqref{r10e11-FB}, $p_HP(I^{\text{FB}})^\alpha=I^{\text{FB}}/\alpha$, so the gross NPV is $I^{\text{FB}}/\alpha-I^{\text{FB}}=\frac{1-\alpha}{\alpha}I^{\text{FB}}$.

\step{2}{Financing condition at the first best} The incentive constraint $(\Delta p)R_b\geq BI$ leaves $\Pc(I)$ to investors, who need $I+K-A$. At $I^{\text{FB}}$:
\begin{align}
  \Pc(I^{\text{FB}})-I^{\text{FB}} &= p_HP(I^{\text{FB}})^\alpha-bI^{\text{FB}}-I^{\text{FB}} = \frac{I^{\text{FB}}}{\alpha}-bI^{\text{FB}}-I^{\text{FB}} = \left[\frac{1-\alpha}{\alpha}-b\right]I^{\text{FB}}, \label{r10e11-gap}
\end{align}
positive by assumption and increasing in $P$. The first best is financeable iff $\Pc(I^{\text{FB}})-I^{\text{FB}}\geq K-A$, i.e.\ iff $P\geq P_0$ with
\begin{align}
  I^{\text{FB}}(P_0) = \frac{K-A}{(1-\alpha)/\alpha-b}, \qquad P_0=\frac{[I^{\text{FB}}(P_0)]^{1-\alpha}}{\alpha p_H}. \label{r10e11-P0}
\end{align}

\step{3}{Second best below $P_0$} For $P<P_0$ the constraint binds; NPV is concave and increasing below $I^{\text{FB}}$, so $I^{\text{SB}}$ is the largest root below $I^{\text{FB}}$ of
\begin{align}
  p_H\left[PI^\alpha-\frac{B}{\Delta p}I\right] = I+K-A. \label{r10e11-SB}
\end{align}
(If $\max_I[\Pc(I)-I]<K-A$, there is no investment.)

\answer{(i)} $I^{\text{FB}}=(\alpha p_HP)^{1/(1-\alpha)}$; $I^{\text{SB}}=I^{\text{FB}}$ for $P\geq P_0$ defined by \eqref{r10e11-P0}, and $I^{\text{SB}}<I^{\text{FB}}$ solves \eqref{r10e11-SB} for $P<P_0$.

\step{4}{Slopes} From \eqref{r10e11-FB}: $dI^{\text{FB}}/dP=I^{\text{FB}}/[(1-\alpha)P]$. Differentiating \eqref{r10e11-SB} totally:
\begin{align}
  p_HI^\alpha\,dP+\bigl[p_HP\alpha I^{\alpha-1}-b-1\bigr]dI &= 0 \nonumber\\
  \frac{dI^{\text{SB}}}{dP} &= \frac{p_HI^\alpha}{1+b-\alpha p_HPI^{\alpha-1}}, \label{r10e11-dSB}
\end{align}
with a positive denominator at the largest root. At $P=P_0$, $\alpha p_HPI^{\alpha-1}=1$ and $p_HI^\alpha=I/(\alpha P)$, so
\begin{align}
  \frac{dI^{\text{SB}}}{dP}\Big|_{P_0^-} = \frac{I}{\alpha Pb}, \qquad \frac{dI^{\text{SB}}/dP}{dI^{\text{FB}}/dP}\Big|_{P_0} = \frac{1-\alpha}{\alpha b}>1, \label{r10e11-ratio}
\end{align}
by the assumption $b<(1-\alpha)/\alpha$. By continuity the second best is steeper on an interval of prices just below $P_0$; the diagram is the book's Figure~7: $I^{\text{SB}}$ lies below $I^{\text{FB}}$ for $P<P_0$ and joins it at $P_0$ from below, hence with a steeper slope.

\step{5}{Example} $\alpha=0.5$, $p_H=0.8$, $B/\Delta p=0.5$ ($b=0.4<1$), $K=0.5$, $A=0.1$: $I^{\text{FB}}=(0.4P)^2$; $I^{\text{FB}}(P_0)=0.4/0.6=2/3$ and $P_0=\sqrt{2/3}/0.4=2.041241$ (gross NPV $2/3\geq K$). Slopes at $P_0$: FB $0.653197$, SB $1.632993$ (ratio $2.5$). At $P=1.9$: $I^{\text{FB}}=0.5776$, $I^{\text{SB}}=0.406557$, slopes $0.608$ vs.\ $2.451628$.

\answer{(ii)} Just below $P_0$, $dI^{\text{SB}}/dP=\frac{1-\alpha}{\alpha b}\,dI^{\text{FB}}/dP>dI^{\text{FB}}/dP$.

\step{6}{Installed base} Let $P=P(p_H[R(I)+R(I_0)])$, $P'<0$, and $I=h(P)$ with $h=I^{\text{FB}}$ or $I^{\text{SB}}$, $h'>0$. Differentiating $I=h(P(p_H[R(I)+R(I_0)]))$:
\begin{align}
  \frac{dI}{dI_0} = \frac{h'P'p_HR'(I_0)}{1-h'P'p_HR'(I)} = -\frac{h'|P'|p_HR'(I_0)}{1+h'|P'|p_HR'(I)}<0, \label{r10e11-crowd}
\end{align}
and total output rises: $\frac{d}{dI_0}p_H[R(I)+R(I_0)]=\frac{p_HR'(I_0)}{1+h'|P'|p_HR'(I)}>0$. The crowding-out $|dI/dI_0|$ increases with $h'$, so it is stronger under credit rationing exactly where $I^{\text{SB}}$ is steeper (prices just below $P_0$).

\answer{(iii)} Endogenise $P=P(p_H[R(I)+R(I_0)])$: past investment crowds out current investment partially (\eqref{r10e11-crowd}), total output still rises, and crowding out is stronger where the second best is more price-sensitive.

\intuition A higher price raises both NPV and pledgeable income. Without rationing only the NPV effect matters. With rationing the price also relaxes the financing constraint --- the net-worth effect --- which near the point where the constraint stops binding makes investment more responsive. The book's \S13.5.1 also shows the opposite ``hindering'' effect under constant returns; decreasing returns let both appear.

\pitfalls
\begin{itemize}
  \item Claiming the second best is always more responsive: the result holds in a region (below $P_0$), not everywhere.
  \item Using the smaller root of \eqref{r10e11-SB}: the entrepreneur wants the largest feasible scale below $I^{\text{FB}}$.
  \item Forgetting the fixed cost $K$ in the investors' outlay.
\end{itemize}

\bookans Matches the book's answer (pp.~617--618): same FB condition, same expression $[\Pc(I^{\text{FB}})-I^{\text{FB}}]=(1-\alpha)[\frac1\alpha-\frac{b}{1-\alpha}]I^{\text{FB}}$, the region to the left of $P_0$ in Figure~7, and the crowding-out argument with $P=P(p_H[R(I)+R(I_0)])$. The slope ratio \eqref{r10e11-ratio} is an addition that makes the diagram argument exact.
\refbib{Tirole \S13.5.1, pp.~489--491.}

% =============================================================================
\exsec{12. Downsizing and aggregate liquidity --- Tirole 15.1, pp.~530--531}
%% BEGIN-STATEMENT 12
\stmtbegin
\textbf{Source:} Tirole, Ch. 15, Exercise 15.1, pp.~530--531 (PDF 543) · Answer in book: yes,
pp.~619--620

\textbf{Exercise 15.1 (downsizing and aggregate liquidity).} Consider the variable-investment model
with \emph{decreasing} returns to scale and a liquidity shock. There is a unit mass of identical
entrepreneurs. The timing for a given entrepreneur is in Figure 15.9.

Figure 15.9 (timeline, left to right; p.~530, PDF 543)

{\def\LTcaptype{none} % do not increment counter
\begin{longtable}[]{@{}
  >{\raggedright\arraybackslash}p{(\linewidth - 6\tabcolsep) * \real{0.2500}}
  >{\raggedright\arraybackslash}p{(\linewidth - 6\tabcolsep) * \real{0.2500}}
  >{\raggedright\arraybackslash}p{(\linewidth - 6\tabcolsep) * \real{0.2500}}
  >{\raggedright\arraybackslash}p{(\linewidth - 6\tabcolsep) * \real{0.2500}}@{}}
\toprule\noalign{}
\begin{minipage}[b]{\linewidth}\raggedright
Date 0
\end{minipage} & \begin{minipage}[b]{\linewidth}\raggedright
Date 1
\end{minipage} & \begin{minipage}[b]{\linewidth}\raggedright
Date 1 → 2
\end{minipage} & \begin{minipage}[b]{\linewidth}\raggedright
Date 2
\end{minipage} \\
\midrule\noalign{}
\endhead
\bottomrule\noalign{}
\endlastfoot
Entrepreneur has wealth \(A\), borrows \(I - A\). & Firm is intact (no reinvestment needed) with
probability \(1 - \lambda\), and distressed (reinvestment \(\rho\) per salvaged unit) with
probability \(\lambda\). & Moral hazard. & Outcome: success (profit \(R(J)\)) with probability
\(p_\bullet\), failure (profit 0) with probability \(1 - p_\bullet\). \\
\end{longtable}
}

At date 1, an amount \(J\), \(0 \leqslant J \leqslant I\), is rescued. In the absence of a liquidity
shock (event has probability \(1 - \lambda\)), of course \(J = I\). But in the face of a liquidity
shock (which has probability \(\lambda\)), the investment is downsized to \(J \leqslant I\) (the
cost of continuation is then \(\rho J\)). The shock is verifiable. Let \(R(J)\) denote the profit in
the case of success.

The moral-hazard stage is described as it usually is: the probability of success is \(p_H\) if the
entrepreneur works and \(p_L = p_H - \Delta p\) if she shirks. The entrepreneur obtains private
benefit \(BJ\) by misbehaving and 0 otherwise. Investors and entrepreneur are risk neutral, and the
latter is protected by limited liability.

Economic agents do not discount the future (which does not imply that rates of interest are always
0!).

From now on, use \(J\) for the amount that is salvaged when there is a liquidity shock (as we noted,
the corresponding amount is \(I\) in the absence of shock).

Assume that \(R(0) = 0\), \(R' > 0\), \(R'' < 0\), \(R'(0) = \infty\), \(R'(\infty) = 0\).

\begin{enumerate}
\def\labelenumi{(\roman{enumi})}
\tightlist
\item
  Assume that there is plenty of liquidity in the economy, so that the firms have access to a store
  of value (by paying \(q = 1\) at date 0, they receive 1 at date 1).
\end{enumerate}

Show that downsizing occurs in the case of a liquidity shock,

\[J^* < I^*,\]

if and only if

\[\rho > \frac{1}{1 - \lambda}.\]

(Hints: (1) write the incentive constraints (the sharing rule can be adjusted to the realization of
the shock) and infer the pledgeable income; (2) maximize the entrepreneur's utility (employ the
usual trick) subject to the investors' breakeven condition, ignoring the constraint
\(J \leqslant I\); let \(\mu\) denote the shadow price of the constraint; (3) derive the stated
result.)

\begin{enumerate}
\def\labelenumi{(\roman{enumi})}
\setcounter{enumi}{1}
\tightlist
\item
  Suppose that the liquidity shocks are \emph{perfectly correlated.}
\end{enumerate}

• What is the minimal number \(L^*\) of outside stores of value (delivering 1 unit of good each at
date 1) needed to support the allocation described in (i)?

• Argue that if \(L < L^*\), then \(q > 1\) and \(J < I\) \emph{a fortiori} if
\(\rho > 1/(1 - \lambda)\).

Derive the equations giving the liquidity premium \((q - 1)\) under these assumptions.

\begin{enumerate}
\def\labelenumi{(\roman{enumi})}
\setcounter{enumi}{2}
\tightlist
\item
  Suppose now that the liquidity shocks are independent across firms.
\end{enumerate}

• Argue that (provided that the entrepreneurs borrow at date 0) there is enough liquidity to support
the allocation derived in (i).

• Suppose that each entrepreneur holds the stock index. When will this provide enough liquidity? How
can one prevent this potential waste of liquidity?

\paragraph{Context needed to solve it}

\emph{From Tirole §15.2.2 ``Wasting Inside Liquidity'', pp.~521--522 (PDF 534--535):}

The sufficiency result by itself only states that the corporate sector produces enough inside
liquidity to support its optimal reinvestment policy. It is silent on how the latter can be
implemented.

Let us first point out that the ``natural implementation,'' namely, the policy that consists in each
firm holding the stock index, that is, a representative portfolio of claims on all firms in the
economy, in general does not work. The date-1 value of the index is \((\rho_0 - \bar\rho)I\). To see
this, note that the corporate sector will return \(\rho_0 I\) to investors at date 2. The average
reinvestment cost, however, is \(((1 - \lambda)\rho_L + \lambda\rho_H)I = \bar\rho I\), which must
be financed at date 1 by issuing new shares on the corporate sector date-2 income and thereby
diluting existing shareholders.

Thus, if all firms hold equal shares in the index, those with a high liquidity shock can meet it by
reselling shares if and only if

\[(\rho_H - \rho_0)I \leqslant (\rho_0 - \bar\rho)I \quad\text{or}\quad \rho_H + \bar\rho \leqslant 2\rho_0. \qquad (15.4)\]

If \(\rho_H + \bar\rho > 2\rho_0\), then holding the stock index does not allow the firms facing the
high liquidity shock to continue.

Why does this ``self-provision'' of liquidity result in a waste of liquidity? Firms that face a low
liquidity shock at date 1 have excess cash for two reasons: first (reasoning in terms of per unit of
investment), they can raise up to \(\rho_0\) when they need only \(\rho_L\); second, they have
invested in the stock index, with resulting value \(\rho_0 - \bar\rho\). This excess liquidity is,
of course, not fully wasted as the extra profit is partly reappropriated by distressed firms that
own a fraction of the stock index and therefore own part of the healthy firms. Still, healthy firms
do have excess liquidity. At date 1, they either redistribute the excess cash to or invest it on
behalf of their owners; they have no incentive to invest in distressed firms, in which the
reinvestment cost exceeds the pledgeable income.

{[}\ldots{]} in order to force healthy firms to redispatch cash to distressed ones at date 1, firms
must at date 0 pool their liquidity and organize a system of (or akin to) credit lines. For example,
shares are deposited with one (or an arbitrary number of) financial institutions; each firm is then
entitled to draw on a credit line up to a cap of \(\rho_H I\). The financial institution can raise
the cash needed to honor the credit lines by selling at date 1 shares it holds in firms to date-1
investors. From the net-borrower assumption, the proceeds of such sales can cover the financing of
the credit lines: \((1 - \lambda)(\rho_0 - \rho_L)I > \lambda(\rho_H - \rho_0)I\).

(In this passage \(\rho_L\), \(\rho_H\), \(\bar\rho = (1 - \lambda)\rho_L + \lambda\rho_H\),
\(\rho_0\) are the per-unit shocks and pledgeable income of the constant-returns model of §15.2.1;
see exercise 3.)

\emph{From Tirole §15.3.1 ``Aggregate Shocks and the Value of Outside Liquidity'', p.~523 (PDF
536):}

{[}\ldots{]} Here we introduce outside stores of value into the picture. For simplicity, there are
\(L^S\) such stores of value. A store of value yields 1 unit of good at date 1 for certain. Its
date-0 price is \(q\); because the market rate of interest is equal to 0 and investors can buy the
store of value, \(q \geqslant 1\). Furthermore, if \(q = 1\), then the corporate sector is able to
implement its efficient allocation. If \(q > 1\), then the corporate sector must hold all stores of
value. {[}\ldots{]} \(q > 1\) only if firms compete with each other for the scarce liquidity. The
date-0 price \(q\) adjusts to the level at which the demand for liquid assets is equal to the
supply: \(L = L^S\).
\stmtend
%% END-STATEMENT

\asked (i) With abundant liquidity, show that distressed firms are downsized ($J^*<I^*$) iff $\rho>1/(1-\lambda)$. (ii) Under perfectly correlated shocks: the outside liquidity $L^*$ needed, why scarcer liquidity means $q>1$ and more downsizing, and the equations for the premium. (iii) Under independent shocks: why inside liquidity suffices, when holding the index suffices, and how to avoid wasting liquidity.

\approach Chapter 5's liquidity model with decreasing returns: maximise NPV subject to investors' breakeven, with state-contingent incentive pay. The comparison of the two first-order conditions is all of (i). (ii) prices liquidity as in \S15.3.1; (iii) uses \S15.2's sufficiency result and the stock-index argument of \S15.2.2.

\notation
\begin{itemize}
  \item $I$: investment (kept intact w.p.\ $1-\lambda$); $J\leq I$: amount salvaged after a shock (w.p.\ $\lambda$), cost $\rho J$.
  \item $R(\cdot)$: profit if success; $b\equiv p_HB/\Delta p$; $\Pc(X)\equiv p_H[R(X)-BX/\Delta p]$: pledgeable income of $X$ units.
  \item $\mu$: multiplier of the breakeven constraint, $m\equiv\mu/(1+\mu)$; $q$: price of a store of value; $N(J)\equiv\rho J-\Pc(J)$: net date-1 need in the shock state.
\end{itemize}

\step{1}{Incentives and breakeven} Pay can depend on the (verifiable) shock: $(\Delta p)R_b^0\geq BI$ without shock, $(\Delta p)R_b^\rho\geq BJ$ with shock. Pledgeable incomes are $\Pc(I)$ and $\Pc(J)$. Investors' breakeven and the entrepreneur's utility (= NPV):
\begin{align}
  (1-\lambda)\Pc(I)+\lambda[\Pc(J)-\rho J] &\geq I-A, \label{r10e12-be}\\
  U_b &= (1-\lambda)p_HR(I)+\lambda[p_HR(J)-\rho J]-I. \label{r10e12-Ub}
\end{align}

\step{2}{First-order conditions} The Lagrangian is $\mathcal L=U_b+\mu\{(1-\lambda)\Pc(I)+\lambda[\Pc(J)-\rho J]-I+A\}$. With respect to $I$:
\begin{align}
  (1-\lambda)p_HR'(I)-1+\mu\Bigl[(1-\lambda)p_H\Bigl(R'(I)-\frac{B}{\Delta p}\Bigr)-1\Bigr] &= 0 \nonumber\\
  (1+\mu)[(1-\lambda)p_HR'(I)-1] &= \mu(1-\lambda)b \nonumber\\
  p_HR'(I) &= \frac{1}{1-\lambda}+m\,b. \label{r10e12-focI}
\end{align}
With respect to $J$ (ignoring $J\leq I$):
\begin{align}
  \lambda[p_HR'(J)-\rho]+\mu\lambda\Bigl[p_H\Bigl(R'(J)-\frac{B}{\Delta p}\Bigr)-\rho\Bigr] &= 0 \nonumber\\
  (1+\mu)[p_HR'(J)-\rho] &= \mu b \nonumber\\
  p_HR'(J) &= \rho+m\,b. \label{r10e12-focJ}
\end{align}

\step{3}{Comparison} $R'$ is decreasing, so $J<I\iff R'(J)>R'(I)\iff\rho+mb>\frac1{1-\lambda}+mb\iff\rho>\frac1{1-\lambda}$. If $\rho\leq1/(1-\lambda)$ the unconstrained $J$ would exceed $I$, so $J=I$. The agency term $mb$ is the same in both states and cancels.

\answer{(i)} $J^*<I^*$ iff $\rho>1/(1-\lambda)$: continuing a unit after the shock costs $\rho$, while a unit of initial investment costs $1/(1-\lambda)$ per unit that survives intact.

\step{4}{Correlated shocks: $L^*$} In the shock state every firm is hit; no healthy firm can supply liquidity. Each firm can still raise $\Pc(J)$ by issuing new claims to date-1 investors, so it needs outside stores for the rest:
\begin{align}
  L^* = N(J^*) = \rho J^*-\Pc(J^*) \quad\text{(if positive; } L^*=0 \text{ otherwise)}. \label{r10e12-L}
\end{align}
The book's answer counts the whole reinvestment, $L^*=\rho J^*$, which corresponds to firms unable to raise anything at date 1.

\step{5}{Scarce liquidity} If $L^S<L^*$, firms compete for stores and $q>1$ (otherwise all firms could buy $L^*$). Buying $N(J)$ stores costs a premium $(q-1)N(J)$ at date 0, subtracted from $U_b$ and added to investors' outlay. The new Lagrangian derivative in $J$, using $\Pc'(J)-\rho=-N'(J)$:
\begin{align}
  \lambda[p_HR'(J)-\rho]-(1+\mu)(q-1)N'(J)-\mu\lambda N'(J) &= 0. \nonumber
\end{align}
Substituting $p_HR'(J)-\rho=b-N'(J)$ and solving for $N'(J)$:
\begin{align}
  \lambda b &= N'(J)[\lambda(1+\mu)+(1+\mu)(q-1)] \nonumber\\
  p_HR'(J) &= \rho+b\left[1-\frac{\lambda}{(1+\mu)(\lambda+q-1)}\right], \tag{4$''$}\label{r10e12-focJq}
\end{align}
which reduces to \eqref{r10e12-focJ} at $q=1$ and has a larger right side for $q>1$: $J$ falls. \eqref{r10e12-focI} is unchanged given $\mu$, so $J<I$ \emph{a fortiori} when $\rho>1/(1-\lambda)$. The equations for the premium are: the modified breakeven $(1-\lambda)\Pc(I)+\lambda[\Pc(J)-\rho J]-(q-1)N(J)=I-A$, \eqref{r10e12-focI}, \eqref{r10e12-focJq}, and market clearing $N(J)=L^S$; four equations in $I,J,\mu,q$. (Book convention $L=\rho J$: $p_HR'(J)=\rho[1+(q-1)/\lambda]+mb$ and $\rho J=L^S$.)

\step{6}{Independent shocks} Exactly a fraction $\lambda$ is hit. At $q=1$, \eqref{r10e12-be} with $I>A$ gives
\begin{align}
  V\equiv(1-\lambda)\Pc(I)+\lambda\Pc(J) > \lambda\rho J, \label{r10e12-V}
\end{align}
so the date-1 value of the corporate sector covers the average reinvestment: no outside liquidity is needed if it is pooled.

\step{7}{Holding the index} A distressed firm raises $\Pc(J)$ on its own and sells its index share, worth $V-\lambda\rho J$ after the reinvestments are financed by dilution. It needs $\rho J$:
\begin{align}
  \Pc(J)+V-\lambda\rho J &\geq \rho J \nonumber\\
  (1-\lambda)\Pc(I)+(1+\lambda)\Pc(J) &\geq (1+\lambda)\rho J \nonumber\\
  (1-\lambda)\Pc(I) &\geq (1+\lambda)[\rho J-\Pc(J)], \label{r10e12-index}
\end{align}
which is not guaranteed. Remedy: pool liquidity through an intermediary and credit lines (\S15.2.2).

\step{8}{Example} $R(X)=2X^{0.8}$, $p_H=0.9$, $B/\Delta p=1$ ($b=0.9$), $\lambda=0.3$, $A=0.2$, $\rho=1.6>1/0.7$: solving \eqref{r10e12-be}, \eqref{r10e12-focI}, \eqref{r10e12-focJ} gives $m=0.160310$, $I^*=0.643244$, $J^*=0.383471<I^*$. With $\rho=1.3<1/0.7$ the unconstrained solution would be $J=0.834348>I=0.552051$, so $J=I$. $L^*=N(J^*)=0.122577$ (book convention $\rho J^*=0.613554$). At $q=1.1$, \eqref{r10e12-focJq} gives $J=0.239918$ (a brute-force maximisation in the script agrees). Index condition \eqref{r10e12-index}: $0.480017\geq0.159350$, satisfied here.

\answer{(ii)--(iii)} $L^*=\rho J^*-\Pc(J^*)$ (book: $\rho J^*$); if $L^S<L^*$, $q>1$ and $J$ falls further, with the premium given by the four equations of Step~5. Independent shocks: enough inside liquidity by \eqref{r10e12-V}; the index suffices iff \eqref{r10e12-index}; otherwise pool via credit lines.

\intuition Downsizing is the decreasing-returns version of Chapter 5's ``do not continue at any cost'': a unit rescued after the shock is worth it only if it is cheaper than a unit of initial scale. Aggregate correlation turns the rescue into a demand for outside liquidity, whose price rises when it is scarce and makes rescues smaller still.

\pitfalls
\begin{itemize}
  \item Using the same incentive pay in both states: pay can be adjusted to the verifiable shock.
  \item Forgetting that the agency term is identical in \eqref{r10e12-focI} and \eqref{r10e12-focJ}, which is why it cancels.
  \item Treating the premium as paid only in the shock state: stores are bought at date 0, so $(q-1)N(J)$ is not multiplied by $\lambda$.
\end{itemize}

\bookans Matches the book's answer (pp.~619--620) for (i) [conditions (3)--(4) and $\rho>1/(1-\lambda)$] and (iii) [inequality (6) and the index condition]. For (ii) the book takes $L^*=\rho J^*$ and $p_HR'(J)=\rho(1+\frac{q-1}{\lambda})+\frac{\mu}{1+\mu}p_HB/\Delta p$; this counts the entire reinvestment as outside liquidity. Allowing distressed firms to issue claims on $\Pc(J)$ at date 1, as \S15.3.1 does, gives \eqref{r10e12-L} and \eqref{r10e12-focJq}; both versions give $q>1$ and $J<I$ \emph{a fortiori}.
\refbib{Tirole Ch.~5 (liquidity needs); \S15.2--15.3, pp.~518--525.}

% =============================================================================
\exsec{13. Complementarity between liquid and illiquid assets --- Tirole 15.4, p.~532}
%% BEGIN-STATEMENT 13
\stmtbegin
\textbf{Source:} Tirole, Ch. 15, Exercise 15.4, p.~532 (PDF 545) · Answer in book: yes, p.~621

\textbf{Exercise 15.4 (complementarity between liquid and illiquid assets).} Go through the analysis
of Section 15.3.1 assuming that entrepreneurs do not want to invest in projects that are
discontinued in the adverse state of nature:

\[(1 - \lambda)\rho_1 < 1 + (1 - \lambda)\rho_L.\]

Show that an increase in the supply \(L^S\) of liquid assets increases the investment \(I\) in
illiquid ones.

\paragraph{Context needed to solve it}

\emph{From Tirole §15.2.1 ``The Sufficiency Result'', pp.~518--520 (PDF 531--533)} --- the model of
Section 15.3.1 is that of Section 15.2:

There are three periods, \(t = 0, 1, 2\). Investors are risk neutral and the market rate of interest
in the economy is 0. The economy is also populated by a large number---technically a continuum of
mass 1---of ex ante identical, risk-neutral entrepreneurs. The representative entrepreneur at date 0
has wealth \(A\), borrows \(I - A\), and invests \(I\). At date 1, a given firm faces liquidity
shock \(\rho I\), with \(\rho = \rho_L\) with probability \(1 - \lambda\) (healthy firm) and
\(\rho = \rho_H > \rho_L\) with probability \(\lambda\) (firm in distress). The firm can continue
only if it finds funds to defray its liquidity shock; otherwise, it is liquidated. We normalize the
liquidation value at 0. In the case of continuation, the firm's date-2 expected income is denoted by
\(\rho_1 I\), of which only \(\rho_0 I < \rho_1 I\) is pledgeable to investors (see Figure 15.1).

Figure 15.1 (p.~519, PDF 532; identical to Figure 15.11): Date 0 --- Representative entrepreneur
borrows \(I - A\), invests \(I\). Date 1 --- Liquidity need \(\rho I\) realized; then ``Continue''
(→ Moral hazard) or ``Liquidate''. Date 2 --- Expected income \(\rho_1 I\), of which \(\rho_0 I\) is
pledgeable.

Let us assume that

\[\rho_L < \rho_0 < \rho_H < \rho_1 \qquad (15.1)\]

and

\[(1 - \lambda)(\rho_H - \rho_L) < 1. \qquad (15.2)\]

{[}\ldots{]} the unit cost of effective investment is
\(c(x) \equiv \dfrac{1 + (1 - \lambda)\rho_L + \lambda\rho_H x}{1 - \lambda + \lambda x}\)
{[}\ldots{]} where \(\bar\rho \equiv (1 - \lambda)\rho_L + \lambda\rho_H\) is the expected liquidity
shock per unit of investment. Lastly, we assume that \(\rho_1 > 1 + \bar\rho\), and so entrepreneurs
prefer investing to consuming \(A\) (the project's NPV is positive). To sum up, the efficient
allocation is given by \(x = 1\) and \(I = \dfrac{A}{(1 + \bar\rho) - \rho_0}\).

\emph{From Tirole §15.3.1 ``Aggregate Shocks and the Value of Outside Liquidity'', pp.~523--525 (PDF
536--538):}

Suppose now that liquidity shocks are perfectly correlated across firms in the model of Section 15.2
(that firms face the same shock is the starkest way of introducing an aggregate shock; Exercise 15.3
studies the more general case of imperfectly correlated shocks). In this polar case, and in the
absence of outside liquidity, firms cannot continue when \(\rho = \rho_H > \rho_0\): \(x = 0\).

We will assume that the entrepreneurs prefer investing to consuming even if the investment is
liquidated in the bad state of nature, i.e., that, per unit of investment, the expected output
\((1 - \lambda)\rho_1\) exceeds the sum of the initial investment cost, 1, and the expected
reinvestment cost, \((1 - \lambda)\rho_L\):

\[(1 - \lambda)\rho_1 > 1 + (1 - \lambda)\rho_L.\]

(This positive-NPV assumption is more stringent than the previous one: \(\rho_1 > 1 + \bar\rho\).)

{[}\ldots{]} Here we introduce outside stores of value into the picture. For simplicity, there are
\(L^S\) such stores of value. A store of value yields 1 unit of good at date 1 for certain. Its
date-0 price is \(q\); because the market rate of interest is equal to 0 and investors can buy the
store of value, \(q \geqslant 1\). Furthermore, if \(q = 1\), then the corporate sector is able to
implement its efficient allocation. If \(q > 1\), then the corporate sector must hold all stores of
value.

In the bad state of nature, the representative firm continues with probability \(x\). Equivalently,
a fraction \(1 - x\) of its investment is liquidated. For ease of exposition, we will work under the
latter interpretation.

The liquidity need in the bad state, \(\rho_H xI\), must not exceed the sum of the amount
\(\rho_0 xI\), which can be raised by returning to the capital market at date 1, and of the income
\(L\) associated with a date-0 purchase of outside liquidity in amount \(L\). And so

\[(\rho_H - \rho_0)xI \leqslant L.\]

The investors' breakeven constraint condition states that the total investor date-0 outlay to pay
for illiquid and liquid assets should be recouped from the pledgeable income:

\[[I + qL] - A \leqslant L + (1 - \lambda)(\rho_0 - \rho_L)I + \lambda(\rho_0 - \rho_H)xI.\]

To make things interesting, we will assume that \(L^S\) is not so large that \(q = 1\) {[}\ldots{]}.
\(q > 1\) only if firms compete with each other for the scarce liquidity. The date-0 price \(q\)
adjusts to the level at which the demand for liquid assets is equal to the supply: \(L = L^S\).

Let

\[U_b = [(1 - \lambda)(\rho_1 - \rho_L) + \lambda(\rho_1 - \rho_H)x]I - (q - 1)(\rho_H - \rho_0)xI - I \qquad (15.5)\]

denote the representative borrower's utility or net present value, given that an amount of liquidity
\(L = (\rho_H - \rho_0)xI\) is required to bring a fraction \(x\) of investment to completion in the
bad state of nature. The investors' breakeven constraint can be rewritten as

\[[(1 - \lambda)(\rho_0 - \rho_L) + \lambda(\rho_0 - \rho_H)x]I - (q - 1)(\rho_H - \rho_0)xI \geqslant I - A. \qquad (15.6)\]

This constraint yields the investment level \(I\) as a function of \(x\) and \(q\). Substituting
into (15.5),

\[U_b = \frac{\rho_1 - c(x, q)}{c(x, q) - \rho_0}A,\]

where

\[c(x, q) = \frac{1 + (1 - \lambda)\rho_L + \lambda\rho_H x + (q - 1)(\rho_H - \rho_0)x}{1 - \lambda + \lambda x}\]

is the cost of effective investment, that is, the average cost of bringing 1 unit of investment to
completion. Let \(\bar q\) be given by

\[\frac{\partial c}{\partial x}(x, \bar q) = 0 \iff (1 - \lambda)(\rho_H - \rho_L) + \frac{(1 - \lambda)}{\lambda}(\rho_H - \rho_0)(\bar q - 1) = 1.\]

Then \(x = 1\) if \(q < \bar q\), \(= 0\) if \(q > \bar q\), \(\in [0, 1]\) if \(q = \bar q\).

For \(q < \bar q\), the demand for liquid assets is given by

\[L^D = (\rho_H - \rho_0)I = \frac{\rho_H - \rho_0}{c(1, q) - \rho_0}A.\]

The equilibrium in the market for liquidity is depicted in Figure 15.3, where

\[\bar L^S \equiv \frac{(\rho_H - \rho_0)A}{1 + \bar\rho - \rho_0}\]

is the lowest amount of outside liquidity such that the market clearing price for liquidity is
\(q = 1\).

Figure 15.3 (p.~524, PDF 537): price \(q\) (vertical axis) against quantity \(L^D, L^S\) (horizontal
axis). The demand \(L^D\) is horizontal at \(q = \bar q\) for small quantities, then decreasing in
\(q\) down to \(q = 1\), which it reaches at \(\bar L^S\); three vertical supply lines
\(L_1^S < L_2^S < L_3^S\) cut it on the flat part, on the decreasing part, and on the floor
\(q = 1\), respectively.

Figure 15.3 illustrates the three regions. When liquidity is very scarce (\(L^S = L_1^S\) for
example), there is liquidation (\(x < 1\)); the price of liquid assets is then the highest price,
\(\bar q\), that firms are willing to pay in order to hold this liquidity and be able to salvage
their assets in the bad state of nature. As the supply of liquidity \(L^S\) increases, \(x\)
increases. Once \(x\) reaches 1, the price of liquid assets adjusts. A lower price increases the
firms' borrowing capacity and, indirectly, the demand for liquid assets, which are a complementary
input into the production process. {[}\ldots{]} Finally, for \(L^S \geqslant \bar L^S\)
{[}\ldots{]}, there is an excess supply of liquidity. The market price is then \(q = 1\), and the
excess liquidity, \(L^S - \bar L^S\), is absorbed by individual investors, who demand a rate of
return equal to 0.

\emph{Does liquidity crowd investment in or out?} {[}\ldots{]} Whether liquid assets crowd in or out
illiquid ones depends on whether the entrepreneurs are willing to invest without continuing in the
adverse state of nature (\((1 - \lambda)\rho_1 > 1 + (1 - \lambda)\rho_L\)) or not. In the former
case, which we have just analyzed, the purchase of liquid assets first comes at the expense of
investment: liquid assets crowd out illiquid ones. For larger amounts of liquidity
(\(L \geqslant \bar L^S\)), \(x = 1\) and liquid assets and investment are necessarily complements:
a larger stock of liquid assets lowers the liquidity premium \((q - 1)\) and thereby the cost of
investment.

When investing yields a negative NPV per unit of investment when the investment is liquidated in the
adverse state (\((1 - \lambda)\rho_1 < 1 + (1 - \lambda)\rho_L\)), entrepreneurs invest only if they
can complement their illiquid investment with liquidity. In this case, liquid and illiquid
investments are always complements (see Exercise 15.4).

\emph{From Tirole footnote 18, p.~525 (PDF 538):} For \(L^S < \underline L^S\), \(x < 1\) and
\(q = \bar q\). Investment is then given by
\([1 + (1 - \lambda)\rho_L - (1 - \lambda)\rho_0]I + [(\bar q - 1) + \lambda]L^S = A\). For
\(L^S \in [\underline L^S, \bar L^S]\), \(x = 1\) and \(q \leqslant \bar q\). Furthermore,
\(L^S = (\rho_H - \rho_0)I\) and
\(I = \dfrac{A}{1 + \bar\rho + (q - 1)(\rho_H - \rho_0) - \rho_0}\).
\stmtend
%% END-STATEMENT

\asked Redo the liquidity-pricing analysis of \S15.3.1 when investing without the option to continue in the bad state is unprofitable, and show that more outside liquidity $L^S$ raises investment $I$.

\approach Liquidity asset pricing (LAPM): under an aggregate shock, firms buy stores of value at price $q$ to continue in the bad state. The unit cost $c(x,q)$ drives everything. Here $x=0$ is not viable, so firms either invest with $x=1$ and buy liquidity, or do not invest: liquidity becomes a necessary complementary input.

\notation
\begin{itemize}
  \item As in exercise 3: $\rho_L<\rho_0<\rho_H<\rho_1$, $\lambda$, $\bar\rho$; $q\geq1$: price of a store yielding 1 at date 1; $L^S$: supply.
  \item $\bar{\bar q}$: price at which entrepreneurs are indifferent about investing; $\bar I=A/(\rho_1-\rho_0)$; $\underline L^S$, $\bar L^S$: region boundaries.
\end{itemize}

\step{1}{NPV per unit} From (15.5), per unit of investment,
\begin{align}
  n(x,q) = (1-\lambda+\lambda x)\rho_1-\bigl[1+(1-\lambda)\rho_L+[\lambda\rho_H+(q-1)(\rho_H-\rho_0)]x\bigr]. \label{r10e13-npv}
\end{align}
At $x=0$, $n=(1-\lambda)\rho_1-1-(1-\lambda)\rho_L<0$ by assumption.

\step{2}{$x=1$} $n$ is linear in $x$ with slope $\partial n/\partial x=\lambda\rho_1-\lambda\rho_H-(q-1)(\rho_H-\rho_0)$. If the slope were $\leq0$, $n<0$ for all $x$: no investment, no demand for liquidity, so $q=1$; but then the slope is $\lambda(\rho_1-\rho_H)>0$, a contradiction. Hence in equilibrium
\begin{align}
  \lambda\rho_1>\lambda\rho_H+(q-1)(\rho_H-\rho_0) \quad\text{and}\quad x=1. \label{r10e13-x1}
\end{align}

\step{3}{Investment and demand for liquidity at $x=1$} With $x=1$, $c(1,q)=1+\bar\rho+(q-1)(\rho_H-\rho_0)$ and, from (15.6) binding,
\begin{align}
  I = \frac{A}{c(1,q)-\rho_0} = \frac{A}{1+\bar\rho+(q-1)(\rho_H-\rho_0)-\rho_0}, \qquad L^D=(\rho_H-\rho_0)I, \label{r10e13-I}
\end{align}
when $c(1,q)<\rho_1$ (strictly positive NPV, so firms borrow to capacity). When $c(1,q)=\rho_1$, entrepreneurs are indifferent and any $I\in[0,\bar I]$ is optimal, where $\bar I$ solves $\rho_0\bar I=\rho_1\bar I-A$:
\begin{align}
  1+\bar\rho+(\bar{\bar q}-1)(\rho_H-\rho_0)=\rho_1 \iff \bar{\bar q}=1+\frac{\rho_1-1-\bar\rho}{\rho_H-\rho_0}, \qquad \bar I=\frac{A}{\rho_1-\rho_0}. \label{r10e13-qq}
\end{align}
At $\bar{\bar q}$, \eqref{r10e13-x1} holds: $\lambda(\rho_1-\rho_H)-(\rho_1-1-\bar\rho)=1+(1-\lambda)\rho_L-(1-\lambda)\rho_1>0$.

\step{4}{Market clearing, three regions} Market clearing $L^D=L^S$ (with $q>1$) gives:
\begin{align}
  L^S\leq\underline L^S\equiv\frac{(\rho_H-\rho_0)A}{\rho_1-\rho_0}&: \quad q=\bar{\bar q},\quad I=\frac{L^S}{\rho_H-\rho_0}, \label{r10e13-r1}\\
  \underline L^S\leq L^S\leq\bar L^S\equiv\frac{(\rho_H-\rho_0)A}{1+\bar\rho-\rho_0}&: \quad I=\frac{L^S}{\rho_H-\rho_0},\quad q \text{ from } \frac{(\rho_H-\rho_0)A}{1+\bar\rho+(q-1)(\rho_H-\rho_0)-\rho_0}=L^S, \label{r10e13-r2}\\
  L^S\geq\bar L^S&: \quad q=1,\quad I=\frac{A}{1+\bar\rho-\rho_0}. \label{r10e13-r3}
\end{align}
In the first region liquidity limits investment and entrepreneurs, indifferent at $\bar{\bar q}$, invest less than their capacity. In the second, $q$ falls as $L^S$ rises, raising borrowing capacity. In the third, liquidity is in excess and $q=1$.

\step{5}{Monotonicity} $I$ is continuous in $L^S$ (at $\underline L^S$, $I=\bar I$ from both sides; at $\bar L^S$, $I=A/(1+\bar\rho-\rho_0)$), strictly increasing up to $\bar L^S$ and constant afterwards.

\step{6}{Example} $\lambda=0.5$, $\rho_L=0.2$, $\rho_0=0.6$, $\rho_H=1$, $\rho_1=2$, $A=1$: $(1-\lambda)\rho_1=1<1.1$, $\rho_1>1+\bar\rho=1.6$. $\bar{\bar q}=1+0.4/0.4=2$, $\underline L^S=0.4/1.4=0.285714$, $\bar L^S=0.4/1=0.4$, $\bar I=0.714286$. $L^S=0.2$: $q=2$, $I=0.5$. $L^S=0.35$: $I=0.875$, $q=1.357143$. $L^S\geq0.4$: $q=1$, $I=1$.

\answer{} $x=1$ in any equilibrium with investment; $I=L^S/(\rho_H-\rho_0)$ for $L^S\leq\bar L^S$ (at $q=\bar{\bar q}$ up to $\underline L^S$, then with $q$ falling) and $I=A/(1+\bar\rho-\rho_0)$ beyond: investment is nondecreasing, strictly increasing until $\bar L^S$ --- liquid and illiquid assets are complements.

\intuition When a project is worthless unless it can be carried through the bad state, a store of value is a required input, like a spare part. More stores directly allow more projects, then make the remaining ones cheaper through a lower liquidity premium. The crowding-out region of \S15.3.1 disappears because firms never trade continuation for scale.

\pitfalls
\begin{itemize}
  \item Allowing $x<1$: with $n(0,q)<0$ and $n$ linear, any investing firm sets $x=1$.
  \item Assuming firms always borrow to capacity: at $q=\bar{\bar q}$ they are indifferent and investment is pinned by $L^S$.
  \item Confusing $\bar q$ of \S15.3.1 (indifference about $x$) with $\bar{\bar q}$ (indifference about investing).
\end{itemize}

\bookans Matches the book's answer (p.~621): $x=1$ by the same contradiction, $L^S=(\rho_H-\rho_0)I$, $\rho_1=1+\bar\rho+(\bar{\bar q}-1)(\rho_H-\rho_0)$, $\underline L^S=(\rho_H-\rho_0)A/(\rho_1-\rho_0)$, $q$ falling and $I=L^S/(\rho_H-\rho_0)$ rising until $\bar L^S$.
\refbib{Tirole \S15.3.1, pp.~523--525; \S15.3.2 (liquidity asset pricing).}

% =============================================================================
\exsec{14. Dynamics of income inequality --- Tirole 13.2, pp.~493--494}
%% BEGIN-STATEMENT 14
\stmtbegin
\textbf{Source:} Tirole, Ch. 13, Exercise 13.2, pp.~493--494 (PDF 506) · Answer in book: yes, p.~617

\textbf{Exercise 13.2 (dynamics of income inequality).} (This exercise builds on the analysis of
Section 13.4 and on Matsuyama (2000).)

\begin{enumerate}
\def\labelenumi{(\roman{enumi})}
\tightlist
\item
  Consider the ``warm-glow'' model: generations are indexed by \(t = 0, 1, \ldots, \infty\). Each
  generation lives for one period; each individual has exactly one heir. A generation-\(t\)
  individual has utility from consumption \(c_t\) and bequest \(\mathcal{L}_t\) equal to
\end{enumerate}

\[\left(\frac{c_t}{1 - a}\right)^{1-a}\left(\frac{\mathcal{L}_t}{a}\right)^a\]

with \(0 < a < 1\).

What is the individual's utility from income \(y_t\)?

\begin{enumerate}
\def\labelenumi{(\roman{enumi})}
\setcounter{enumi}{1}
\tightlist
\item
  Consider the entrepreneurship model of Section 13.4, with two twists:
\end{enumerate}

• variable-size investment (instead of a fixed-size one),

• intraperiod rate of interest \(r\) (so investors demand \((1 + r)\) times their outlay, in
expectation); \(r\) is assumed constant for simplicity.

One will assume that \(p_H = 1\) and that each generation \(t\) is born with endowment \(\hat A\)
(to which is added bequest \(\mathcal{L}_{t-1}\), so \(A_t = \hat A + \mathcal{L}_{t-1}\)). See
Figure 13.12.

Figure 13.12 (timeline of generation \(t\), left to right; p.~494, PDF 507)

{\def\LTcaptype{none} % do not increment counter
\begin{longtable}[]{@{}
  >{\raggedright\arraybackslash}p{(\linewidth - 6\tabcolsep) * \real{0.2500}}
  >{\raggedright\arraybackslash}p{(\linewidth - 6\tabcolsep) * \real{0.2500}}
  >{\raggedright\arraybackslash}p{(\linewidth - 6\tabcolsep) * \real{0.2500}}
  >{\raggedright\arraybackslash}p{(\linewidth - 6\tabcolsep) * \real{0.2500}}@{}}
\toprule\noalign{}
\begin{minipage}[b]{\linewidth}\raggedright
Start of period
\end{minipage} & \begin{minipage}[b]{\linewidth}\raggedright
Moral hazard
\end{minipage} & \begin{minipage}[b]{\linewidth}\raggedright
Outcome
\end{minipage} & \begin{minipage}[b]{\linewidth}\raggedright
End of period
\end{minipage} \\
\midrule\noalign{}
\endhead
\bottomrule\noalign{}
\endlastfoot
\(\mathcal{L}_{t-1}\) received. Asset \(\hat A + \mathcal{L}_{t-1}\). Investment in project of
(endogenous) size \(I\). & Behaves: Pr(success) 1, private benefit 0. Misbehaves: Pr(success)
\(p_L\), private benefit \(BI\). & Outcome: success/failure. Income \(y_t\). & Consumption. Bequest
(to generation \(t + 1\)). \\
\end{longtable}
}

A successful project delivers \(RI \geqslant (1 + r)I\), an unsuccessful one 0. The private benefit
from misbehaving, \(BI\), is also proportional to investment.

Let

\[\rho_1 \equiv R \quad\text{and}\quad \rho_0 \equiv R - \frac{B}{\Delta p}.\]

Assume that

\[a(\rho_1 - \rho_0) < 1 - \frac{\rho_0}{1 + r}.\]

Show that each dynasty's long-term wealth converges to

\[A_\infty \equiv \frac{\hat A}{1 - a(\rho_1 - \rho_0)/(1 - \rho_0/(1 + r))},\]

regardless of its initial total wealth \(A_0\) (that is, \(\hat A\) plus the bequest from generation
\(-1\), if any).

\begin{enumerate}
\def\labelenumi{(\roman{enumi})}
\setcounter{enumi}{2}
\tightlist
\item
  Now assume that there is a minimal investment scale \(\underline I > 0\) below which nothing can
  be produced. For \(I \geqslant \underline I\), the technology is as above (constant returns to
  scale, profit \(RI\) in the case of success, private benefit \(BI\) in the case of misbehavior,
  etc.).
\end{enumerate}

Compute the threshold \(A_0^*\) under which the dynasty remains one of lenders (at rate \(r\)) and
never makes it to entrepreneurship.

What is the limit wealth \(A_\infty^L\) of these poor dynasties? (The limit wealth of dynasties
starting with \(A_0 \geqslant A_0^*\) is still \(A_\infty\).)

\begin{enumerate}
\def\labelenumi{(\roman{enumi})}
\setcounter{enumi}{3}
\tightlist
\item
  Finally, close the model by assuming that investors are domestic investors and by describing the
  equilibrium in the loan market. Focus on \emph{steady states}. Show that multiple steady states
  may coexist:
\end{enumerate}

• one in which everyone (investors, entrepreneurs) has the same wealth and

\[\rho_1 = 1 + r,\]

• others, with unequal wealth distribution, in which \(\rho_1 > 1 + r\), a fraction \(\kappa\) of
the population is poor (lends), and a fraction \(1 - \kappa\) is rich (borrows to undertake
projects).

\paragraph{Context needed to solve it}

\emph{From Tirole §13.4.2.1 ``The Warm-Glow Model'', pp.~484--485 (PDF 497--498):}

We study long-lived lineages of short-lived family members. Parents become entrepreneurs if they
have sufficient funds, earn income, and finally leave wealth as a bequest to their children, who may
then use this wealth to undertake projects of their own, and so forth. {[}\ldots{]} For the purposes
of this section, we first bypass the difficulty in a somewhat ad hoc way by using the warm-glow
model, which enables us to discuss dynamics without worrying about dynamic programming. Namely,
suppose that the following conditions hold.

• Individuals live for one period. An individual living at date \(t\) has exactly one heir who lives
at date \(t + 1\), and so on.

• Individuals are ``altruistic'' in a rather specific way. Rather than caring about the utility of
their heirs, they derive utility from the bequest they make to their heirs. We will assume that a
generation-\(t\) individual derives utility from her own consumption \(c_t\) and from the bequest
\(\mathcal{L}_t\) to her heir. Assume further that the utility function is a Cobb--Douglas utility
function:

\[\left(\frac{c_t}{1 - a}\right)^{1-a}\left(\frac{\mathcal{L}_t}{a}\right)^a,\]

where \(a \in (0, 1)\) is the (impure) altruism parameter.

Then the utility from income \(y_t\) is (taking logs)

\[\log U_t(y_t) = \max_{\{c_t, \mathcal{L}_t\}}\{(1 - a)\log c_t + a\log\mathcal{L}_t\} \quad\text{s.t.}\quad c_t + \mathcal{L}_t = y_t.\]

\emph{From Tirole §13.4.2.2 ``Lineages of Entrepreneurs in the Warm-Glow Model'', pp.~485--487 (PDF
498--500):}

Let us now consider a ``warm-glow lineage'' in which each generation \(t\) is a would-be
entrepreneur, who

• is born with some exogenous endowment \(\hat A\), to which is added the bequest
\(\mathcal{L}_{t-1}\) made by generation \(t - 1\);

• invests this total asset either in a storage technology yielding an interest rate equal to 0
(i.e., preserving the wealth) or in a fixed-size project as described in Section 13.2.1; and

• finally uses the proceed of this investment (her ``income'') for consumption \(c_t\) and bequest
\(\mathcal{L}_t\) to the next generation.

The timing is summarized in Figure 13.7.

Figure 13.7 (p.~486, PDF 499): same layout as Figure 13.12, with: Asset
\(\hat A + \mathcal{L}_{t-1}\); Storage (return: 0) \emph{or} investment in project of size \(I\) →
Moral hazard (if investment): Behaves --- Pr(success) \(p_H\), private benefit 0; Misbehaves ---
Pr(success) \(p_L\), private benefit \(B\) → Outcome: success/failure (if investment); Income
\(y_t\) → Consumption \(c_t = (1 - a)y_t\) → Bequest \(\mathcal{L}_t = ay_t\).

Let us assume that the intraperiod rate of interest in the economy is equal to 0.

As in the rest of the book, the private benefit \(B\) obtained by misbehaving is expressed in terms
of money. So in the case of misbehavior the entrepreneur's utility from income \(y_t\) and private
benefit \(B\) is \(y_t + B\). As usual, we will assume that investment has a positive NPV if and
only if the entrepreneur is induced to behave: \(p_H R > I > p_L R + B\). It will also prove
convenient to assume that success is a sure thing in the case of good behavior: \(p_H = 1\).

Lastly, let \(\bar A\) be defined (as in equation (13.4)) by the equality between the pledgeable
income and the investors' outlay:

\[p_H\left(R - \frac{B}{\Delta p}\right) = I - \bar A \quad\text{or (using } p_H = 1\text{)}\quad \bar A = I - \left(R - \frac{B}{\Delta p}\right).\]

\textbf{Assumption 13.1.} \(\dfrac{\hat A + a(R - I)}{1 - a} > \bar A > \dfrac{\hat A}{1 - a}\).

\begin{enumerate}
\def\labelenumi{(\alph{enumi})}
\tightlist
\item
  \emph{Lineage stuck in a poverty trap.} Suppose that generation \(t\) receives bequest
  \(\mathcal{L}_{t-1} < \bar A - \hat A\). Generation \(t\)'s total wealth is then insufficient to
  have access to funds. Generation \(t\) must therefore invest \([\hat A + \mathcal{L}_{t-1}]\) into
  the low-return storage technology, and so \(y_t = \hat A + \mathcal{L}_{t-1}\). With warm-glow
  preferences, bequests to generation \(t + 1\) are
  \(\mathcal{L}_t = ay_t = a(\hat A + \mathcal{L}_{t-1})\), {[}\ldots{]} The dynasty's total wealth
  per generation converges to \(A_\infty < \bar A\), given by
\end{enumerate}

\[A_\infty = \hat A + aA_\infty \quad\text{or}\quad A_\infty = \frac{\hat A}{1 - a} < \bar A.\]

The lineage is stuck in a poverty trap.

\begin{enumerate}
\def\labelenumi{(\alph{enumi})}
\setcounter{enumi}{1}
\tightlist
\item
  \emph{Rich, entrepreneurial lineage.} By contrast, suppose that generation \(t\)'s initial wealth
  exceeds \(\bar A\): \(\mathcal{L}_{t-1} > \bar A - \hat A\) or
  \(A_t \equiv \hat A + \mathcal{L}_{t-1} > \bar A\). Generation \(t\) has enough pledgeable income
  to offset the inventors' outlay \(I - (\hat A + \mathcal{L}_{t-1})\). Under risk neutrality,
  generation \(t\) selects the highest NPV solution and therefore prefers becoming an entrepreneur
  to investing in the storage technology. The NPV is then \(p_H R - I = R - I > 0\), and so the
  entrepreneur's end-of-period income after reimbursing lenders is \(y_t = (R - I) + A_t\) (recall
  that the capital market is competitive, and so the entire NPV goes to the entrepreneur).
  Generation \((t + 1)\)'s total wealth at the beginning of period \(t + 1\) is therefore
  \(A_{t+1} = \hat A + a(R - I + A_t)\). {[}\ldots{]} The lineage's beginning-of-period wealth
  converges to \(A_\infty\), where
\end{enumerate}

\[A_\infty = \hat A + a(R - I + A_\infty) \quad\text{or}\quad A_\infty = \frac{\hat A + a(R - I)}{1 - a}.\]

Figure 13.8 (p.~487, PDF 500): next generation's wealth \(A_{t+1}\) against current wealth \(A_t\).
Two increasing lines with slope \(a\): \(A_{t+1} = \hat A + aA_t\) for \(A_t < \bar A\) (fixed point
\(\hat A/(1 - a)\)) and \(A_{t+1} = \hat A + a(R - I + A_t)\) for \(A_t > \bar A\) (fixed point
\([\hat A + a(R - I)]/(1 - a)\)); two initial wealths \(A^-\) just below \(\bar A\) and \(A^+\) just
above it converge to the low and the high fixed point, respectively.

Figure 13.8 illustrates cases (a) and (b) and shows that a small difference in initial wealth
(points \(A^-\) and \(A^+\), respectively, in Figure 13.8) can make a big difference: for the
current generation (as we know from Chapter 3) and even more for further generations.

Note, finally, that in the Arrow--Debreu world of no agency cost (\(B = 0\)), long-term incomes of
different lineages would converge to \([\hat A + a(R - I)]/(1 - a)\) regardless of the lineage's
initial wealth (rather than diverge as in Figure 13.8). A stronger investor protection, for example,
reduces the dependency on wealth and generates a more equal long-run income distribution.
\stmtend
%% END-STATEMENT

\asked (i) The indirect utility of income in the warm-glow model. (ii) Wealth dynamics of a dynasty of entrepreneurs with variable investment at rate $r$, and convergence to $A_\infty$. (iii) With a minimum scale $\underline I$: the wealth below which a dynasty stays a lender forever, and its limit wealth. (iv) Closing the loan market: show that an equal-wealth steady state with $\rho_1=1+r$ and unequal steady states with $\rho_1>1+r$ can coexist.

\approach Warm-glow bequests make utility linear in income (Tirole \S13.4.2.1), so each generation maximises income. Borrowing capacity is the equity multiplier of \S3.4 at rate $r$; income is the entrepreneur's rent $(\rho_1-\rho_0)I$. A minimum scale creates a poverty trap; general equilibrium in the loan market endogenises $r$ (Matsuyama 2000).

\notation
\begin{itemize}
  \item $a\in(0,1)$: bequest share; $c_t$, $\mathcal L_t$, $y_t$: consumption, bequest, income of generation $t$; $\hat A$: endowment; $A_t=\hat A+\mathcal L_{t-1}$.
  \item $\rho_1=R$, $\rho_0=R-B/\Delta p$ ($p_H=1$); $r$: interest rate; $\beta\equiv a(\rho_1-\rho_0)/(1-\rho_0/(1+r))$.
  \item $\underline I$: minimum scale; $A_0^*$: entry threshold; $\kappa$: fraction of lenders.
\end{itemize}

\step{1}{Warm glow} Maximise $(1-a)\log c+a\log\mathcal L$ s.t.\ $c+\mathcal L=y$. The first-order conditions $(1-a)/c=a/\mathcal L$ and the budget give $c=(1-a)y$, $\mathcal L=ay$, so
\begin{align}
  U = \left(\frac{(1-a)y}{1-a}\right)^{1-a}\left(\frac{ay}{a}\right)^a = y^{1-a}y^a = y. \label{r10e14-U}
\end{align}

\answer{(i)} $U_t(y_t)=y_t$.

\step{2}{Borrowing capacity} With $p_H=1$, the incentive constraint $(\Delta p)R_b\geq BI$ leaves investors $(R-B/\Delta p)I=\rho_0I$; they require $(1+r)(I-A_t)$:
\begin{align}
  \rho_0I_t = (1+r)(I_t-A_t) \;\Longrightarrow\; I_t = \frac{A_t}{1-\rho_0/(1+r)}. \label{r10e14-I}
\end{align}
The NPV $[\rho_1-(1+r)]I_t\geq0$, so investing to capacity beats lending.

\step{3}{Income and bequest} Income after repaying investors, using \eqref{r10e14-I}:
\begin{align}
  y_t = \rho_1I_t-\rho_0I_t = (\rho_1-\rho_0)I_t, \qquad A_{t+1} = \hat A+ay_t = \hat A+\beta A_t. \label{r10e14-dyn}
\end{align}
The assumption $a(\rho_1-\rho_0)<1-\rho_0/(1+r)$ is $\beta<1$, so \eqref{r10e14-dyn} is a contraction: $A_t\to\hat A/(1-\beta)=A_\infty$ from any $A_0$.

\answer{(ii)} $A_{t+1}=\hat A+\beta A_t$ with $\beta<1$, so $A_t\to A_\infty=\hat A/[1-a(\rho_1-\rho_0)/(1-\rho_0/(1+r))]$ for every $A_0$.

\step{4}{Minimum scale} Entrepreneurship needs $I_t\geq\underline I$, i.e.\ by \eqref{r10e14-I},
\begin{align}
  A_t\geq A_0^*\equiv\left(1-\frac{\rho_0}{1+r}\right)\underline I. \label{r10e14-A0}
\end{align}
A lender earns $y_t=(1+r)A_t$, so $A_{t+1}=\hat A+a(1+r)A_t$, which converges (if $a(1+r)<1$) to
\begin{align}
  A_\infty^L = \frac{\hat A}{1-a(1+r)}. \label{r10e14-AL}
\end{align}
If $A_\infty^L<A_0^*$, the lenders' map lies below the 45$^\circ$ line at $A_0^*$, so a dynasty starting below $A_0^*$ stays below forever: $A_0^*$ is the threshold. (If $A_\infty^L\geq A_0^*$, every dynasty eventually becomes entrepreneurial.)

\answer{(iii)} $A_0^*=(1-\rho_0/(1+r))\underline I$; poor dynasties converge to $A_\infty^L=\hat A/[1-a(1+r)]$ (a trap when $A_\infty^L<A_0^*\leq A_\infty$).

\step{5}{Equal-wealth steady state} If $\rho_1=1+r$, lending and investing pay the same, so everyone has the same wealth. Then $\beta=a\rho_1(\rho_1-\rho_0)/(\rho_1-\rho_0)=a\rho_1$ and
\begin{align}
  A_\infty = \frac{\hat A}{1-a\rho_1}, \qquad I_\infty=\frac{A_\infty}{1-\rho_0/\rho_1}=\frac{\rho_1\hat A}{(1-a\rho_1)(\rho_1-\rho_0)}. \label{r10e14-eq}
\end{align}
The fraction $\kappa$ of lenders clears the loan market:
\begin{align}
  \kappa A_\infty &= (1-\kappa)(I_\infty-A_\infty) \nonumber\\
  \kappa &= (1-\kappa)\left(\frac{\rho_1}{\rho_1-\rho_0}-1\right) = (1-\kappa)\frac{\rho_0}{\rho_1-\rho_0} \nonumber\\
  \kappa(\rho_1-\rho_0) &= (1-\kappa)\rho_0 \iff \kappa=\frac{\rho_0}{\rho_1}. \label{r10e14-keq}
\end{align}
It is a steady state if $I_\infty\geq\underline I$.

\step{6}{Unequal steady states} If $\rho_1>1+r$ everyone wants to invest, so lenders must be unable to: their wealth is $A_\infty^L<A_0^*$, while the rich have $A_\infty\geq A_0^*$. Loan-market clearing:
\begin{align}
  \kappa A_\infty^L = (1-\kappa)(I_\infty-A_\infty), \qquad I_\infty=\frac{A_\infty}{1-\rho_0/(1+r)} \;\Longrightarrow\; \kappa=\frac{I_\infty-A_\infty}{A_\infty^L+I_\infty-A_\infty}. \label{r10e14-kun}
\end{align}
Each $r$ satisfying $A_\infty^L(r)<A_0^*(r)\leq A_\infty(r)$ and $\beta(r)<1$ supports such a steady state, with $\kappa$ from \eqref{r10e14-kun}.

\step{7}{Example of coexistence} $a=0.4$, $\hat A=1$, $\rho_1=1.5$, $\rho_0=0.5$, $\underline I=3.5$. Equal: $r=0.5$, $A_\infty=1/(1-0.6)=2.5$, $I_\infty=1.5/(0.4\cdot1)=3.75\geq3.5$, $A_0^*=(1-1/3)(3.5)=2.333333\leq2.5$, $\kappa=1/3$. Unequal at $r=0.05$: $\beta=0.4/(1-0.5/1.05)=0.763636$, $A_\infty=1/(1-0.763636)=4.230769$, $I_\infty=8.076923$, $A_0^*=1.833333$, $A_\infty^L=1/(1-0.42)=1.724138<A_0^*$, $\kappa=3.846154/(1.724138+3.846154)=0.690476$. Both are steady states for the same primitives.

\answer{(iv)} Equal steady state: $\rho_1=1+r$, common wealth $\hat A/(1-a\rho_1)$, $\kappa=\rho_0/\rho_1$. Unequal ones: $\rho_1>1+r$, poor wealth $A_\infty^L<A_0^*$, rich wealth $A_\infty$, $\kappa$ from \eqref{r10e14-kun}; several can coexist (example).

\intuition Without a minimum scale, credit constraints slow but do not stop convergence: every dynasty invests proportionally to its wealth. With a minimum scale, being poor means lending at $r$ rather than earning the entrepreneurial rent, and low wealth perpetuates itself. Closing the model makes inequality self-sustaining: many poor lenders keep $r$ low, which is what keeps them from saving enough to enter.

\pitfalls
\begin{itemize}
  \item Using $y_t=(1+r)A_t+$NPV and then forgetting that, with binding borrowing, it equals $(\rho_1-\rho_0)I_t$ --- both are right, mixing them double counts.
  \item Writing the poor dynasty's dynamics with storage at 0 instead of lending at $r$.
  \item Forgetting $a(1+r)<1$ for convergence of poor dynasties.
\end{itemize}

\bookans Matches the book's answer (p.~617): $U_t=y_t$; $I=A_t/(1-\rho_0/(1+r))$, $y_t=(\rho_1-\rho_0)I_t$, convergence to $A_\infty$; $A_0^*/(1-\rho_0/(1+r))=\underline I$, $A_\infty^L=\hat A/[1-a(1+r)]$; the equal-wealth steady state with $\kappa(\rho_1-\rho_0)=(1-\kappa)\rho_0$ and the unequal one with $\kappa A_\infty^L=(1-\kappa)(I_\infty-A_\infty)$. The example of coexistence is an addition.
\refbib{Tirole \S13.4.2, pp.~484--487; \S13.4.3; Matsuyama (2000).}

% =============================================================================
\exsec{15. Political economy of exchange-rate policies --- Tirole 16.3, pp.~565--566}
%% BEGIN-STATEMENT 15
\stmtbegin
\textbf{Source:} Tirole, Ch. 16, Exercise 16.3, pp.~565--566 (PDF 578) · Answer in book: yes, p.~623

\textbf{Exercise 16.3 (political economy of exchange rate policies).} Consider a country that has
liberalized its capital account. There are two goods: a tradable good (the only one consumed by
foreigners) and a nontradable good.

• The only investors are foreign investors, with preferences over date-0 and date-1 consumptions

\[c_0^* + c_1^*,\]

where an asterisk refers to the tradable good.

• The country is populated by a unit mass of domestic entrepreneurs endowed with a
constant-returns-to-scale technology. The representative entrepreneur (1) invests \(I\) units of
tradables in equipment (where \(I\) is endogenous), (2) produces \(RI\) units of tradables in the
case of success, and 0 in the case of failure, and \(SI\) units of nontradables for certain. We
assume that firms' outcomes are independent (there is no macroeconomic shock).

The model is a variation on the standard variable-investment model:

• Each entrepreneur is initially endowed with \(A\) units of tradables (her only wealth), borrows
\(I - A\).

• There is moral hazard. The probability of success in the tradable-good activity is \(p_H\) if the
entrepreneur behaves, and \(p_L\) otherwise. The entrepreneur receives private benefit \(BI\) in
\emph{tradables} by misbehaving and 0 otherwise.

• An entrepreneur's utility is \(c_1 + u(c_1^*) + v(g^*)\), where \(c_1\) is the consumption of
nontradables, \(c_1^*\) the consumption of the tradable good, and \(g^*\) the level of public good
supplied by the government. \(u\) and \(v\) are concave.

We add a government. The government has international reserves \(\mathcal{R}^*\), of which it
consumes \(g^*\) to produce a public good. The rest, \(\mathcal{R}^* - g^*\), is dumped on the
currency market at the end. So, \(e\), the price of tradables in terms of nontradables, is given
\emph{ex post} by

\[p_H RI + \mathcal{R}^* - g^* = c_1^*(e) + d^* + \frac{d}{e},\]

where \(I\) is the representative entrepreneur's investment, \(d^*\) is the entrepreneurs' average
reimbursed debt in tradables and \(d\) is the average reimbursed debt in nontradables. The
government cares only about the welfare of entrepreneurs, i.e., does not internalize that of the
foreigners.

The timing is summarized in Table 16.1, where ``t'' and ``nt'' stand for ``tradables'' and
``nontradables,'' respectively.

Table 16.1 (p.~565, PDF 578)

{\def\LTcaptype{none} % do not increment counter
\begin{longtable}[]{@{}
  >{\raggedright\arraybackslash}p{(\linewidth - 8\tabcolsep) * \real{0.2000}}
  >{\raggedright\arraybackslash}p{(\linewidth - 8\tabcolsep) * \real{0.2000}}
  >{\raggedright\arraybackslash}p{(\linewidth - 8\tabcolsep) * \real{0.2000}}
  >{\raggedright\arraybackslash}p{(\linewidth - 8\tabcolsep) * \real{0.2000}}
  >{\raggedright\arraybackslash}p{(\linewidth - 8\tabcolsep) * \real{0.2000}}@{}}
\toprule\noalign{}
\begin{minipage}[b]{\linewidth}\raggedright
\end{minipage} & \begin{minipage}[b]{\linewidth}\raggedright
Date 0
\end{minipage} & \begin{minipage}[b]{\linewidth}\raggedright
Date 1
\end{minipage} & \begin{minipage}[b]{\linewidth}\raggedright
\end{minipage} & \begin{minipage}[b]{\linewidth}\raggedright
\end{minipage} \\
\midrule\noalign{}
\endhead
\bottomrule\noalign{}
\endlastfoot
(government) & Government selects \(g^*\) (commitment) & Government selects \(g^*\) (noncommitment)
& & \\
nt & & & \(SI\) & Exchange rate determined \\
t & Entrepreneur borrows \(I - A\) & Entrepreneur takes private benefit \(BI\), or none & Outcome
\(RI\) or 0 & Exchange rate determined \\
\end{longtable}
}

Consider financing contracts in which

• investors receive \(R_l^* = RI - R_b^*\) in tradables in the case of success, and 0 in the case of
failure;

• investors have nominal claims \(R_l^S\) and \(R_l^F\) in nontradables in the cases of success and
failure, respectively.

\begin{enumerate}
\def\labelenumi{(\roman{enumi})}
\item
  Relate \((d, d^*)\) and \((R_l^*, R_l^S, R_l^F)\).
\item
  Fixing an expected exchange rate \(e\), determine the investment \(I\) of the representative
  entrepreneur in this constant-returns-to-scale model assuming that
  \(\rho_0 \equiv p_H(R - B/\Delta p) < 1 - (S/e)\) in the relevant range.
\end{enumerate}

Show that \(R_l^F = SI\) and that

\[I = \frac{A}{1 - [(S/e) + \rho_0]}.\]

\begin{enumerate}
\def\labelenumi{(\roman{enumi})}
\setcounter{enumi}{2}
\tightlist
\item
  Compare the exchange rate and the welfare of entrepreneurs when the government chooses \(g^*\)
  after the private sector borrows abroad (``noncommitment'') and when the government can commit to
  \(g^*\) before entrepreneurs borrow abroad (``commitment'').
\end{enumerate}

Assume that the exchange rate depreciates as government expenditures \(g^*\) grow. Show that
\(v'(g^*) > e\) under commitment (underspending) and \(v'(g^*) < e\) under noncommitment
(overspending).

\begin{enumerate}
\def\labelenumi{(\roman{enumi})}
\setcounter{enumi}{3}
\tightlist
\item
  Show that there is an externality among borrowers when the government cannot commit.
\end{enumerate}

\paragraph{Context needed to solve it}

None beyond Table 16.1, transcribed above (the variable-investment model is restated in the
statement).
\stmtend
%% END-STATEMENT

\asked (i) Express average debt in tradables $d^*$ and nontradables $d$ in terms of the contract. (ii) For a given expected exchange rate $e$, investment and the optimal nontradable claim in failure. (iii) Compare public spending, the exchange rate and welfare with and without commitment: underspending ($v'>e$) vs.\ overspending ($v'<e$). (iv) The externality among borrowers without commitment.

\approach Variable-investment model with two goods. Pledgeable income now includes the certain nontradable output $SI$, worth $SI/e$ to foreign investors. The government spends reserves; spending depreciates the currency ($e\uparrow$), which reduces the tradable value of nontradable debt --- a transfer from foreigners that a government without commitment exploits (Tirole \S16.3; Caballero--Krishnamurthy, discussed in \S15.2.2).

\notation
\begin{itemize}
  \item $e$: price of tradables in nontradables; $R_l^*$, $R_l^S$, $R_l^F$: investors' tradable claim if success and nontradable claims in success/failure; $R_b^*=RI-R_l^*$, $R_b^S=SI-R_l^S$, $R_b^F=SI-R_l^F$.
  \item $d^*$, $d$: average repaid debt in tradables / nontradables; $\Rc^*$: reserves; $g^*$: public good; $u$, $v$: concave utilities; $c_1^*(e)$: tradable consumption with $u'(c_1^*)=e$.
  \item $\rho_0=p_H(R-B/\Delta p)$.
\end{itemize}

\step{1}{Average debts} Outcomes are independent, so averages equal expectations:
\begin{align}
  d^* = p_HR_l^*, \qquad d = p_HR_l^S+(1-p_H)R_l^F. \label{r10e15-d}
\end{align}

\answer{(i)} $d^*=p_HR_l^*$, $d=p_HR_l^S+(1-p_H)R_l^F$.

\step{2}{Constraints in tradable units} Nontradables are worth $1/e$ tradables. Incentive constraint and foreign investors' breakeven:
\begin{align}
  (\Delta p)\left[R_b^*+\frac{R_b^S-R_b^F}{e}\right] &\geq BI, \label{r10e15-ic}\\
  d^*+\frac de = p_H(RI-R_b^*)+\frac{p_H(SI-R_b^S)+(1-p_H)(SI-R_b^F)}{e} &\geq I-A. \label{r10e15-be}
\end{align}
From \eqref{r10e15-ic}, $p_HR_b^*\geq p_HBI/\Delta p-p_H(R_b^S-R_b^F)/e$. Substituting into the middle of \eqref{r10e15-be}:
\begin{align}
  d^*+\frac de &\leq p_HRI-\frac{p_HBI}{\Delta p}+\frac{p_H(R_b^S-R_b^F)}{e}+\frac{SI-p_HR_b^S-(1-p_H)R_b^F}{e} \nonumber\\
  &= \rho_0I+\frac{SI-R_b^F}{e} = \rho_0I+\frac{R_l^F}{e}. \label{r10e15-pl}
\end{align}
Pledgeable income is maximised by $R_l^F=SI$ ($R_b^F=0$): rewarding failure only tightens incentives. Breakeven binding:
\begin{align}
  \left(\rho_0+\frac Se\right)I = I-A \;\Longrightarrow\; I(e) = \frac{A}{1-[(S/e)+\rho_0]}, \label{r10e15-I}
\end{align}
positive by $\rho_0<1-S/e$; $I(e)$ falls as $e$ rises (depreciation).

\answer{(ii)} $R_l^F=SI$ and $I=A/\{1-[(S/e)+\rho_0]\}$.

\step{3}{Commitment} With $g^*$ chosen first, investment is $I(e)$ and the breakeven gives $d^*+d/e=I-A$. Entrepreneurs consume (on average) the nontradables $SI$ net of transfers and the tradables left after paying foreigners; from market clearing,
\begin{align}
  c_1^* = p_HRI(e)+\Rc^*-g^*-[I(e)-A] = \Rc^*-g^*+(p_HR-1)I(e)+A. \label{r10e15-mc}
\end{align}
The government maximises $W=SI(e)+u(c_1^*)+v(g^*)$ subject to \eqref{r10e15-I} and \eqref{r10e15-mc}, with $e=e(g^*)$, $e'>0$. Differentiating, with $u'=e$:
\begin{align}
  S\,I'(e)e'(g^*)+e\bigl[-1+(p_HR-1)I'(e)e'(g^*)\bigr]+v'(g^*) &= 0 \nonumber\\
  v'(g^*) &= e\left[1-\left(\frac Se+p_HR-1\right)I'(e)\frac{de}{dg^*}\right]>e, \label{r10e15-c}
\end{align}
since $S/e+p_HR-1>0$ (positive NPV per unit in tradables), $I'(e)<0$ and $de/dg^*>0$: \emph{underspending}. The government holds back because spending depreciates the currency and shrinks borrowing capacity.

\step{4}{Noncommitment} Now $\bar I$, $d^*$, $d$ are fixed when $g^*$ is chosen; the ex post clearing condition determines $e$ and
\begin{align}
  W^{\text{nc}} &= S\bar I+u\Bigl(\Rc^*-g^*+p_HR\bar I-d^*-\frac{d}{e(g^*)}\Bigr)+v(g^*), \nonumber\\
  v'(g^*) &= e\left[1-\frac{d}{e^2}\frac{de}{dg^*}\right]<e, \label{r10e15-nc}
\end{align}
\emph{overspending}: depreciation lowers $d/e$, the tradable value of debt owed to foreigners, a gain the government counts while ignoring foreigners' loss. Investors anticipate it, so in equilibrium $e$ is higher and $I(e)$ lower.

\step{5}{Example} $u(c)=6\log c$ (so $c_1^*=6/e$), $v(g)=\log g$, $p_H=0.8$, $R=2$, $\rho_0=0.4$, $S=0.3$, $A=1$, $\Rc^*=2$; contracts with $d^*=\rho_0I$, $d=SI$. Commitment: $g^*=0.508893$, $e=1.504114$, $I=2.496585$, $v'=1.965050>e$. Noncommitment (rational-expectations fixed point): $g^*=0.701486$, $e=1.598417$, $I=2.425335$, $v'=1.425545<e$. Welfare $8.374789$ vs.\ $8.309522$. Both first-order conditions \eqref{r10e15-c}--\eqref{r10e15-nc} hold numerically.

\answer{(iii)} Commitment: $v'(g^*)>e$, lower spending, stronger currency, higher investment. Noncommitment: $v'(g^*)<e$, more spending, depreciation, lower investment and welfare.

\step{6}{Externality} Under noncommitment, \eqref{r10e15-nc} shows that $g^*$ rises with aggregate nontradable debt $d$ (the gain from depreciating is $d\,e'/e^2$). Each borrower, by pledging $SI$ in nontradables, raises $d$, hence $g^*$ and $e$; the depreciation lowers the value of every foreigner's nontradable claims, which foreigners price in, reducing all borrowers' capacity $I(e)$. Individually, though, pledging nontradables maximises one's own borrowing capacity (Step~2).

\answer{(iv)} Without commitment, each borrower's nontradable debt invites overspending and depreciation that hurts all borrowers: every one would gain if the others issued fewer nontradable claims, yet each has an individual incentive to issue them.

\intuition A government that cannot commit treats debt owed to foreigners in domestic goods as a tax base it can inflate away. Lenders understand this and charge for it, so the country as a whole borrows less. Commitment makes the government internalise that its spending feeds back into the exchange rate and the country's borrowing capacity, so it spends less than the static optimum $v'=e$.

\pitfalls
\begin{itemize}
  \item Valuing nontradable claims at face value in tradables: they are worth $1/e$.
  \item Keeping $I$ variable under noncommitment: it is fixed when $g^*$ is chosen.
  \item Reading ``underspending'' as $v'<e$: underspending means the marginal value of the public good exceeds its cost $e$.
\end{itemize}

\bookans Matches the book's answer (p.~623): $d^*=p_HR_l^*$, $d=p_HR_l^S+(1-p_H)R_l^F$; $R_b^F=0$, $R_l^F=SI$, $I(e)=A/(1-[(S/e)+\rho_0])$; $v'(g^*)=e[1-(S/e+p_HR-1)\frac{dI}{de}\frac{de}{dg^*}]>e$ under commitment and $v'(g^*)=e[1-\frac{d}{e^2}\frac{de}{dg^*}]<e$ under noncommitment; the externality through nontradable debt. The example is an addition.
\refbib{Tirole \S16.3 (property rights, time inconsistency), pp.~544--551; \S15.2.2 (Caballero--Krishnamurthy discussion), p.~522.}

\end{document}
